\documentclass[12pt, openany]{book}

%% Fonts
\usepackage{fontspec}
\setmainfont{Latin Modern Roman}
\setmonofont{DejaVu Sans Mono}[Scale=0.82]

%% Page Layout
\usepackage[a4paper, top=2.8cm, bottom=2.8cm, left=3cm, right=3cm]{geometry}
\usepackage{microtype}

%% Math
\usepackage{amsmath, amssymb, amsthm, mathtools}

%% Colors
\usepackage[dvipsnames, svgnames, x11names]{xcolor}
\definecolor{navyblue}{HTML}{1a2744}
\definecolor{tealaccent}{HTML}{2a7f8f}
\definecolor{defgreen}{HTML}{1e5c2e}
\definecolor{defbg}{HTML}{eaf5ec}
\definecolor{thmbg}{HTML}{e8f0fb}
\definecolor{thmblue}{HTML}{1a3a6e}
\definecolor{codebg}{HTML}{f5f5f0}
\definecolor{codefg}{HTML}{1a1a2e}
\definecolor{codekw}{HTML}{0a4b8c}
\definecolor{codecomment}{HTML}{5a7a5a}
\definecolor{codestr}{HTML}{7a2020}
\definecolor{exbg}{HTML}{fdf6e3}
\definecolor{exborder}{HTML}{c9a227}
\definecolor{remarkbg}{HTML}{f5f0ff}
\definecolor{remarkborder}{HTML}{7c4daa}

%% Section formatting - use sectsty, NOT titlesec
\usepackage{sectsty}
\partfont{\color{navyblue}\bfseries\Large}
\chapterfont{\color{navyblue}\bfseries}
\sectionfont{\color{tealaccent}\bfseries}
\subsectionfont{\color{tealaccent}}

%% Boxes
\usepackage[many, breakable]{tcolorbox}
\usepackage{etoolbox}

\newtcolorbox{defbox}[2][]{breakable,enhanced,fontupper=\color{black},
  colbacktitle=defgreen,colframe=defgreen,colback=defbg,coltitle=white,
  fonttitle=\bfseries,title={Definition: #2},arc=3pt,boxrule=0.6pt,
  before upper={\color{black}},#1}
\newtcolorbox{thmbox}[2][]{breakable,enhanced,fontupper=\color{black},
  colbacktitle=thmblue,colframe=thmblue,colback=thmbg,coltitle=white,
  fonttitle=\bfseries,title={Theorem: #2},arc=3pt,boxrule=0.6pt,
  before upper={\color{black}},#1}
\newtcolorbox{lembox}[2][]{breakable,enhanced,fontupper=\color{black},
  colbacktitle=thmblue!80!black,colframe=thmblue!80!black,colback=thmbg,coltitle=white,
  fonttitle=\bfseries,title={Lemma: #2},arc=3pt,boxrule=0.6pt,
  before upper={\color{black}},#1}
\newtcolorbox{corbox}[2][]{breakable,enhanced,fontupper=\color{black},
  colbacktitle=thmblue!70!black,colframe=thmblue!70!black,colback=thmbg,coltitle=white,
  fonttitle=\bfseries,title={Corollary: #2},arc=3pt,boxrule=0.6pt,
  before upper={\color{black}},#1}
\newtcolorbox{propbox}[2][]{breakable,enhanced,fontupper=\color{black},
  colbacktitle=thmblue!60!black,colframe=thmblue!60!black,colback=thmbg,coltitle=white,
  fonttitle=\bfseries,title={Proposition: #2},arc=3pt,boxrule=0.6pt,
  before upper={\color{black}},#1}
\newtcolorbox{exbox}[1][]{breakable,enhanced,fontupper=\color{black},
  colback=exbg,colframe=exborder,colbacktitle=exborder!25!white,coltitle=black,
  fonttitle=\bfseries,title={Exercises},arc=3pt,boxrule=0.6pt,
  before upper={\color{black}},#1}
\newtcolorbox{remarkbox}[1][]{breakable,enhanced,fontupper=\color{black},
  colback=remarkbg,colframe=remarkborder,coltitle=remarkborder,
  colbacktitle=remarkbg,fonttitle=\bfseries\color{remarkborder},title={Remark},
  arc=3pt,boxrule=0.6pt,before upper={\color{black}},#1}
\newtcolorbox{leanbox}[1][]{breakable,enhanced,fontupper=\color{codefg},
  colback=codebg,colframe=thmblue!40!black,coltitle=white,
  colbacktitle=thmblue!70!black,fonttitle=\bfseries\color{white},title={Lean~4},
  arc=3pt,boxrule=0.6pt,before upper={\color{codefg}},#1}

\AfterEndEnvironment{defbox}{\color{black}}
\AfterEndEnvironment{thmbox}{\color{black}}
\AfterEndEnvironment{lembox}{\color{black}}
\AfterEndEnvironment{corbox}{\color{black}}
\AfterEndEnvironment{propbox}{\color{black}}
\AfterEndEnvironment{exbox}{\color{black}}
\AfterEndEnvironment{remarkbox}{\color{black}}
\AfterEndEnvironment{leanbox}{\color{black}}

%% Code listings
\usepackage{listings}
\lstdefinelanguage{lean4}{
  morekeywords={def,theorem,lemma,example,instance,class,structure,
    inductive,where,do,let,fun,match,with,if,then,else,
    return,namespace,open,variable,axiom,noncomputable,abbrev,
    by,exact,apply,intro,intros,rfl,simp,ring,linarith,omega,
    decide,constructor,cases,induction,rcases,obtain,use,have,
    show,calc,suffices,contradiction,trivial,assumption,ext,
    funext,congr,refl,symm,trans,and,or,not,deriving,
    Type,Prop,Sort,protected,private,mutual,repeat,first,
    norm_num,push_neg,positivity,field_simp,aesop,tauto},
  sensitive=true,
  morecomment=[l]{--},
  morecomment=[s]{/-}{-/},
  morestring=[b]",
  literate=
    {->}{{$\to$}}1
    {<-}{{$\leftarrow$}}1
    {=>}{{$\Rightarrow$}}1
    {<=}{{$\leq$}}1
    {>=}{{$\geq$}}1
    {<>}{{$\neq$}}1
    {:=}{{$:=$}}2
    {forall}{{$\forall$\,}}1
    {exists}{{$\exists$\,}}1
    {/\\}{{$\land$}}1
    {\\/}{{$\lor$}}1
    {⟨}{{$\langle$}}1
    {⟩}{{$\rangle$}}1
    {∀}{{$\forall$\,}}1
    {∃}{{$\exists$\,}}1
    {→}{{$\to$}}1
    {←}{{$\leftarrow$}}1
    {↔}{{$\leftrightarrow$}}1
    {≤}{{$\leq$}}1
    {≥}{{$\geq$}}1
    {≠}{{$\neq$}}1
    {∧}{{$\land$}}1
    {∨}{{$\lor$}}1
    {¬}{{$\lnot$}}1
    {∈}{{$\in$}}1
    {∉}{{$\notin$}}1
    {·}{{$\cdot$}}1
    {×}{{$\times$}}1
    {∣}{{$\mid$}}1
}
\lstset{
  language=lean4,
  basicstyle=\ttfamily\small\color{codefg},
  keywordstyle=\color{codekw}\bfseries,
  commentstyle=\color{codecomment}\itshape,
  stringstyle=\color{codestr},
  backgroundcolor=\color{codebg},
  frame=none,breaklines=true,keepspaces=true,
  showstringspaces=false,tabsize=2,
  xleftmargin=4pt,xrightmargin=4pt,
  aboveskip=0pt,belowskip=0pt,
}

%% Headers/Footers
\usepackage{fancyhdr}
\pagestyle{fancy}\fancyhf{}
\fancyhead[LE]{\color{navyblue}\small\itshape Verified Calculus}
\fancyhead[RO]{\color{navyblue}\small\itshape\leftmark}
\fancyfoot[C]{\color{navyblue}\thepage}
\renewcommand{\headrulewidth}{0.4pt}

%% Misc
\usepackage{booktabs,enumitem,hyperref}
\usepackage{tikz}
\usetikzlibrary{arrows.meta,positioning,shapes.geometric,fit,backgrounds,decorations.pathreplacing,calc}
\pgfdeclarelayer{background}
\pgfsetlayers{background,main}
\tikzset{every node/.append style={text=black}}
\AtBeginEnvironment{tikzpicture}{\color{black}}
\hypersetup{colorlinks=true,linkcolor=thmblue,urlcolor=tealaccent}

%% Theorem environments
\theoremstyle{plain}
\newtheorem{theorem}{Theorem}[chapter]
\newtheorem{lemma}[theorem]{Lemma}
\newtheorem{corollary}[theorem]{Corollary}
\newtheorem{proposition}[theorem]{Proposition}
\theoremstyle{definition}
\newtheorem{definition}[theorem]{Definition}
\newtheorem{example}[theorem]{Example}
\theoremstyle{remark}
\newtheorem{remark}[theorem]{Remark}

%% Macros
\newcommand{\N}{\mathbb{N}}
\newcommand{\Z}{\mathbb{Z}}
\newcommand{\Q}{\mathbb{Q}}
\newcommand{\R}{\mathbb{R}}
\newcommand{\abs}[1]{\left|#1\right|}
\newcommand{\norm}[1]{\left\|#1\right\|}
\newcommand{\eps}{\varepsilon}
\newcommand{\del}{\delta}
\newcommand{\li}[1]{{\ttfamily\small\color{tealaccent}#1}}

\begin{document}
\color{black}

%% Title Page
\begin{titlepage}
\pagecolor{navyblue}\color{white}\centering\vspace*{3.5cm}
{\fontsize{38}{48}\selectfont\bfseries Verified Calculus\par}\vspace{0.7cm}
{\fontsize{17}{24}\selectfont\color{tealaccent!80!white}
From the Natural Numbers to Symbolic Solvers in Lean~4\par}\vspace{1.8cm}
{\color{white!50!navyblue}\rule{0.55\textwidth}{0.5pt}}\vspace{1.8cm}
{\large\itshape A zero-to-hero journey through real analysis\par
with every definition proved from scratch in Lean~4\par}
\vfill
{\small\color{white!60!navyblue}No Mathlib. No shortcuts. Every theorem earned.\par}
\vspace{0.5cm}
{\small\color{white!50!navyblue}Lean~4 $\cdot$ Real Analysis $\cdot$ Symbolic Computation}
\end{titlepage}
\nopagecolor\color{black}

\frontmatter

\chapter*{Preface}
\addcontentsline{toc}{chapter}{Preface}
\color{black}

This book rebuilds calculus from a single inductive type in Lean~4 to a working command-line symbolic solver. No external mathematics libraries are imported---every definition, from addition on the natural numbers to the Fundamental Theorem of Calculus, is constructed explicitly and its properties proved.

\medskip\noindent\textbf{Prerequisites.} Familiarity with proof-based mathematics (induction, contradiction, epsilon-delta arguments seen before). No prior Lean~4 experience required.

\medskip\noindent\textbf{Structure.}
\begin{itemize}[noitemsep]
  \item \textbf{Part~I} constructs $\N$, $\Z$, $\Q$, and $\R$ via Cauchy completion.
  \item \textbf{Part~II} develops the topology of $\R$: sequences, limits, continuity.
  \item \textbf{Part~III} covers differential calculus through Taylor's theorem.
  \item \textbf{Part~IV} covers integral calculus and the Fundamental Theorem.
  \item \textbf{Part~V} treats infinite series and power series.
  \item \textbf{Part~VI} extends to $\R^n$: multivariable calculus and vector analysis.
  \item \textbf{Part~VII} is the capstone: \textit{Symboliculus}, a step-by-step symbolic calculus solver as a runnable Lean~4 CLI.
\end{itemize}

\tableofcontents

\mainmatter\color{black}

%%=============================================================================
\part{Building the Real Numbers from Scratch}
%%=============================================================================
\color{black}

%%-----------------------------------------------------------------------------
\chapter{The Natural Numbers}
%%-----------------------------------------------------------------------------
\color{black}

Every edifice needs a foundation. For mathematics, that foundation is the natural numbers. Rather than treating them as self-evident, we ask what they \emph{are}, and answer with five axioms due to Giuseppe Peano.

\begin{remarkbox}
\textbf{Historical Note: What Are the Natural Numbers?}

For most of human history, natural numbers were treated as primitives --- self-evident objects that needed no definition. The ancient Greeks distinguished between \emph{arithmetic} (the study of number) and \emph{logistic} (calculation), but never asked what a number fundamentally \emph{is}. This question only became pressing in the nineteenth century, when mathematicians began demanding rigorous foundations for analysis.

The crisis was triggered by analysis itself. Cauchy, Bolzano, and Weierstrass had built calculus on a foundation of limits and convergence, but their definitions ultimately rested on an intuitive notion of ``real number'' that nobody had made precise. To define $\R$ rigorously, you first need $\Q$; to define $\Q$ you need $\Z$; to define $\Z$ you need $\N$. And $\N$, it turned out, was the hardest of all.

Several mathematicians tackled the problem simultaneously in the 1870s--1880s. \textbf{Richard Dedekind} (1872) defined the integers and rationals via equivalence classes --- the same construction we use in Chapters~2 and~3 --- and proposed defining $\N$ itself via what he called ``simply infinite systems,'' an abstract structure satisfying conditions equivalent to Peano's axioms. \textbf{Gottlob Frege} (1884) attempted to define natural numbers as \emph{logical objects}: the number $2$ is the class of all two-element sets. Frege's program was far more ambitious --- he wanted to show arithmetic was pure logic, with no mathematical primitives at all (a position called \textbf{logicism}). \textbf{Giuseppe Peano} (1889) took the more pragmatic approach we follow here: simply axiomatise $\N$ directly, accepting $0$ and $S$ as primitives, and derive everything else.

The logicist program collapsed when Bertrand Russell sent Frege a letter in 1902 pointing out that the unrestricted comprehension principle underlying Frege's set theory was contradictory --- this is \textbf{Russell's paradox}: the set of all sets that do not contain themselves can neither contain itself nor not contain itself. Frege's life's work was destroyed. Russell and Whitehead subsequently rebuilt logicism on a more careful type-theoretic foundation in \emph{Principia Mathematica} (1910--1913), requiring hundreds of pages to prove that $1 + 1 = 2$.

Peano's axiomatic approach, by contrast, proved robust. Combined with Zermelo--Fraenkel set theory (1908--1922), it became the standard foundation for mathematics. The five axioms we state below are the distillation of this century of foundational work.
\end{remarkbox}
\color{black}

\section{The Peano Axioms}

\begin{defbox}{The Peano Axioms}
The \textbf{natural numbers} are a set $\N$ with a distinguished element $0 \in \N$ and a function $S : \N \to \N$ (\textbf{successor}) satisfying:
\begin{enumerate}[label=\textbf{P\arabic*}.,noitemsep]
  \item $0 \in \N$.
  \item If $n \in \N$ then $S(n) \in \N$.
  \item $S$ is injective: $S(m) = S(n) \Rightarrow m = n$.
  \item $0$ is not a successor: $S(n) \neq 0$ for all $n$.
  \item \textbf{Induction:} If $P(0)$ holds and $P(n) \Rightarrow P(S(n))$, then $P(n)$ holds for all $n$.
\end{enumerate}
\end{defbox}
\color{black}

We write $1 := S(0)$, $2 := S(1)$, $3 := S(2)$, etc. Axiom~P5 ensures no ``extra'' elements exist beyond what is reachable from $0$ by applying $S$.

\color{black}\begin{center}
\begin{tikzpicture}[
  node distance=1.3cm,
  every node/.style={circle, draw=thmblue, fill=white, thick,
                     minimum size=0.75cm, font=\small\bfseries},
  arr/.style={-{Stealth[length=6pt]}, thick, color=tealaccent}
]
  \node (0) {\textcolor{black}{$0$}};
  \node (1) [right=of 0] {\textcolor{black}{$1$}};
  \node (2) [right=of 1] {\textcolor{black}{$2$}};
  \node (3) [right=of 2] {\textcolor{black}{$3$}};
  \node (4) [right=of 3] {\textcolor{black}{$4$}};
  \node (dots) [right=of 4, draw=none, fill=none, font=\Large] {\textcolor{black}{$\cdots$}};

  \draw[arr] (0) -- node[above, draw=none, fill=none, font=\scriptsize\itshape] {\textcolor{tealaccent}{$S$}} (1);
  \draw[arr] (1) -- node[above, draw=none, fill=none, font=\scriptsize\itshape] {\textcolor{tealaccent}{$S$}} (2);
  \draw[arr] (2) -- node[above, draw=none, fill=none, font=\scriptsize\itshape] {\textcolor{tealaccent}{$S$}} (3);
  \draw[arr] (3) -- node[above, draw=none, fill=none, font=\scriptsize\itshape] {\textcolor{tealaccent}{$S$}} (4);
  \draw[arr] (4) -- (dots);

  % P4 annotation: 0 is not a successor
  \node[draw=none, fill=none, below=0.5cm of 0, font=\scriptsize\itshape]
    {\textcolor{red!70}{no predecessor (P4)}};

  % P3 annotation
  \node[draw=none, fill=none, below=0.5cm of 3, font=\scriptsize\itshape]
    {\textcolor{defgreen}{injective (P3)}};
\end{tikzpicture}
\end{center}

It is worth pausing on what makes these five axioms so powerful.  P1 and P2 say the natural numbers are closed under $S$.  P3 says $S$ never collapses two different numbers into one --- every natural number has a unique predecessor.  P4 says $0$ is genuinely the starting point, not in the middle of some cycle.  Without P3 and P4 together, we could have a ``number system'' like $\{0, 1, 2\}$ where $S(2) = 0$, which satisfies P1, P2, and P5 but is not the natural numbers.  P5 --- the induction axiom --- is the one that rules out such finite imposters by insisting every element is reachable from $0$.

In Lean~4, an inductive type with two constructors captures exactly these axioms: \li{zero} corresponds to $0$, and \li{succ} corresponds to $S$.  Crucially, the Lean kernel enforces P3 and P4 automatically: distinct constructors are always disjoint (\li{zero} $\neq$ \li{succ n}), and constructors are always injective (\li{succ m = succ n} $\Rightarrow$ \li{m = n}).  P5 is the elimination principle generated by the \li{inductive} declaration --- pattern matching on a \li{MyNat} is literally induction.

\begin{leanbox}
\begin{lstlisting}
-- Our own natural numbers, no imports needed
inductive MyNat : Type where
  | zero : MyNat
  | succ : MyNat -> MyNat
  deriving Repr

-- Naming the first few
def MyNat.one   : MyNat := MyNat.succ MyNat.zero
def MyNat.two   : MyNat := MyNat.succ MyNat.one
def MyNat.three : MyNat := MyNat.succ MyNat.two

-- The induction principle is generated automatically.
-- Pattern matching IS structural induction in Lean.
theorem mynat_ind (P : MyNat -> Prop)
    (h0 : P MyNat.zero)
    (hs : forall n, P n -> P (MyNat.succ n)) :
    forall n, P n := fun n => by
  induction n with
  | zero      => exact h0
  | succ n ih => exact hs n ih
\end{lstlisting}
\end{leanbox}
\color{black}

\begin{remarkbox}
\textbf{Kronecker, Brouwer, and the Constructivist Challenge}

Leopold Kronecker famously declared in 1886: ``God made the integers; all else is the work of man'' (\emph{Die ganzen Zahlen hat der liebe Gott gemacht, alles andere ist Menschenwerk}). He meant this as a compliment to $\Z$ but a criticism of the irrational numbers, infinite sets, and non-constructive proofs that were proliferating in analysis. For Kronecker, a mathematical object had to be \emph{explicitly constructible} from the integers in finitely many steps; anything else was metaphysics.

This position was radicalised by \textbf{L.E.J.\ Brouwer} (1907--) into \textbf{intuitionism}: mathematics is a mental activity, and a mathematical object exists only when it can be \emph{mentally constructed}. On this view, the law of excluded middle ($P \lor \lnot P$) is not automatically valid, because there may be propositions $P$ for which we can neither construct a proof of $P$ nor a proof of $\lnot P$. A proof of $\lnot(\lnot P)$ (not-not-$P$) is \emph{not} the same as a proof of $P$, because it merely rules out a disproof without providing a construction.

Lean~4 is, at its core, a \textbf{dependent type theory} based on the Calculus of Inductive Constructions, which is broadly constructive: proofs in Lean are \emph{programs}, and every proof of an existential $\exists x, P(x)$ must exhibit a specific witness. However, Lean also provides \texttt{Classical.choice}, which non-constructively selects a witness from any nonempty type, enabling the full law of excluded middle. The theorems in this book use classical logic freely --- we prove existence by contradiction where convenient, we use the least upper bound property which requires choice, and so on.

A fully constructive treatment of real analysis --- one that Brouwer would endorse --- is substantially harder. Bishop's \emph{Constructive Analysis} (1967) showed it can be done, but the proofs are longer and some classical results must be weakened. The Lean/Mathlib library contains both classical and constructive portions, and flags which results require \texttt{Classical} axioms.

The key point for us: when we write \texttt{inductive MyNat}, we are defining $\N$ \emph{constructively} --- each natural number is literally a finite sequence of \texttt{succ} constructors. But when we later invoke the least upper bound property of $\R$ to prove the Intermediate Value Theorem, we are using classical logic. Both are legitimate mathematics; they just live in different epistemic universes.
\end{remarkbox}
\color{black}

\section{Arithmetic by Recursion}

The Peano axioms give us $\N$ as a bare set with a successor function.  We have not yet said what $+$ or $\times$ mean --- those must be \emph{defined}, not assumed.  The tool for doing so is the recursion principle: to define a function $f : \N \to X$, it suffices to specify $f(0)$ and a rule computing $f(S(n))$ from $f(n)$.  This is just the constructive reading of induction: define the function by building it bottom-up from the base case.

The definitions below may look strange at first --- why define $n + S(m)$ rather than $S(n) + m$?  The answer is that we recurse on the \emph{second} argument because Lean's pattern matching must terminate, and it terminates when we peel off one \li{succ} constructor at a time from a fixed argument.  The commutativity $m + n = n + m$ is then a non-trivial \emph{theorem} that requires induction to prove, not something built into the definition.

\begin{defbox}{Addition and Multiplication on $\N$}
\textbf{Addition} is defined by recursion on the second argument:
\[ n + 0 := n, \qquad n + S(m) := S(n + m). \]
\textbf{Multiplication}:
\[ n \cdot 0 := 0, \qquad n \cdot S(m) := (n \cdot m) + n. \]
\end{defbox}
\color{black}

\begin{leanbox}
\begin{lstlisting}
@[simp] def MyNat.add : MyNat -> MyNat -> MyNat
  | n, .zero   => n
  | n, .succ m => .succ (n.add m)

-- Two cases only. Overlapping patterns (e.g. a zero+n shortcut) break
-- simp's equation lemmas, so we keep the definition non-overlapping
-- and prove zero_add as a theorem instead.

@[simp] def MyNat.mul : MyNat -> MyNat -> MyNat
  | _, .zero   => .zero
  | n, .succ m => (n.mul m).add n

instance : Add MyNat := ⟨MyNat.add⟩
instance : Mul MyNat := ⟨MyNat.mul⟩

-- Bridge lemmas: connect + and * notation through the typeclass
-- instance to the raw definitions so simp can unfold them.
-- Without these, simp [MyNat.add] sees + but can't unfold it.
@[simp] theorem MyNat.add_def (m n : MyNat) : m + n = m.add n := rfl
@[simp] theorem MyNat.mul_def (m n : MyNat) : m * n = m.mul n := rfl

-- zero + n = n. Not immediate from the definition since add
-- recurses on the SECOND argument; needs induction on n.
theorem MyNat.zero_add (n : MyNat) : MyNat.zero.add n = n := by
  induction n with
  | zero      => rfl
  | succ n ih => simp [ih]

-- n.succ.add m = (n.add m).succ. Succ is on the LEFT argument
-- which the definition doesn't recurse on; needs induction on m.
theorem MyNat.succ_add (n m : MyNat) : n.succ.add m = (n.add m).succ := by
  induction m with
  | zero      => rfl
  | succ m ih => simp [ih]

-- Commutativity: uses both zero_add and succ_add as helpers.
theorem MyNat.add_comm (m n : MyNat) : m.add n = n.add m := by
  induction n with
  | zero      => simp [MyNat.zero_add]
  | succ n ih => simp [MyNat.succ_add, ih]

-- Associativity: follows directly from the definition by induction on c.
theorem MyNat.add_assoc (a b c : MyNat) : (a.add b).add c = a.add (b.add c) := by
  induction c with
  | zero      => rfl
  | succ c ih => simp [ih]
\end{lstlisting}
\end{leanbox}
\color{black}

\section{The Order on \texorpdfstring{$\N$}{N}}

Arithmetic alone does not give us a way to compare natural numbers --- we need an ordering.  The right definition turns out to be purely additive: $m \leq n$ means there is some $k$ that witnesses the gap $n - m$ (without subtraction!).  This is the only way to define order from addition without assuming subtraction exists.

The \emph{well-ordering} property is what distinguishes $\N$ from other ordered sets like $\Z$ or $\Q$: every nonempty subset has a smallest element.  This is the principle underlying nearly every induction argument we will make in the rest of the book --- if something fails, it fails at some \emph{smallest} $n$, and analysing that $n$ gives the contradiction.  Strong induction and well-ordering are two faces of the same coin, as the proof below shows.

The three forms of induction are logically equivalent; choose whichever gives the most natural proof:

\color{black}\begin{center}
\begin{tikzpicture}[
  box/.style={draw=thmblue, fill=white, rounded corners=4pt,
              minimum width=3.8cm, minimum height=1.5cm,
              align=center, font=\small, thick},
  arr/.style={<->, thick, color=tealaccent, shorten >=3pt, shorten <=3pt}
]
  \node[box] (weak) at (0,0)
    {\textcolor{black}{\textbf{Weak Induction}}\\[3pt]
     \textcolor{black}{Assume $P(n)$,}\\
     \textcolor{black}{prove $P(n{+}1)$}};

  \node[box] (strong) at (5,0)
    {\textcolor{black}{\textbf{Strong Induction}}\\[3pt]
     \textcolor{black}{Assume $\forall k{<}n,\,P(k)$,}\\
     \textcolor{black}{prove $P(n)$}};

  \node[box] (wo) at (2.5,-2.6)
    {\textcolor{black}{\textbf{Well-Ordering}}\\[3pt]
     \textcolor{black}{Every nonempty $S\!\subseteq\!\N$}\\
     \textcolor{black}{has a least element}};

  \draw[arr] (weak) -- node[above, font=\scriptsize\itshape] {\textcolor{tealaccent}{equivalent}} (strong);
  \draw[arr] (weak) -- node[left,  font=\scriptsize\itshape, xshift=-2pt] {\textcolor{tealaccent}{equivalent}} (wo);
  \draw[arr] (strong) -- node[right, font=\scriptsize\itshape, xshift=2pt] {\textcolor{tealaccent}{equivalent}} (wo);

  % use-when labels
  \node[font=\scriptsize, align=center] at (0, -1.1)
    {\textcolor{defgreen}{use when step needs}\\
     \textcolor{defgreen}{only $P(n)$}};
  \node[font=\scriptsize, align=center] at (5, -1.1)
    {\textcolor{defgreen}{use when step needs}\\
     \textcolor{defgreen}{earlier values}};
  \node[font=\scriptsize, align=center] at (2.5, -4.0)
    {\textcolor{defgreen}{use for minimality /}\\
     \textcolor{defgreen}{existence arguments}};
\end{tikzpicture}
\end{center}

\begin{defbox}{Natural Order}
$m \leq n$ if there exists $k \in \N$ with $m + k = n$.
$m < n$ means $m \leq n$ and $m \neq n$.
\end{defbox}
\color{black}

\begin{thmbox}{$(\N, \leq)$ is a Well-Ordered Total Order}
$\leq$ is reflexive, antisymmetric, transitive, and total. Moreover, every nonempty subset of $\N$ has a least element (\textbf{well-ordering}).
\end{thmbox}
\color{black}

\begin{proof}
Reflexivity: $m + 0 = m$. Antisymmetry: if $m + k = n$ and $n + j = m$ then $m + k + j = m$ so $k = j = 0$, giving $m = n$. Transitivity: $m + k = n$ and $n + j = p$ gives $m + (k+j) = p$. Totality: by induction on $n$. Well-ordering follows from the equivalence theorem below.
\end{proof}

\begin{thmbox}{The Three Forms of Induction Are Equivalent}
The following three principles are logically equivalent --- each can be derived from the others using only the definitions:
\begin{enumerate}[noitemsep]
  \item \textbf{Weak induction (P5):} If $P(0)$ and $\forall n,\, P(n) \Rightarrow P(n{+}1)$, then $\forall n,\, P(n)$.
  \item \textbf{Strong induction:} If $\forall n,\, (\forall k < n,\, P(k)) \Rightarrow P(n)$, then $\forall n,\, P(n)$.
  \item \textbf{Well-ordering:} Every nonempty $S \subseteq \N$ has a least element.
\end{enumerate}

\bigskip
\textbf{Proof: (1) $\Rightarrow$ (2).}\quad
Suppose strong induction's hypothesis holds: $(\forall k < n,\, P(k)) \Rightarrow P(n)$ for all $n$.
Define $Q(n) := \forall k \leq n,\, P(k)$.  We prove $\forall n,\, Q(n)$ by \emph{weak} induction.

\textit{Base case $Q(0)$:} We need $P(0)$.  The strong hypothesis at $n = 0$ reads $(\forall k < 0,\, P(k)) \Rightarrow P(0)$.  The premise is vacuously true (no $k < 0$ exists), so $P(0)$ holds.

\textit{Step:} Assume $Q(n)$, i.e., $\forall k \leq n,\, P(k)$.  For $Q(n{+}1)$ we need $P(n{+}1)$ additionally.  The strong hypothesis at $n{+}1$ reads $(\forall k \leq n,\, P(k)) \Rightarrow P(n{+}1)$, and $Q(n)$ is exactly the premise.  So $P(n{+}1)$ holds and $Q(n{+}1)$ follows.

By weak induction $\forall n,\, Q(n)$, so $\forall n,\, P(n)$. \qquad$\square$

\bigskip
\textbf{Proof: (2) $\Rightarrow$ (3).}\quad
Let $S \subseteq \N$ be nonempty; assume for contradiction it has no least element.
Set $P(n) := n \notin S$ and apply strong induction.

\textit{Step:} Suppose $\forall k < n,\, k \notin S$.  If $n \in S$, then $n$ would be the least element of $S$ (all smaller elements are excluded).  But $S$ has no least element --- contradiction.  So $n \notin S$, i.e., $P(n)$.

By strong induction $\forall n,\, n \notin S$, so $S = \emptyset$ --- contradicting $S$ nonempty. \qquad$\square$

\bigskip
\textbf{Proof: (3) $\Rightarrow$ (1).}\quad
Assume well-ordering.  Suppose $P(0)$ holds and $\forall n,\, P(n) \Rightarrow P(n{+}1)$.
Let $S = \{n \in \N : \lnot P(n)\}$ be the set of \emph{counterexamples}.

Suppose $S \neq \emptyset$.  By well-ordering, $S$ has a least element $m$.
Since $P(0)$ holds, $0 \notin S$, so $m \neq 0$, giving $m = k{+}1$ for some $k$.
Since $k < m$ and $m$ is the \emph{least} element of $S$, we have $k \notin S$, so $P(k)$ holds.
By the step hypothesis $P(k) \Rightarrow P(k{+}1) = P(m)$, so $m \notin S$ --- contradiction.

Therefore $S = \emptyset$ and $\forall n,\, P(n)$. \qquad$\square$
\end{thmbox}
\color{black}

\begin{leanbox}
\begin{lstlisting}
-- We can encode all three in Lean and derive each from the others.

-- (1) Weak induction -- built into MyNat's eliminator.

-- (2) Strong induction -- derived from weak by the Q(n) trick.
-- Define Q n := forall k <= n, P k, prove Q by weak induction,
-- then P follows since P(n) is implied by Q(n).
theorem MyNat.strong_ind (P : MyNat -> Prop)
    (step : forall n, (forall k, k < n -> P k) -> P n) :
    forall n, P n := by
  suffices h : forall n, forall k, k <= n -> P k from
    fun n => h n n (le_refl n)
  intro n
  induction n with
  | zero =>
    intro k hk
    have hkz : k = MyNat.zero := Nat.le_zero.mp hk
    subst hkz
    exact step .zero (fun k hlt => absurd hlt (Nat.not_lt_zero k))
  | succ n ih =>
    intro k hk
    rcases Nat.lt_or_eq_of_le hk with hlt | rfl
    · exact ih k (Nat.lt_succ_iff.mp hlt)
    · exact step (.succ n) (fun j hj => ih j (Nat.lt_succ_iff.mpr hj))

-- Key lemma: n + k.succ ≠ n (can't add a positive and return to start)
theorem MyNat.no_overflow (n m : MyNat) : n.add m.succ ≠ n := by
  induction n with
  | zero      => simp [MyNat.add]
  | succ n ih => simp [MyNat.add]; exact ih

-- Corollary: if n + k = n then k = 0
theorem MyNat.add_eq_self (n k : MyNat) : n.add k = n -> k = .zero := by
  cases k with
  | zero   => intros; rfl
  | succ k => intro h; exact absurd h (MyNat.no_overflow n k)

-- If a + b = 0 then a = 0 (and by symmetry b = 0)
theorem MyNat.add_eq_zero_left {a b : MyNat}
    (h : a.add b = .zero) : a = .zero := by
  cases a with
  | zero   => rfl
  | succ a => cases b <;> simp [MyNat.add] at h

-- Antisymmetry: uses add_eq_self to avoid case-splitting on witnesses
theorem MyNat.le_antisymm {n m : MyNat}
    (h1 : n <= m) (h2 : m <= n) : n = m := by
  obtain ⟨k, hk⟩ := h1
  obtain ⟨j, hj⟩ := h2
  -- n + k = m and m + j = n, so n + (k + j) = n
  have hkj : n.add (k.add j) = n := by
    calc n.add (k.add j)
        = (n.add k).add j := (MyNat.add_assoc n k j).symm
      _ = m.add j         := by rw [show n.add k = m from hk]
      _ = n               := hj
  -- Therefore k + j = 0, so k = 0, so n = m
  have h0 : k.add j = .zero := MyNat.add_eq_self n _ hkj
  have hk0 : k = .zero     := MyNat.add_eq_zero_left h0
  simp [hk0, MyNat.add] at hk
  exact hk.symm

-- (3) Well-ordering -- every nonempty S ⊆ MyNat has a least element.
theorem MyNat.well_ordering (S : MyNat -> Prop) (hne : exists n, S n) :
    exists m, S m /\ forall k, S k -> m <= k := by
  by_contra hall
  push_neg at hall
  apply MyNat.strong_ind (fun n => ¬ S n)
  · intro n ih hSn
    obtain ⟨k, hSk, hkn⟩ := hall n hSn
    exact ih k hkn hSk
  · obtain ⟨n, hn⟩ := hne
    exact absurd hn (MyNat.strong_ind _
      (fun n ih hSn => (hall n hSn).elim
        (fun ⟨k, hk, hlt⟩ => ih k hlt hk)) n)
\end{lstlisting}
\end{leanbox}
\color{black}

\begin{exbox}
\begin{enumerate}
  \item Prove $a \cdot (b + c) = a \cdot b + a \cdot c$ for $\N$ by induction on $c$.
  \item Prove the cancellation law: $a + c = b + c \Rightarrow a = b$.
  \item Define exponentiation recursively and prove $n^{j+k} = n^j \cdot n^k$.
  \item Implement truncated subtraction $m \mathbin{\dot{-}} n$ (returns $0$ when $n > m$) and prove $m \mathbin{\dot{-}} m = 0$.
  \item Prove the division algorithm: for $b > 0$, there exist unique $q, r$ with $r < b$ and $a = bq + r$.
  \item Prove that $\N$ satisfies P3 (injectivity of succ) and P4 ($0$ is not a successor) from the inductive type definition.
\end{enumerate}
\end{exbox}
\color{black}

%%-----------------------------------------------------------------------------
\chapter{The Integers}
%%-----------------------------------------------------------------------------
\color{black}

The natural numbers have no subtraction: $3 + x = 2$ has no solution in $\N$. We introduce negative numbers via a quotient construction---an integer is a \emph{difference} of two naturals.

\section{Constructing \texorpdfstring{$\Z$}{Z}}

The key idea is that a negative number is not a new primitive thing but rather a \emph{relationship between two natural numbers}: the integer $-3$ is just the claim that something is $3$ less than something else.  More precisely, the pair $(2, 5)$ represents $2 - 5 = -3$, but so does $(0, 3)$, $(7, 10)$, and $(100, 103)$.  All of these pairs should be identified as the same integer.

This is our first encounter with the \textbf{quotient construction}: we have a ``too big'' set ($\N \times \N$, with many redundant representatives for each integer) and an equivalence relation that collapses the redundancy.  The integers are what remains after making all equivalent pairs identical.  The same idea appeared implicitly in everyday arithmetic --- when you learned that $\frac{2}{4} = \frac{1}{2}$, you were working with a quotient --- and we will use it once more in Chapter~3 to build $\Q$.

Notice the elegance of the equivalence condition $(a, b) \sim (c, d) \iff a + d = b + c$: it avoids subtraction entirely, expressing $a - b = c - d$ using only addition, which we already have on $\N$.

\color{black}\begin{center}
\begin{tikzpicture}[
  pairnode/.style={draw=thmblue, fill=white, rounded corners=3pt,
                   font=\small, inner sep=5pt, thick},
  grp/.style={draw=defgreen!70, fill=defbg, rounded corners=5pt,
              dashed, thick, inner sep=8pt},
  znode/.style={draw=thmblue, fill=white, rounded corners=4pt,
                minimum width=1.1cm, minimum height=0.6cm,
                font=\small\bfseries, thick},
  arr/.style={-{Stealth[length=6pt]}, thick, color=tealaccent}
]
  % N x N label
  \node[font=\footnotesize\bfseries, text=navyblue] at (-3.0, 2.3)
    {\textcolor{navyblue}{$\N\times\N$}};

  % --- equivalence class of -1 ---
  % Pair nodes first, then fit-box drawn BEHIND them on background layer
  \node[pairnode] (p01) at (-4.5, 1.0)  {\textcolor{black}{$(0,1)$}};
  \node[pairnode] (p12) at (-3.0, 1.0)  {\textcolor{black}{$(1,2)$}};
  \node[pairnode] (p23) at (-1.5, 1.0)  {\textcolor{black}{$(2,3)$}};
  \node[font=\small] (d1) at (-0.4, 1.0) {\textcolor{black}{$\cdots$}};
  \begin{pgfonlayer}{background}
    \node[grp, fit=(p01)(p12)(p23)(d1)] (g1) {};
  \end{pgfonlayer}
  \node[font=\scriptsize\itshape] at (-3.0, 0.35)
    {\textcolor{defgreen}{all represent $-1$}};

  % --- equivalence class of 0 ---
  \node[pairnode] (p00) at (-4.5,-0.6)  {\textcolor{black}{$(0,0)$}};
  \node[pairnode] (p11) at (-3.0,-0.6)  {\textcolor{black}{$(1,1)$}};
  \node[pairnode] (p22) at (-1.5,-0.6)  {\textcolor{black}{$(2,2)$}};
  \node[font=\small] (d0) at (-0.4,-0.6) {\textcolor{black}{$\cdots$}};
  \begin{pgfonlayer}{background}
    \node[grp, fit=(p00)(p11)(p22)(d0)] (g0) {};
  \end{pgfonlayer}
  \node[font=\scriptsize\itshape] at (-3.0,-1.25)
    {\textcolor{defgreen}{all represent $0$}};

  % --- equivalence class of +1 ---
  \node[pairnode] (p10) at (-4.5,-2.1)  {\textcolor{black}{$(1,0)$}};
  \node[pairnode] (p21) at (-3.0,-2.1)  {\textcolor{black}{$(2,1)$}};
  \node[pairnode] (p32) at (-1.5,-2.1)  {\textcolor{black}{$(3,2)$}};
  \node[font=\small] (d2) at (-0.4,-2.1) {\textcolor{black}{$\cdots$}};
  \begin{pgfonlayer}{background}
    \node[grp, fit=(p10)(p21)(p32)(d2)] (g2) {};
  \end{pgfonlayer}
  \node[font=\scriptsize\itshape] at (-3.0,-2.75)
    {\textcolor{defgreen}{all represent $+1$}};

  % --- Z on the right ---
  \node[font=\footnotesize\bfseries] at (2.8, 2.3)
    {\textcolor{navyblue}{$\Z = (\N\times\N)/{\sim}$}};
  \node[znode] (z1) at (2.8, 1.0)  {\textcolor{black}{$-1$}};
  \node[znode] (z0) at (2.8,-0.6)  {\textcolor{black}{$0$}};
  \node[znode] (z2) at (2.8,-2.1)  {\textcolor{black}{$+1$}};
  \node[font=\small] at (2.8,-3.0)  {\textcolor{black}{$\vdots$}};

  % Quotient arrows
  \draw[arr] (g1.east) --
    node[above, font=\scriptsize\itshape] {\textcolor{tealaccent}{$/{\sim}$}}
    (z1.west);
  \draw[arr] (g0.east) -- (z0.west);
  \draw[arr] (g2.east) -- (z2.west);
\end{tikzpicture}
\end{center}

\begin{defbox}{The Integer Equivalence Relation}
Define $\sim$ on $\N \times \N$ by:
\[ (a, b) \sim (c, d) \iff a + d = b + c. \]
Intuitively, $(a,b)$ represents $a - b$, and two pairs represent the same integer iff their differences agree.
The integers are $\Z := (\N \times \N)/{\sim}$.
\end{defbox}
\color{black}

\begin{thmbox}{$\sim$ is an Equivalence Relation}
\textbf{Proof.}
\textbf{Refl:} $(a,b)\sim(a,b)$ since $a+b=b+a$ by commutativity in $\N$.
\textbf{Symm:} $a+d=b+c \Rightarrow c+b=d+a$, so $(c,d)\sim(a,b)$.
\textbf{Trans:} From $a+d=b+c$ and $c+f=d+e$, add both equations to get
$a+d+c+f=b+c+d+e$; cancel $c+d$ (cancellation law for $\N$) to obtain
$a+f=b+e$, so $(a,b)\sim(e,f)$. \qed
\end{thmbox}

\begin{leanbox}
\begin{lstlisting}
-- We build Z on top of the MyNat from Chapter 1.
-- (a, b) represents the integer a - b.
-- Note: IntRel uses Nat (built-in) so omega and simp work cleanly.
def IntRel : (Nat × Nat) -> (Nat × Nat) -> Prop
  | (a, b), (c, d) => a + d = b + c

theorem IntRel.refl : forall p : Nat × Nat, IntRel p p
  | (a, b) => by simp [IntRel, Nat.add_comm]

theorem IntRel.symm : forall p q, IntRel p q -> IntRel q p := by
  intro (a, b) (c, d) h; simp [IntRel] at *; omega

-- Transitivity: uses the cancellation law for MyNat
theorem IntRel.trans : forall p q r,
    IntRel p q -> IntRel q r -> IntRel p r := by
  intro (a, b) (c, d) (e, f) h1 h2
  simp [IntRel] at *
  omega

-- MyInt as the quotient
def MyInt : Type := Quot IntRel

def mkInt (a b : Nat) : MyInt := Quot.mk IntRel (a, b)

-- The canonical embedding Nat -> MyInt
def intOfNat (n : Nat) : MyInt := mkInt n 0

-- Familiar constants
def intZero : MyInt := mkInt 0 0
def intOne  : MyInt := mkInt 1 0

-- Negation: -(a - b) = b - a
def intNeg (x : MyInt) : MyInt :=
  Quot.lift (fun p => mkInt p.2 p.1)
    (by intro (a, b) (c, d) h
        apply Quot.sound
        simp [IntRel] at *; omega) x
\end{lstlisting}
\end{leanbox}
\color{black}

\section{Ring Axioms for \texorpdfstring{$\Z$}{Z}}

Now that we have $\Z$ as a set, we must equip it with arithmetic.  The definitions look natural once you remember that $(a, b)$ means $a - b$: adding two differences adds numerators and denominators componentwise, and multiplying uses the FOIL identity $(a-b)(c-d) = (ac+bd) - (ad+bc)$.

The subtle point is \textbf{well-definedness}.  When we write $[(a,b)] + [(c,d)] := [(a+c, b+d)]$, we are defining addition on \emph{equivalence classes} by choosing \emph{representatives}.  This is only legitimate if swapping in a different representative for $[(a,b)]$ or $[(c,d)]$ gives the same equivalence class on the right-hand side.  Checking well-definedness is not optional: it is the mathematical price of the quotient construction, and it must be paid for every operation we define.  In Lean~4 this is enforced by the type system via \li{Quot.lift} and \li{Quot.lift2}, which demand a proof of well-definedness before they will build the lifted function.

A ring satisfying $ab = 0 \Rightarrow a = 0$ or $b = 0$ is called an \textbf{integral domain}.  This property says there are no ``zero divisors'' --- no two nonzero elements whose product is zero.  $\Z$ is an integral domain; this fact was used in the transitivity proof for $\approx$ in Chapter~3, and it is what allows us to cancel a common nonzero factor.

\begin{defbox}{Addition and Multiplication on $\Z$}
$[(a,b)] + [(c,d)] := [(a+c,\; b+d)]$

$[(a,b)] \cdot [(c,d)] := [(ac+bd,\; ad+bc)]$

$-[(a,b)] := [(b,a)]$

(The multiplication formula follows from $(a-b)(c-d) = ac+bd - (ad+bc)$.)
\end{defbox}
\color{black}

\begin{thmbox}{$\Z$ is a Commutative Ring}
$\Z$ satisfies all ring axioms; in particular it is an integral domain: $ab = 0 \Rightarrow a = 0$ or $b = 0$.
\end{thmbox}
\color{black}

The proof of each axiom proceeds by lifting from $\N \times \N$ representatives (where $\N$ here means our \li{MyNat} from Chapter~1), reducing each axiom to the arithmetic facts proved there. Well-definedness requires checking that the result is independent of the chosen representative.

\begin{exbox}
\begin{enumerate}
  \item Implement \li{addInt} and \li{mulInt} in Lean~4 with well-definedness proofs.
  \item Prove $a \cdot 0 = 0$ in $\Z$ directly from the definition.
  \item Prove $(-1) \cdot a = -a$ in $\Z$.
  \item Define the order on $\Z$ and prove it is a total order compatible with addition and multiplication by positives.
  \item Prove $\Z$ has no zero divisors by reducing to the analogous fact for $\N$.
\end{enumerate}
\end{exbox}
\color{black}

%%-----------------------------------------------------------------------------
\chapter{The Rationals}
%%-----------------------------------------------------------------------------
\color{black}

The integers lack division: $2x = 3$ has no solution in $\Z$. We apply the same quotient trick a second time, this time forming fractions.

\section{Constructing \texorpdfstring{$\Q$}{Q}}

We are repeating the same trick as Chapter~2, but one level up.  There we turned $\N$ into $\Z$ by formally introducing differences $a - b$.  Now we turn $\Z$ into $\Q$ by formally introducing quotients $a/b$.  The equivalence relation $\frac{a}{b} = \frac{c}{d} \iff ad = bc$ is the cross-multiplication criterion you learned in school, now elevated to a definition.

The requirement that the denominator be nonzero ($b \in \Z^*$) is essential: $\frac{a}{0}$ is not a number.  In Lean~4 we enforce this by storing a proof \li{den\_pos : 0 < den} inside the \li{MyRat} structure --- the type system makes it impossible to construct a rational number with denominator zero.  This is an example of \textbf{dependent types} at work: the type of a value depends on a proof.

One subtlety: we defined \li{MyRat} as a \li{structure} (a concrete record) rather than a quotient type.  To use it correctly we must always work with the equivalence relation \li{ratEquiv} explicitly when stating theorems, rather than appealing to definitional equality.  An alternative design using \li{Quot} would match our $\Z$ construction more closely, and is left as an exercise.

\begin{defbox}{The Rational Equivalence Relation}
Let $\Z^* := \Z \setminus \{0\}$. Define $\approx$ on $\Z \times \Z^*$ by:
\[ (a, b) \approx (c, d) \iff a \cdot d = b \cdot c. \]
The rationals are $\Q := (\Z \times \Z^*)/{\approx}$. The class $[(a,b)]$ represents $a/b$.
\end{defbox}
\color{black}

\begin{thmbox}{$\approx$ is an Equivalence Relation}
\textbf{Proof.}
Reflexivity ($ad = da$) and symmetry are immediate.
\textbf{Trans:} From $ad = bc$ and $cf = de$, multiply to get $adf = bcf = bde$;
cancel $d \neq 0$ (using that $\Z$ is an integral domain) to obtain $af = be$. \qed
\end{thmbox}
\color{black}

\begin{defbox}{Field Operations on $\Q$}
\[ \frac{a}{b} + \frac{c}{d} := \frac{ad+bc}{bd}, \qquad \frac{a}{b} \cdot \frac{c}{d} := \frac{ac}{bd}, \qquad \left(\frac{a}{b}\right)^{-1} := \frac{b}{a}\; (a \neq 0). \]
\end{defbox}
\color{black}

\begin{thmbox}{$\Q$ is an Ordered Field}
$\Q$ is a field under these operations. It admits a total order compatible with addition and multiplication: $a/b < c/d$ (for $b, d > 0$) iff $ad < bc$ in $\Z$. This order makes $\Q$ an ordered field.
\end{thmbox}
\color{black}

\begin{thmbox}{Density of $\Q$}
Between any two distinct rationals there is another: for $p < q$ in $\Q$, the midpoint $(p+q)/2$ satisfies $p < (p+q)/2 < q$.
\end{thmbox}
\color{black}

\begin{thmbox}{Archimedean Property of $\Q$}
For any $q > 0$ in $\Q$, there exists $n \in \N$ with $1/n < q$.
\end{thmbox}
\color{black}

Despite density and the Archimedean property, $\Q$ has gaps---the subject of Chapter~4.

\begin{leanbox}
\begin{lstlisting}
-- Rational numbers using integers
structure MyRat where
  num : Int
  den : Int
  den_pos : 0 < den

def ratEquiv (p q : MyRat) : Prop :=
  p.num * q.den = q.num * p.den

def addRat (p q : MyRat) : MyRat :=
  { num     := p.num * q.den + q.num * p.den,
    den     := p.den * q.den,
    den_pos := Int.mul_pos p.den_pos q.den_pos }

def mulRat (p q : MyRat) : MyRat :=
  { num     := p.num * q.num,
    den     := p.den * q.den,
    den_pos := Int.mul_pos p.den_pos q.den_pos }

def invRat (p : MyRat) (h : p.num ≠ 0) : MyRat :=
  if hp : 0 < p.num then
    { num := p.den, den := p.num, den_pos := hp }
  else
    { num := -p.den, den := -p.num,
      den_pos := by omega }

-- The midpoint of two rationals
def midpoint (p q : MyRat) : MyRat :=
  { num     := p.num * q.den + q.num * p.den,
    den     := 2 * p.den * q.den,
    den_pos := by positivity }
\end{lstlisting}
\end{leanbox}
\color{black}

\begin{exbox}
\begin{enumerate}
  \item Prove that $\Q$ has characteristic $0$: $n \cdot 1_\Q \neq 0$ for all $n > 0$.
  \item Prove the field axioms (additive inverse, multiplicative inverse for nonzero) for $\Q$.
  \item Prove well-definedness of \li{addRat} with respect to \li{ratEquiv}.
  \item Show $\Q$ is countable: exhibit an injection $\Q \hookrightarrow \N \times \N$.
\end{enumerate}
\end{exbox}
\color{black}

%%-----------------------------------------------------------------------------
\chapter{Why \texorpdfstring{$\Q$}{Q} Is Not Enough}
%%-----------------------------------------------------------------------------
\color{black}

\begin{remarkbox}
\textbf{Historical Note: The Crisis of the Irrational}

The discovery that $\sqrt{2}$ is irrational reportedly shook the Pythagorean brotherhood to its foundations. The Pythagoreans believed that all of reality was expressible in ratios of whole numbers --- ``all is number,'' in their slogan. The proof that the diagonal of a unit square cannot be so expressed was not merely an abstract theorem; it was an existential threat to their worldview.

The legend that Hippasus of Metapontum was drowned for revealing this result to outsiders is probably apocryphal, but it captures the genuine philosophical shock. The Greeks called such lengths \emph{alogos} (without ratio) or \emph{arrhetos} (inexpressible) --- words that eventually became our ``irrational.'' Euclid includes the proof in Book X of the \emph{Elements} (c.\ 300 BCE), buried deep enough that most readers never reached it.

The proof we give below is essentially Euclid's, though he phrased it geometrically rather than algebraically. The key ingredient is what we would call the \textbf{fundamental theorem of arithmetic} --- the unique factorisation of integers into primes --- which Euclid also proved in Book IX. The modern algebraic form (if $p$ is prime and $p \mid n^2$ then $p \mid n$) was made explicit by Gauss in the \emph{Disquisitiones Arithmeticae} (1801).

For two thousand years after the Greeks, mathematicians worked with irrational numbers as geometric magnitudes (lengths, areas, volumes) rather than as numbers in their own right. It was only the arithmetisation of analysis in the nineteenth century that forced mathematicians to ask: what \emph{is} $\sqrt{2}$, precisely? The answer --- Dedekind cuts, or Cauchy sequences, or any other rigorous construction of $\R$ --- did not arrive until 1872.
\end{remarkbox}
\color{black}

\section{The Irrationality of \texorpdfstring{$\sqrt{2}$}{sqrt 2}}

The simplest demonstration that $\Q$ is incomplete is the equation $x^2 = 2$.  This has an obvious geometric solution --- the diagonal of a unit square has length exactly $\sqrt{2}$ --- yet no fraction $p/q$ satisfies it.  The ancient Greeks found this result shocking: the Pythagoreans reportedly tried to keep it secret, and legend has it that Hippasus of Metapontum was drowned for revealing it.

The proof is a classic argument by contradiction using the unique factorisation of integers.  It generalises immediately: the same argument shows $\sqrt{n} \notin \Q$ whenever $n$ is not a perfect square.  In Lean~4 we can give the full formal proof using only the divisibility machinery of \li{Nat}, which Lean's standard library provides --- we do not need our full \li{MyInt} or \li{MyRat} constructions for this particular theorem.

\begin{thmbox}{$\sqrt{2} \notin \Q$}
No rational number squares to $2$.

\textbf{Proof.} Suppose $p/q = \sqrt{2}$ with $\gcd(p,q)=1$. Then $p^2 = 2q^2$, so $p$ is even; write $p = 2k$. Then $4k^2 = 2q^2$, so $q^2 = 2k^2$, making $q$ even. But then $\gcd(p,q) \geq 2$ --- contradiction. \qed
\end{thmbox}

\begin{leanbox}
\begin{lstlisting}
theorem even_of_sq_even (n : Nat) (h : 2 ∣ n * n) : 2 ∣ n := by
  by_contra hodd
  rw [Nat.not_dvd_iff_odd (by norm_num)] at hodd
  obtain ⟨k, rfl⟩ := hodd
  simp [Nat.mul_add, Nat.add_mul] at h
  omega

theorem sqrt2_irrational :
    forall p q : Nat, q ≠ 0 -> Nat.Coprime p q -> p * p ≠ 2 * (q * q) := by
  intro p q hq hcop heq
  have h2p : 2 ∣ p := even_of_sq_even p ⟨q * q, by linarith⟩
  obtain ⟨k, rfl⟩ := h2p
  have hq2 : q * q = 2 * (k * k) := by linarith
  have h2q : 2 ∣ q := even_of_sq_even q ⟨k * k, by linarith⟩
  have : 2 ∣ Nat.gcd (2 * k) q := Nat.dvd_gcd ⟨k, rfl⟩ h2q
  rw [hcop] at this
  exact absurd this (by norm_num)
\end{lstlisting}
\end{leanbox}
\color{black}

\section{Cauchy Sequences Without Rational Limits}

Irrationality of $\sqrt{2}$ shows $\Q$ is missing individual points.  But the deeper failure of $\Q$ is \emph{sequential}: you can be moving steadily toward a destination, getting closer and closer to it with every step, and yet the destination simply does not exist inside $\Q$.  This is the incompleteness that Cauchy sequences expose.

The Cauchy condition $|q_m - q_n| < \eps$ for all large $m, n$ captures the intuition of ``the terms are bunching together''.  It is a condition purely on the sequence itself --- it does not mention any limit.  A sequence can be Cauchy without us being able to name what it converges to.  In a \emph{complete} space, being Cauchy is enough to guarantee convergence; in $\Q$, it is not.

Newton's method for $\sqrt{2}$ is the archetypal example.  Starting from $q_1 = 1$, each step $q_{n+1} = (q_n + 2/q_n)/2$ is a rational operation on a rational number, so every term is rational.  The terms get exponentially close to each other (the error squares at each step --- this is quadratic convergence).  Yet the sequence has nowhere to go inside $\Q$, because its limit $\sqrt{2}$ does not live there.

\begin{defbox}{Cauchy Sequence in $\Q$}
$(q_n)$ is \textbf{Cauchy} if $\forall \eps > 0 \in \Q,\; \exists N,\; \forall m, n > N,\; |q_m - q_n| < \eps$.
\end{defbox}
\color{black}

\begin{remarkbox}
\textbf{Historical Note: Bolzano Before Cauchy}

The Cauchy sequence is named after Augustin-Louis Cauchy, who used the convergence criterion in his landmark textbook \emph{Cours d'analyse} (1821). But the concept was understood, at least implicitly, earlier. \textbf{Bernard Bolzano} --- a Bohemian priest and mathematician working in near-total obscurity in Prague --- stated and proved essentially the same criterion in 1817, four years before Cauchy's textbook. Bolzano's work went unread for decades; he is now recognised as one of the founders of rigorous analysis, but only posthumously.

Cauchy himself was not entirely rigorous by modern standards. His 1821 proof that the limit of a convergent series of continuous functions is continuous contains a gap (he implicitly assumed \emph{uniform} convergence, a concept that would only be articulated by Weierstrass decades later). Abel found a counterexample in 1826. This illustrates the iterative nature of mathematical rigour: each generation tightens the screws left loose by the previous one.

The name ``Cauchy sequence'' also obscures an irony: Cauchy believed the real numbers were already well-understood, and never felt the need to \emph{construct} $\R$ from $\Q$. He simply assumed the reals existed and used convergence as a criterion for a sequence to have a limit \emph{in} $\R$. The explicit construction of $\R$ using Cauchy sequences was carried out by \textbf{Georg Cantor} (1872), independently of Dedekind. Cantor was motivated by his work on trigonometric series, where he needed to understand exactly which sets of points could be the ``exceptional sets'' of a Fourier series. This led him directly to the question of what a real number fundamentally is --- and from there to his revolutionary set theory, transfinite cardinals, and one of the most productive and turbulent careers in the history of mathematics.
\end{remarkbox}
\color{black}

\begin{example}
Newton's method for $\sqrt{2}$: $q_1 = 1$, $q_{n+1} = (q_n + 2/q_n)/2$. The sequence
$1,\; 3/2,\; 17/12,\; 577/408,\;\ldots$ is Cauchy in $\Q$ but converges to $\sqrt{2} \notin \Q$.
\end{example}

\begin{thmbox}{$\Q$ is Not Complete}
There exists a Cauchy sequence in $\Q$ with no limit in $\Q$.
\end{thmbox}
\color{black}

\begin{proof}
The Newton sequence above is Cauchy (the error halves quadratically at each step) and its limit satisfies $L^2 = 2$, which has no rational solution.
\end{proof}

\begin{remarkbox}
Cauchy completion generalizes: given any metric space, we can form its completion by the same quotient of Cauchy sequences. This construction is used to define $p$-adic numbers, $L^2$ spaces, and completions of normed rings.
\end{remarkbox}
\color{black}

\begin{exbox}
\begin{enumerate}
  \item Prove $\sqrt{3} \notin \Q$ and generalize: $\sqrt{p} \notin \Q$ for any prime $p$.
  \item Show the sequence $q_n = \sum_{k=0}^n 1/k!$ is Cauchy in $\Q$.
  \item Prove: any convergent sequence in $\Q$ is Cauchy.
  \item Prove: any Cauchy sequence in $\Q$ is bounded.
  \item Define a Dedekind cut and verify that $L = \{q \in \Q : q < 0 \lor q^2 < 2\}$ is a valid cut representing $\sqrt{2}$.
\end{enumerate}
\end{exbox}
\color{black}

%%-----------------------------------------------------------------------------
\chapter{The Real Numbers}
%%-----------------------------------------------------------------------------
\color{black}

\begin{remarkbox}
\textbf{Historical Note: The Year 1872 and the Construction of the Reals}

The year 1872 saw a remarkable simultaneous publication of rigorous constructions of the real numbers. \textbf{Richard Dedekind}'s \emph{Stetigkeit und irrationale Zahlen} (Continuity and Irrational Numbers) introduced Dedekind cuts --- the approach of defining a real number as a partition of $\Q$ into a lower set and an upper set. \textbf{Georg Cantor}'s paper, published the same year, defined real numbers as equivalence classes of Cauchy sequences of rationals --- the approach we follow. \textbf{Karl Weierstrass} had been teaching a similar construction in his Berlin lectures since the 1860s, though he did not publish it.

Dedekind was motivated by a moment of dissatisfaction in the classroom. While teaching calculus in G\"{o}ttingen in 1858, he found himself unable to give a rigorous proof that a monotone bounded sequence converges --- the very theorem we prove in Part~II. He realised the proof ultimately rested on geometric intuition about the ``continuity of the line,'' and set out to make that continuity precise without appealing to geometry. His cuts were the result, written down immediately in 1858 but not published for fourteen years.

The two approaches --- cuts and Cauchy sequences --- produce isomorphic structures, as expected from the uniqueness theorem. Mathematicians have personal preferences: Dedekind cuts are more immediately geometric and make the order structure transparent; Cauchy sequences are more algebraic and generalise directly to any metric space (not just $\Q$). Dedekind's approach is used in most undergraduate analysis textbooks (Hardy, Rudin's \emph{Principles}); Cantor's approach appears in more algebraic treatments and is the natural one for $p$-adic numbers.

A third construction, due to \textbf{Weierstrass}, builds $\R$ as infinite decimal expansions with a careful equivalence relation ($0.999\ldots = 1.000\ldots$). This is the most computationally concrete but the most painful to make rigorous. A fourth approach, due to \textbf{John Conway} (1976), constructs the ``surreal numbers'' --- a vast ordered field containing $\R$, the ordinals, and infinitesimal quantities --- via an elegant game-theoretic definition. Conway's construction is arguably the most beautiful but the least practical for analysis.

\textbf{A note on non-standard analysis.} In 1960, Abraham Robinson showed that it is consistent (relative to ZFC) to have a complete ordered field containing genuine infinitesimals --- numbers $\eps$ with $0 < \eps < 1/n$ for all $n \in \N$. This is called \textbf{non-standard analysis} (or \textbf{hyperreal analysis}), and it vindicates the intuitions of Leibniz, who always thought of $dx$ as a literal infinitesimal rather than a limit. Non-standard analysis gives rigorous proofs of many calculus theorems that are shorter and more intuitive than their standard counterparts, though the logical overhead of setting up the framework is substantial.
\end{remarkbox}
\color{black}

We now close the gaps in $\Q$ by formally adjoining the limits of all Cauchy sequences.

\color{black}\begin{center}
\begin{tikzpicture}[
  set/.style={ellipse, draw=thmblue, fill=white, thick,
              minimum width=#1, minimum height=1.1cm, align=center,
              font=\small\bfseries},
  emb/.style={-{Stealth[length=7pt]}, thick, color=tealaccent},
]
  % Draw ellipses largest-first on background layer so inner labels show through
  \begin{pgfonlayer}{background}
    \node[set=6.4cm, fill=exbg, draw=exborder]  (R) at (0,0) {};
    \node[set=4.4cm]  (Q) at (0,0) {};
    \node[set=2.8cm]  (Z) at (0,0) {};
    \node[set=1.4cm]  (N) at (0,0) {};
  \end{pgfonlayer}
  % Labels drawn on main layer, on top of all ellipses
  \node[font=\small\bfseries] at ( 0.00, 0) {\textcolor{navyblue}{$\N$}};
  \node[font=\small\bfseries] at ( 0.85, 0) {\textcolor{navyblue}{$\Z$}};
  \node[font=\small\bfseries] at ( 1.65, 0) {\textcolor{navyblue}{$\Q$}};
  \node[font=\small\bfseries] at ( 2.65, 0) {\textcolor{navyblue}{$\R$}};

  % Embedding arrows (between label positions)
  \draw[emb] (0.42, -0.75) -- node[below, font=\scriptsize\itshape]
        {\textcolor{tealaccent}{$n \mapsto [(n,0)]$}} (0.82, -0.75);
  \draw[emb] (1.22, -0.5) -- node[below, font=\scriptsize\itshape]
        {\textcolor{tealaccent}{$a/b \mapsto [(a,b)]$}} (1.62, -0.5);
  \draw[emb] (2.05, -0.25) -- node[below, font=\scriptsize\itshape]
        {\textcolor{tealaccent}{$q \mapsto [(q,q,\ldots)]$}} (2.55,-0.25);

  % Gap label
  \node[font=\scriptsize, align=center] at (3.8, 0.9)
        {\textcolor{red!70}{$\sqrt{2}$ lives here}\\
         \textcolor{red!70}{(gap in $\Q$)}};
  \draw[->, red!60, dashed] (3.5, 0.65) -- (2.8, 0.15);
\end{tikzpicture}
\end{center}

\section{Cauchy Completion of \texorpdfstring{$\Q$}{Q}}

The Cauchy completion is a bold move: instead of hunting for the missing points and trying to insert them one by one, we \emph{declare} that every Cauchy sequence \emph{is} a real number.  If the sequence $(1, 1.4, 1.41, 1.414, \ldots)$ wants to converge somewhere but $\Q$ has no such place, we simply \emph{create} a new object and call it $\sqrt{2}$.  Every Cauchy sequence that ``should'' converge to the same value defines the same real number.

This is the third and final quotient construction in Part~I.  The pattern is always the same:
\begin{enumerate}[noitemsep]
  \item Start with a large set of ``candidates'' (pairs of naturals; pairs of integers; Cauchy sequences).
  \item Declare an equivalence relation that collapses redundant candidates (same difference; same ratio; same limit).
  \item Define the new number system as the quotient.
  \item Lift arithmetic to equivalence classes and check well-definedness.
\end{enumerate}

The one genuinely new ingredient here, compared to Chapters~2 and~3, is that the candidate set $\mathcal{C}$ is \emph{infinite-dimensional}: its elements are sequences, not pairs.  The well-definedness checks therefore involve limit arguments (show that sums and products of Cauchy sequences are Cauchy), not just algebraic identities.

\begin{defbox}{Equivalence of Cauchy Sequences}
Two Cauchy sequences $(a_n)$ and $(b_n)$ in $\Q$ are equivalent---written $(a_n) \sim (b_n)$---if:
\[ \forall \eps > 0 \in \Q,\; \exists N,\; \forall n > N,\; |a_n - b_n| < \eps. \]
The real numbers are $\R := \mathcal{C}/{\sim}$, where $\mathcal{C}$ is the set of all Cauchy sequences in $\Q$.
\end{defbox}
\color{black}

\begin{thmbox}{$\sim$ is an Equivalence Relation on $\mathcal{C}$}
\textbf{Proof.}
\textbf{Refl:} $|a_n - a_n| = 0 < \eps$ for all $n$.
\textbf{Symm:} $|b_n - a_n| = |a_n - b_n| < \eps$.
\textbf{Trans:} $|a_n - c_n| \leq |a_n - b_n| + |b_n - c_n|$; apply each
hypothesis with $\eps/2$. \qed
\end{thmbox}
\color{black}

\begin{defbox}{Arithmetic on $\R$}
$[(a_n)] + [(b_n)] := [(a_n + b_n)]$ \quad and \quad $[(a_n)] \cdot [(b_n)] := [(a_n b_n)]$.

The embedding $\Q \hookrightarrow \R$ sends $q$ to the constant sequence $[(q, q, q, \ldots)]$.
\end{defbox}
\color{black}

\begin{thmbox}{Products and Sums of Cauchy Sequences Are Cauchy}
If $(a_n)$ and $(b_n)$ are Cauchy in $\Q$, so are $(a_n + b_n)$ and $(a_n b_n)$.
\end{thmbox}
\color{black}

\begin{proof}
\textit{Sum:} Triangle inequality with $\eps/2$. \textit{Product:} Use boundedness (Cauchy $\Rightarrow$ bounded, shown separately). With $|a_n| \leq M_a$ and $|b_n| \leq M_b$:
\[ |a_m b_m - a_n b_n| \leq M_a|b_m - b_n| + M_b|a_m - a_n|. \]
Choose $N$ so both terms are $< \eps/2$.
\end{proof}

\begin{leanbox}
\begin{lstlisting}
-- Our real numbers as a quotient of Cauchy sequences over Rat
-- (Using Lean's built-in Rat for the field of fractions)
def IsCauchy (a : Nat -> Rat) : Prop :=
  forall eps : Rat, eps > 0 ->
    exists N : Nat, forall m n : Nat, m > N -> n > N ->
      |a m - a n| < eps

structure CauchySeq where
  seq    : Nat -> Rat
  cauchy : IsCauchy seq

def CauchyEquiv (a b : CauchySeq) : Prop :=
  forall eps : Rat, eps > 0 ->
    exists N : Nat, forall n : Nat, n > N ->
      |a.seq n - b.seq n| < eps

-- MyReal is the quotient
def MyReal : Type := Quot CauchyEquiv

-- Constant sequence embeds Q into R
def ofRat (q : Rat) : MyReal :=
  Quot.mk CauchyEquiv
    { seq    := fun _ => q,
      cauchy := fun eps heps => ⟨0, fun m n _ _ => by simp; exact heps⟩ }

-- Addition
def addCauchy (a b : CauchySeq) : CauchySeq :=
  { seq    := fun n => a.seq n + b.seq n,
    cauchy := fun eps heps => by
      obtain ⟨Na, ha⟩ := a.cauchy (eps/2) (by linarith)
      obtain ⟨Nb, hb⟩ := b.cauchy (eps/2) (by linarith)
      refine ⟨max Na Nb, fun m n hm hn => ?_⟩
      have hma := Nat.lt_of_lt_of_le hm (Nat.le_max_left Na Nb)
      have hmb := Nat.lt_of_lt_of_le hm (Nat.le_max_right Na Nb)
      have hna := Nat.lt_of_lt_of_le hn (Nat.le_max_left Na Nb)
      have hnb := Nat.lt_of_lt_of_le hn (Nat.le_max_right Na Nb)
      have h1 := ha m n hma hna
      have h2 := hb m n hmb hnb
      have : |(a.seq m + b.seq m) - (a.seq n + b.seq n)| <=
             |a.seq m - a.seq n| + |b.seq m - b.seq n| := by
        calc |(a.seq m + b.seq m) - (a.seq n + b.seq n)|
            = |(a.seq m - a.seq n) + (b.seq m - b.seq n)| := by ring_nf
          _ <= |a.seq m - a.seq n| + |b.seq m - b.seq n| := abs_add _ _
      linarith }

def addReal : MyReal -> MyReal -> MyReal :=
  Quot.lift2
    (fun a b => Quot.mk CauchyEquiv (addCauchy a b))
    (by
      intro a1 a2 b h
      apply Quot.sound
      intro eps heps
      obtain ⟨N, hN⟩ := h eps heps
      exact ⟨N, fun n hn => by simp [addCauchy]; linarith [hN n hn]⟩)
\end{lstlisting}
\end{leanbox}
\color{black}

\section{Completeness and Key Properties of \texorpdfstring{$\R$}{R}}

We built $\R$ specifically to be complete, so the first thing to verify is that we succeeded.  Completeness of $\R$ --- every Cauchy sequence converges --- follows almost by definition: a Cauchy sequence in $\R$ is a sequence of equivalence classes of Cauchy sequences in $\Q$, and we can always ``diagonalise'' to extract a single Cauchy sequence in $\Q$ that represents the limit.

The \textbf{least upper bound property} (every nonempty bounded-above set has a supremum) is equivalent to completeness, and is often taken as the definition of completeness in analysis textbooks.  We derive it here from the Cauchy-sequence version using a bisection argument.  This property is what powers the proofs of the Intermediate Value Theorem, the Extreme Value Theorem, and the existence of definite integrals --- essentially every deep theorem of calculus ultimately rests on it.

The \textbf{Archimedean property} says $\R$ has no ``infinitely large'' natural numbers and no ``infinitely small'' positive reals.  This might seem obvious but it is not a consequence of the ordered field axioms alone --- there exist non-Archimedean ordered fields.  In $\R$ it follows from the least upper bound property: if $\N$ were bounded above, it would have a supremum, leading immediately to a contradiction.

Finally, a note on uniqueness: the complete ordered field axioms characterise $\R$ up to isomorphism.  Any two complete ordered fields are isomorphic.  This means the specific construction we used --- Cauchy sequences of rationals --- is one of many possible routes to the same destination.  Dedekind cuts lead to the same place.  So does any other sensible completion of $\Q$.  This is reassuring: the real numbers are not an artifact of our particular construction but a mathematical object with an intrinsic identity.

\medskip
\textbf{The constructivist objection.} Brouwer and his followers reject the least upper bound property as stated, because its proof requires the \emph{law of excluded middle}: either a given rational $q$ is an upper bound for $S$ or it is not. A constructivist demands a \emph{decision procedure} --- an algorithm that, given any $q$, tells you which case applies. For a general subset $S \subseteq \R$ defined by an arbitrary predicate, no such algorithm may exist.

Bishop's constructive analysis resolves this by weakening the statement: instead of ``every bounded-above set has a supremum,'' one proves ``every \emph{located} set'' (one for which the distance from any point can be computed) has a supremum. This is a strictly stronger hypothesis and gives a strictly weaker conclusion, but it is computable. The resulting theory is consistent with classical analysis (every classically valid theorem is also constructively valid if you add enough hypotheses), but proofs are longer.

For our purposes, we work classically throughout. The Lean~4 theorems in this book use \texttt{Classical.em} (law of excluded middle) implicitly wherever we appeal to the LUB property or prove existence by contradiction. If you wanted a fully constructive development, you would need to replace our \texttt{MyReal} with a type of \emph{Cauchy reals with modulus} --- sequences equipped with an explicit convergence rate function --- and all existence proofs would need to be replaced by explicit constructions. This is the approach taken in Lean's \texttt{Mathlib.Topology.Algebra.Order} when it works in the constructive fragment.

\begin{thmbox}{$\R$ is Complete}
Every Cauchy sequence in $\R$ converges to an element of $\R$.

\textbf{Proof sketch.} Let $(x_k)$ be Cauchy in $\R$, where each $x_k = [(a_n^{(k)})_n]$. For each $k$, choose a ``diagonal'' rational $r_k$ within $1/k$ of $x_k$ (possible since the Cauchy sequence for $x_k$ eventually stays close to its limit). Then $(r_k)$ is Cauchy in $\Q$, and $[(r_k)] = \lim_{k \to \infty} x_k$ in $\R$. \qed
\end{thmbox}

\begin{thmbox}{The Least Upper Bound Property}
Every nonempty subset of $\R$ that is bounded above has a \textbf{supremum} (least upper bound) in $\R$.
\end{thmbox}
\color{black}

\begin{proof}[Proof sketch]
This follows from completeness. Given $S \subseteq \R$ bounded above by $M$: pick any $a_0 \in S$ and set $b_0 = M$. Inductively, let $m_k = (a_k + b_k)/2$. If some element of $S$ exceeds $m_k$, set $a_{k+1} = m_k$, $b_{k+1} = b_k$; otherwise $a_{k+1} = a_k$, $b_{k+1} = m_k$. The $a_k$ form a Cauchy sequence whose limit is $\sup S$.
\end{proof}

\begin{remarkbox}
\textbf{Classical Logic and the Least Upper Bound}

The completeness of $\R$---often stated as ``every Cauchy sequence converges'' or equivalently ``every bounded-above set has a least upper bound''---is \emph{not} a theorem but an axiom that distinguishes $\R$ from $\Q$. We have \emph{constructed} $\R$ to have this property. The remarkable fact is that this property, together with the ordered field axioms, uniquely determines $\R$ up to isomorphism.
\end{remarkbox}
\color{black}

\begin{thmbox}{The Archimedean Property of $\R$}
For any $x \in \R$ with $x > 0$, there exists $n \in \N$ with $1/n < x$.
Equivalently, $\N$ is not bounded above in $\R$.
\end{thmbox}
\color{black}

\begin{proof}
If $\N$ were bounded above, by the LUB property it would have a supremum $s$. But $s - 1/2$ is not an upper bound (since $s$ is the least), so some $n > s - 1/2$, giving $n + 1 > s$---contradicting $s$ being an upper bound.
\end{proof}

\begin{exbox}
\begin{enumerate}
  \item Prove that $\R$ is uncountable (Cantor's diagonal argument).
  \item Prove the \textbf{nested interval theorem}: if $[a_n, b_n]$ is a nested sequence of closed intervals with $b_n - a_n \to 0$, then $\bigcap [a_n, b_n]$ is a single point.
  \item Show the ordering on $\R$ is well-defined: $[(a_n)] > 0$ iff there exists $\eps \in \Q$, $\eps > 0$, and $N$ such that $a_n > \eps$ for all $n > N$.
  \item Prove that $\Q$ is dense in $\R$: for any $x < y$ in $\R$, there exists $q \in \Q$ with $x < q < y$.
  \item Implement \li{mulReal} in Lean~4, proving well-definedness with respect to \li{CauchyEquiv}.
\end{enumerate}
\end{exbox}
\color{black}

%%=============================================================================
\part{Topology of \texorpdfstring{$\R$}{R} and Limits}
%%=============================================================================
\color{black}

%%-----------------------------------------------------------------------------
\chapter{Sequences and Convergence}
%%-----------------------------------------------------------------------------
\color{black}

With $\R$ in hand, we can give rigorous meaning to ``approaching a limit.'' This chapter establishes the language of sequences and convergence that underlies all of calculus.

\begin{remarkbox}
\textbf{Historical Note: The Arithmetisation of Analysis}

For most of the eighteenth century, analysis was built on intuitive notions of ``infinitely small'' quantities. Newton called them \emph{fluxions}; Leibniz called them \emph{differentials}. Both inventors knew these objects were problematic --- Newton hedged by speaking of ``ultimate ratios'' and ``first and last quantities,'' Leibniz by treating $dx$ as a formal symbol that obeyed algebraic rules. Bishop Berkeley savaged both approaches in 1734 in \emph{The Analyst}, mocking infinitesimals as ``the ghosts of departed quantities.''

The response, over the following century, was the \textbf{arithmetisation of analysis}: the project of rebuilding calculus entirely in terms of numbers and their properties, without appealing to geometric intuition or infinitesimal quantities. The heroes of this programme are three: \textbf{Bolzano} (1817), who first stated the rigorous definition of convergence and continuity; \textbf{Cauchy} (1821--1823), whose three textbooks --- \emph{Cours d'analyse}, \emph{R\'{e}sum\'{e} des le\c{c}ons}, and \emph{Le\c{c}ons sur le calcul diff\'{e}rentiel} --- gave the first systematic $\varepsilon$-$\delta$ treatment of limits, continuity, and derivatives; and \textbf{Weierstrass} (1860s--1880s), whose Berlin lectures completed the programme by filling the gaps Cauchy had left (uniform convergence, the pathological functions, and the rigorous construction of $\R$).

The $\varepsilon$-$N$ definition we use in this chapter is essentially Cauchy's, rendered in modern notation by Weierstrass. The notation $\varepsilon$ for the ``error bound'' and $\delta$ for the input tolerance appear to have been introduced by Weierstrass; earlier writers used prose descriptions. The symbol $\lim$ for limit is due to Simon L'Huilier (1786), though it was not universally adopted until the early twentieth century.

An important subtlety: Cauchy's definition of limit was slightly different from ours. He defined $\lim_{n\to\infty} a_n = L$ to mean ``$a_n - L$ becomes infinitely small,'' which required a separate (and never fully rigorous) account of what ``infinitely small'' meant. Weierstrass replaced this with the purely quantitative $\forall \varepsilon > 0, \exists N, \forall n > N, |a_n - L| < \varepsilon$, which contains no undefined terms at all. This is the version in our Definition box below.
\end{remarkbox}
\color{black}

\section{Convergence of Sequences}

With $\R$ constructed and its completeness established, we can now give rigorous meaning to the phrase ``approaching a value.''  Informally, a sequence converges to $L$ if the terms eventually stay as close to $L$ as we like.  The $\varepsilon$-$N$ definition formalises this: $\varepsilon$ is the precision we demand, and $N$ is the threshold beyond which the sequence delivers it.  The quantifier order matters: $\forall \varepsilon,\, \exists N$ means the threshold $N$ is allowed to depend on how tight a precision we ask for.  The smaller the $\varepsilon$, the larger $N$ may need to be.

\color{black}\begin{center}
\begin{tikzpicture}[xscale=0.7, yscale=1.1]
  % Axes
  \draw[->] (-0.3,0) -- (10.5,0) node[right, font=\small] {\textcolor{black}{$n$}};
  \draw[->] (0,-0.3) -- (0,3.2) node[above, font=\small] {\textcolor{black}{$a_n$}};

  % Limit line L
  \def\L{2.0}
  \def\eps{0.55}
  \draw[dashed, color=thmblue, thick] (-0.1,\L) -- (10.3,\L)
    node[right, font=\small] {\textcolor{thmblue}{$L$}};

  % eps band
  \fill[tealaccent!12] (-0.1, \L-\eps) rectangle (10.3, \L+\eps);
  \draw[dashed, color=tealaccent] (-0.1,\L+\eps) -- (10.3,\L+\eps)
    node[right, font=\scriptsize] {\textcolor{tealaccent}{$L+\varepsilon$}};
  \draw[dashed, color=tealaccent] (-0.1,\L-\eps) -- (10.3,\L-\eps)
    node[right, font=\scriptsize] {\textcolor{tealaccent}{$L-\varepsilon$}};

  % Sequence terms (oscillating, settling)
  \foreach \x/\y in {
    0.5/0.4, 1/3.0, 1.5/0.8, 2/2.8, 2.5/0.6, 3/2.6,
    3.5/1.0, 4/2.5, 4.5/1.3, 5/2.35, 5.5/1.55,
    6/2.2, 6.5/1.75, 7/2.15, 7.5/1.88, 8/2.08,
    8.5/1.95, 9/2.04, 9.5/1.98, 10/2.01
  }{
    \fill[thmblue] (\x, \y) circle (2.4pt);
  }

  % N threshold
  \def\Nval{6.0}
  \draw[-, color=red!65, thick, dashed] (\Nval,-0.15) -- (\Nval, 3.1);
  \node[below, font=\scriptsize] at (\Nval, -0.15) {\textcolor{red!70}{$N$}};
  \node[font=\scriptsize, align=center] at (8.3, 3.0)
    {\textcolor{red!70}{$|a_n - L| < \varepsilon$}\\
     \textcolor{red!70}{for all $n > N$}};
  \draw[->, red!50, thin] (7.5, 2.78) -- (8.2, 2.25);

  % Brace for eps
  \draw[decorate, decoration={brace,amplitude=5pt}, color=tealaccent]
    (10.6, \L) -- (10.6, \L+\eps) node[midway, right=4pt, font=\scriptsize] {\textcolor{tealaccent}{$\varepsilon$}};
\end{tikzpicture}
\end{center}

\begin{defbox}{Convergence of a Sequence}
A sequence $(a_n)_{n=1}^\infty$ in $\R$ \textbf{converges} to $L \in \R$, written $a_n \to L$ or $\lim_{n\to\infty} a_n = L$, if:
\[ \forall \eps > 0,\; \exists N \in \N,\; \forall n > N,\; |a_n - L| < \eps. \]
A sequence that converges to some $L$ is \textbf{convergent}; otherwise \textbf{divergent}.
\end{defbox}
\color{black}

The definition unwinds as: ``for every desired precision $\eps$, there is a threshold $N$ beyond which the sequence stays within $\eps$ of $L$.'' The key point is that $N$ may depend on $\eps$---smaller $\eps$ may require larger $N$.

\begin{thmbox}{Limits Are Unique}
If $a_n \to L$ and $a_n \to M$, then $L = M$.
\end{thmbox}
\color{black}

\begin{proof}
Suppose $L \neq M$; set $\eps = |L - M|/2 > 0$. Choose $N_1$ so $|a_n - L| < \eps$ for $n > N_1$ and $N_2$ so $|a_n - M| < \eps$ for $n > N_2$. For $n > \max(N_1, N_2)$:
\[ |L - M| \leq |L - a_n| + |a_n - M| < \eps + \eps = |L - M|. \]
Contradiction.
\end{proof}

\begin{leanbox}
\begin{lstlisting}
def SeqLimit (a : Nat -> Real) (L : Real) : Prop :=
  forall eps : Real, eps > 0 ->
    exists N : Nat, forall n : Nat, n > N -> |a n - L| < eps

-- Limits are unique
theorem limit_unique (a : Nat -> Real) (L M : Real)
    (hL : SeqLimit a L) (hM : SeqLimit a M) : L = M := by
  by_contra h
  have hne : L ≠ M := h
  set eps := |L - M| / 2
  have heps : eps > 0 := by positivity
  obtain ⟨N1, hN1⟩ := hL eps heps
  obtain ⟨N2, hN2⟩ := hM eps heps
  have hn := max N1 N2 + 1
  have h1 := hN1 (max N1 N2 + 1) (by omega)
  have h2 := hN2 (max N1 N2 + 1) (by omega)
  have : |L - M| < |L - M| := by
    calc |L - M| = |L - a _ + (a _ - M)| := by ring_nf
      _ <= |L - a _| + |a _ - M| := abs_add _ _
      _ < eps + eps := by linarith
      _ = |L - M| := by ring
  linarith
\end{lstlisting}
\end{leanbox}
\color{black}

\section{Algebra of Limits}

Once we know individual sequences converge, we want to know whether building new sequences from old ones preserves convergence.  The algebra of limits says yes: limits interact well with the four arithmetic operations.  This is the machinery that will let us compute derivatives and integrals without returning to $\varepsilon$-$N$ arguments every time --- we reduce everything to known limits by algebraic manipulation.

\begin{thmbox}{Algebra of Limits}
If $a_n \to L$ and $b_n \to M$ in $\R$, then:
\begin{enumerate}[noitemsep]
  \item $a_n + b_n \to L + M$
  \item $c \cdot a_n \to cL$ for any constant $c$
  \item $a_n \cdot b_n \to L \cdot M$
  \item $a_n / b_n \to L/M$ provided $M \neq 0$ (and $b_n \neq 0$ eventually)
\end{enumerate}
\end{thmbox}
\color{black}

\begin{proof}
\textbf{(1):} Given $\eps > 0$, choose $N_1, N_2$ so that $|a_n - L| < \eps/2$ and $|b_n - M| < \eps/2$ for $n > N_1, N_2$ respectively. For $n > \max(N_1, N_2)$:
\[ |(a_n + b_n) - (L + M)| \leq |a_n - L| + |b_n - M| < \eps. \]

\textbf{(3):} Write $a_n b_n - LM = (a_n - L)b_n + L(b_n - M)$. Use boundedness of $(b_n)$ and convergence of both sequences.

\textbf{(4):} Reduce to showing $1/b_n \to 1/M$. Since $M \neq 0$, for large $n$ we have $|b_n| \geq |M|/2 > 0$, and:
\[ \abs{\frac{1}{b_n} - \frac{1}{M}} = \frac{|b_n - M|}{|b_n||M|} \leq \frac{2}{|M|^2}|b_n - M|. \]
\end{proof}

\section{Important Limit Theorems}

Three theorems deserve special attention because they do not follow directly from the algebra of limits but are used constantly.  The Squeeze Theorem handles limits we cannot compute directly by trapping the sequence between two simpler ones.  The Monotone Convergence Theorem is the sequential counterpart of the least upper bound property --- it is the reason that bounded increasing sequences always converge in $\R$ but not in $\Q$.  Bolzano--Weierstrass is the compactness statement for sequences: any bounded sequence, however wildly it oscillates, must have a convergent subsequence hiding inside it.

\begin{thmbox}{Squeeze Theorem}
If $a_n \leq c_n \leq b_n$ for all $n$ and $a_n \to L$, $b_n \to L$, then $c_n \to L$.
\end{thmbox}
\color{black}

\begin{proof}
Given $\eps > 0$, choose $N$ so $|a_n - L| < \eps$ and $|b_n - L| < \eps$ for $n > N$. Then $L - \eps < a_n \leq c_n \leq b_n < L + \eps$, so $|c_n - L| < \eps$.
\end{proof}

\begin{thmbox}{Monotone Convergence Theorem}
A bounded, monotone sequence in $\R$ converges. Specifically:
\begin{itemize}[noitemsep]
  \item An increasing sequence bounded above converges to its supremum.
  \item A decreasing sequence bounded below converges to its infimum.
\end{itemize}
\end{thmbox}
\color{black}

\begin{proof}
Let $(a_n)$ be increasing and bounded above. By the LUB property, $L := \sup\{a_n\}$ exists in $\R$. Given $\eps > 0$, since $L - \eps$ is not an upper bound, there exists $N$ with $a_N > L - \eps$. For $n > N$: $L - \eps < a_N \leq a_n \leq L$, so $|a_n - L| < \eps$.
\end{proof}

\begin{thmbox}{Bolzano--Weierstrass Theorem}
Every bounded sequence in $\R$ has a convergent subsequence.
\end{thmbox}
\color{black}

\begin{proof}
Let $(a_n)$ be bounded by $[m, M]$. Use the bisection method: at each step, one half of the current interval contains infinitely many terms of the sequence; pick any one as the next subsequence element. This produces a Cauchy subsequence, which converges by completeness.
\end{proof}

\begin{exbox}
\begin{enumerate}
  \item Prove parts (2) and (4) of the Algebra of Limits theorem.
  \item Show that $1/n^2 \to 0$ directly from the definition.
  \item Prove that any convergent sequence is bounded.
  \item Use the Squeeze Theorem to show $\lim_{n\to\infty} \sin(n)/n = 0$.
  \item Prove the Bolzano--Weierstrass theorem in full detail (the bisection argument sketched above).
  \item Formalize the Monotone Convergence Theorem in Lean~4 using \li{SeqLimit}.
\end{enumerate}
\end{exbox}
\color{black}

%%-----------------------------------------------------------------------------
\chapter{Limits of Functions}
%%-----------------------------------------------------------------------------
\color{black}

\section{The \texorpdfstring{$\eps$--$\del$}{epsilon-delta} Definition}

Sequential limits described convergence of infinite lists of numbers.  Functional limits describe what happens to $f(x)$ as $x$ moves continuously toward a point $a$.  The definition is deliberately agnostic about the value of $f$ at $a$ itself --- limits capture approach, not arrival.  This separation between the limit and the function value is what makes the concept useful: it lets us study where a function is \emph{heading} even at points where it is undefined or discontinuous.

\color{black}\begin{center}
\begin{tikzpicture}[xscale=1.5, yscale=1.3]
  % Axes
  \draw[->] (-0.2,0) -- (4.0,0) node[right, font=\small] {\textcolor{black}{$x$}};
  \draw[->] (0,-0.2) -- (0,3.0) node[above, font=\small] {\textcolor{black}{$f(x)$}};

  % a and L
  \def\a{2.0}
  \def\L{1.8}
  \def\del{0.55}
  \def\eps{0.5}

  % delta band on x-axis (shaded region)
  \fill[tealaccent!15] (\a-\del, -0.05) rectangle (\a+\del, 2.85);
  \fill[exbg] (\L-\eps+0.02, -0.05) rectangle (\L+\eps-0.02, 2.85);
  % (overdraw to show eps band on y-axis properly)

  % eps band on y-axis
  \fill[exbg] (-0.05, \L-\eps) rectangle (3.9, \L+\eps);
  \fill[tealaccent!15] (\a-\del,-0.05) rectangle (\a+\del, 3.0);

  % Grid lines for eps and delta
  \draw[dashed, color=tealaccent, thick] (\a-\del, -0.05) -- (\a-\del, 2.85);
  \draw[dashed, color=tealaccent, thick] (\a+\del, -0.05) -- (\a+\del, 2.85);
  \draw[dashed, color=exborder, thick]  (-0.05, \L-\eps) -- (3.9, \L-\eps);
  \draw[dashed, color=exborder, thick]  (-0.05, \L+\eps) -- (3.9, \L+\eps);

  % Function curve
  \draw[very thick, color=thmblue, domain=0.2:3.8, samples=80]
    plot (\x, {1.8 + 0.35*(\x-2) + 0.4*sin(180*(\x-2)/1.2)});

  % Hole at a
  \fill[white] (\a, \L) circle (3.5pt);
  \draw[thmblue, thick] (\a, \L) circle (3.5pt);

  % Labels on axes
  \draw[thick] (\a, 0.04) -- (\a, -0.04)
    node[below, font=\small] {\textcolor{black}{$a$}};
  \draw[thick] (-0.04, \L) -- (0.04, \L)
    node[left, font=\small] {\textcolor{black}{$L$}};

  % delta label
  \draw[decorate, decoration={brace,amplitude=4pt, mirror}, color=tealaccent]
    (\a, -0.18) -- (\a+\del, -0.18)
    node[midway, below=4pt, font=\scriptsize] {\textcolor{tealaccent}{$\delta$}};

  % eps label
  \draw[decorate, decoration={brace,amplitude=4pt}, color=exborder]
    (3.95, \L) -- (3.95, \L+\eps)
    node[midway, right=4pt, font=\scriptsize] {\textcolor{exborder}{$\varepsilon$}};

  % "f(a) not needed" annotation
  \node[font=\scriptsize\itshape, align=center] at (2.0, 2.65)
    {\textcolor{red!60}{$f(a)$ need not}\\
     \textcolor{red!60}{be defined}};
  \draw[->, red!50, thin] (2.0, 2.45) -- (2.02, \L+0.12);
\end{tikzpicture}
\end{center}

\begin{defbox}{Limit of a Function at a Point}
Let $f : \R \to \R$ and $a, L \in \R$. We say $\lim_{x \to a} f(x) = L$ if:
\[ \forall \eps > 0,\; \exists \del > 0,\; \forall x \in \R,\; 0 < |x - a| < \del \Rightarrow |f(x) - L| < \eps. \]
Note that $f(a)$ need not be defined, and even if it is, the limit may differ from $f(a)$.
\end{defbox}
\color{black}

\begin{leanbox}
\begin{lstlisting}
def FuncLimit (f : Real -> Real) (a L : Real) : Prop :=
  forall eps : Real, eps > 0 ->
    exists delta : Real, delta > 0 /\
      forall x : Real, 0 < |x - a| -> |x - a| < delta -> |f x - L| < eps

-- Example: lim_{x -> 2} (x^2) = 4
-- We need: for |x - 2| < delta, |x^2 - 4| < eps
-- |x^2 - 4| = |x+2||x-2|.  If |x-2| < 1 then |x| < 3, |x+2| < 5.
-- So take delta = min(1, eps/5).
example : FuncLimit (fun x => x * x) 2 4 := by
  intro eps heps
  refine ⟨min 1 (eps / 5), by positivity, fun x hxpos hxd => ?_⟩
  have hx2 : |x - 2| < 1     := lt_of_lt_of_le hxd (min_le_left _ _)
  have hx5 : |x - 2| < eps/5 := lt_of_lt_of_le hxd (min_le_right _ _)
  have hbound : |x + 2| < 5 := by
    have : |x| < 3 := by linarith [abs_lt.mp hx2]
    linarith [abs_le.mp (le_of_lt (abs_lt.mpr (by linarith)))]
  calc |x * x - 4| = |(x - 2) * (x + 2)| := by ring_nf
    _ = |x - 2| * |x + 2| := abs_mul _ _
    _ < (eps / 5) * 5 := by nlinarith
    _ = eps := by ring
\end{lstlisting}
\end{leanbox}
\color{black}

\section{One-sided Limits and Limits at Infinity}

The standard limit requires $x$ to approach $a$ from both sides simultaneously.  Sometimes a function behaves differently from the left and the right --- the absolute value $|x|/x$ is $-1$ for $x < 0$ and $+1$ for $x > 0$, so its one-sided limits at $0$ differ.  Limits at infinity formalise the notion of a horizontal asymptote: $f(x) \to L$ as $x \to \infty$ means the graph of $f$ flattens out toward the line $y = L$.

\begin{defbox}{One-sided Limits}
$\lim_{x \to a^+} f(x) = L$ (right-hand limit): replace $0 < |x-a| < \del$ by $0 < x - a < \del$.

$\lim_{x \to a^-} f(x) = L$ (left-hand limit): replace by $0 < a - x < \del$.

$\lim_{x \to \infty} f(x) = L$: $\forall \eps > 0,\; \exists M,\; \forall x > M,\; |f(x) - L| < \eps$.
\end{defbox}
\color{black}

\begin{thmbox}{Sequential Characterization of Limits}
$\lim_{x \to a} f(x) = L$ if and only if for every sequence $(x_n)$ with $x_n \to a$ and $x_n \neq a$ for all $n$, we have $f(x_n) \to L$.
\end{thmbox}
\color{black}

\begin{proof}
$(\Rightarrow)$: If $f(x) \to L$ and $x_n \to a$ ($x_n \neq a$), given $\eps > 0$ choose $\del$ for the limit, then $N$ so $|x_n - a| < \del$ for $n > N$; then $|f(x_n) - L| < \eps$.

$(\Leftarrow)$: Suppose the limit does not hold. Then there exists $\eps > 0$ such that for every $\del > 0$, there exists $x$ with $0 < |x - a| < \del$ and $|f(x) - L| \geq \eps$. Taking $\del = 1/n$ gives a sequence $x_n \to a$ with $|f(x_n) - L| \geq \eps$ for all $n$---contradicting $f(x_n) \to L$.
\end{proof}

\section{Algebra of Functional Limits}

The algebra of functional limits is a direct corollary of the sequential version: we proved limits behave well under arithmetic for sequences, and the sequential characterisation of functional limits (the next theorem) transfers this immediately.  The key bridge theorem below is worth internalising: it converts every question about functional limits into a question about sequential limits, making the two notions equivalent and letting us move freely between them.

\begin{thmbox}{Algebra of Functional Limits}
If $\lim_{x \to a} f(x) = L$ and $\lim_{x \to a} g(x) = M$, then:
$(f + g)(x) \to L + M$, $(fg)(x) \to LM$, $(f/g)(x) \to L/M$ (if $M \neq 0$), all as $x \to a$.
\end{thmbox}
\color{black}

\begin{proof}
Reduce to the sequential characterization and apply the Algebra of Sequential Limits from Chapter~6.
\end{proof}

\begin{exbox}
\begin{enumerate}
  \item Prove from the definition: $\lim_{x \to 3} (2x + 1) = 7$.
  \item Prove that $\lim_{x \to 0} \sin(x)/x = 1$ using the squeeze theorem (hint: geometric argument gives $\cos\theta \leq \sin\theta/\theta \leq 1$ for $\theta \in (0, \pi/2)$).
  \item Prove the sequential characterization of limits in full detail.
  \item Show that $\lim_{x \to 0} |x|/x$ does not exist by finding sequences approaching $0$ from different sides.
  \item Prove: if $\lim_{x \to a} f(x) = L > 0$, then $f(x) > 0$ for all $x$ sufficiently close to $a$.
\end{enumerate}
\end{exbox}
\color{black}

%%-----------------------------------------------------------------------------
\chapter{Continuity}
%%-----------------------------------------------------------------------------
\color{black}

\begin{remarkbox}
\textbf{Historical Note: Continuity and the Pathological Functions}

The nineteenth century produced a series of increasingly disturbing examples that forced mathematicians to rethink their intuitions about continuity.

\textbf{Fourier} (1822) showed that a discontinuous function --- one with jumps --- could be represented as an infinite sum of continuous sine and cosine functions. This raised an immediate question: if the \emph{limit} of continuous functions is not necessarily continuous, what exactly does ``continuous'' mean? Cauchy claimed (incorrectly) that the uniform limit of continuous functions is continuous. \textbf{Abel} found a counterexample in 1826; \textbf{Seidel} and \textbf{Stokes} independently clarified the issue in 1847--48 by distinguishing uniform from pointwise convergence.

The next shock came from \textbf{Dirichlet} (1829), who defined the function that is $1$ on rationals and $0$ on irrationals. This function is nowhere continuous --- the rationals and irrationals are both dense --- yet it is perfectly well-defined pointwise. It is not Riemann integrable (Riemann's criterion requires the oscillation to be controllable on each subinterval, which fails everywhere for Dirichlet's function). It \emph{is} Lebesgue integrable, with integral $0$, since the rationals form a set of measure zero.

The most spectacular example came from \textbf{Weierstrass} in 1872 (announced to the Berlin Academy; published 1875): a function $f(x) = \sum_{n=0}^\infty a^n \cos(b^n \pi x)$ where $0 < a < 1$, $b$ is a positive odd integer, and $ab > 1 + 3\pi/2$. This function is \emph{continuous everywhere} but \emph{differentiable nowhere}. Weierstrass's construction demolished the intuition that a continuous curve must have a tangent ``almost everywhere.'' Hermite called it ``a lamentable scourge'' (\emph{un fl\'{e}au d\'{e}plorable}). Poincar\'{e} complained that such functions were ``invented to show that our fathers' reasoning was at fault.'' Today we know that, in a precise sense, the ``typical'' continuous function is nowhere differentiable --- the differentiable functions are a small exception.

These examples are what forced the rigorous $\varepsilon$-$\delta$ definitions in this chapter. Without them, you cannot even \emph{state} precisely what it means to be continuous, let alone prove that Dirichlet's function is discontinuous everywhere or that Weierstrass's function is continuous everywhere but nowhere differentiable.
\end{remarkbox}
\color{black}

\section{Continuous Functions}

A function is continuous at $a$ when the limit and the function value agree --- when there is no jump, hole, or oscillation at $a$.  Informally, you can draw the graph without lifting your pen.  Formally, continuity is the condition that $f$ \emph{preserves} the notion of closeness: points near $a$ map to points near $f(a)$.  This is the natural condition for a function to interact well with limits, and nearly every theorem in analysis will require continuity as a hypothesis.

\begin{defbox}{Continuity at a Point}
$f : \R \to \R$ is \textbf{continuous at} $a \in \R$ if $\lim_{x \to a} f(x) = f(a)$, i.e.:
\[ \forall \eps > 0,\; \exists \del > 0,\; |x - a| < \del \Rightarrow |f(x) - f(a)| < \eps. \]
$f$ is \textbf{continuous on} $S \subseteq \R$ if it is continuous at every point of $S$.
\end{defbox}
\color{black}

The three conditions the definition packages together: $f(a)$ is defined, $\lim_{x\to a} f(x)$ exists, and they are equal.

\begin{leanbox}
\begin{lstlisting}
def ContinuousAt (f : Real -> Real) (a : Real) : Prop :=
  forall eps : Real, eps > 0 ->
    exists delta : Real, delta > 0 /\
      forall x : Real, |x - a| < delta -> |f x - f a| < eps

-- Polynomials are continuous: x^2 is continuous at every a
theorem sq_continuous_at (a : Real) : ContinuousAt (fun x => x * x) a := by
  intro eps heps
  -- |x^2 - a^2| = |x-a||x+a|. If |x-a| < 1, |x+a| <= |x-a| + 2|a| < 1 + 2|a|
  refine ⟨min 1 (eps / (1 + 2 * |a| + 1)), by positivity, fun x hx => ?_⟩
  have hd1 : |x - a| < 1 := lt_of_lt_of_le hx (min_le_left _ _)
  have hde : |x - a| < eps / (1 + 2 * |a| + 1) :=
    lt_of_lt_of_le hx (min_le_right _ _)
  have hxa : |x + a| < 1 + 2 * |a| := by
    have : |x + a| <= |x - a| + 2 * |a| := by
      calc |x + a| = |(x - a) + 2 * a| := by ring_nf
        _ <= |x - a| + |2 * a| := abs_add _ _
        _ = |x - a| + 2 * |a| := by rw [abs_mul]; norm_num
    linarith
  calc |x * x - a * a|
      = |x - a| * |x + a| := by rw [← abs_mul]; ring_nf
    _ < (eps / (1 + 2 * |a| + 1)) * (1 + 2 * |a|) := by nlinarith
    _ < eps := by
        apply div_mul_lt_iff.mpr
        · ring_nf; linarith
        · linarith
\end{lstlisting}
\end{leanbox}
\color{black}

\section{Properties of Continuous Functions}

Continuous functions are closed under the basic constructions of analysis: sums, products, quotients, and compositions.  This means that once we know a handful of basic functions are continuous --- polynomials, $\sin$, $\cos$, $e^x$ --- the entire zoo of functions built from them by algebra and composition is automatically continuous wherever it is defined.  We never need to return to the $\varepsilon$-$\delta$ definition for a composed function.

\begin{thmbox}{Algebra of Continuous Functions}
If $f$ and $g$ are continuous at $a$, then $f + g$, $fg$, $cf$ (for $c \in \R$), and $f/g$ (when $g(a) \neq 0$) are continuous at $a$.
\end{thmbox}
\color{black}

\begin{thmbox}{Composition of Continuous Functions}
If $g$ is continuous at $a$ and $f$ is continuous at $g(a)$, then $f \circ g$ is continuous at $a$.
\end{thmbox}
\color{black}

\begin{proof}
Given $\eps > 0$, choose $\del_1 > 0$ so $|y - g(a)| < \del_1 \Rightarrow |f(y) - f(g(a))| < \eps$. Choose $\del > 0$ so $|x - a| < \del \Rightarrow |g(x) - g(a)| < \del_1$. Then $|x - a| < \del \Rightarrow |f(g(x)) - f(g(a))| < \eps$.
\end{proof}

\section{The Intermediate Value Theorem}

The Intermediate Value Theorem captures the intuition that a continuous function cannot skip values: if it starts at $f(a)$ and ends at $f(b)$, it must pass through every value in between.  This is a purely topological fact --- it follows from the connectedness of $[a,b]$ --- and it is one of the deepest consequences of the completeness of $\R$.  The same theorem fails for $\Q$: the function $f(q) = q^2 - 2$ satisfies $f(1) < 0 < f(2)$, yet $f(q) = 0$ has no rational solution.

\begin{thmbox}{Intermediate Value Theorem (IVT)}
If $f : [a, b] \to \R$ is continuous and $f(a) < v < f(b)$ (or $f(b) < v < f(a)$), then there exists $c \in (a, b)$ with $f(c) = v$.
\end{thmbox}
\color{black}

\begin{proof}
Assume $f(a) < v < f(b)$. Consider $S = \{x \in [a,b] : f(x) \leq v\}$. $S$ is nonempty (contains $a$) and bounded above by $b$, so $c := \sup S$ exists. We claim $f(c) = v$.

If $f(c) < v$: by continuity, $f(x) < v$ for $x$ near $c$, so some $c' > c$ has $f(c') < v$, contradicting $c = \sup S$.

If $f(c) > v$: similarly, $f(x) > v$ near $c$, so $c - \del$ is an upper bound for $S$ for small $\del$---contradicting $c = \sup S$.

Hence $f(c) = v$.
\end{proof}

\begin{corbox}{Existence of $n$-th Roots}
For any $y > 0$ and $n \in \N$, $y^{1/n}$ exists as a real number.
\end{corbox}
\color{black}

\begin{proof}
Apply IVT to $f(x) = x^n$ on $[0, 1 + y]$: $f(0) = 0 < y \leq (1+y)^n = f(1+y)$.
\end{proof}

\section{Uniform Continuity}

Ordinary continuity says the tolerance $\delta$ for a given precision $\varepsilon$ may depend on the point $x$ we are near.  Uniform continuity says $\delta$ can be chosen to work \emph{everywhere on the domain at once}.  This is a strictly stronger condition: $f(x) = 1/x$ is continuous on $(0,1)$ but not uniformly so, because near $0$ you need a much smaller $\delta$ than near $1$.  Uniform continuity is what integration needs: to approximate $\int_a^b f$ by Riemann sums, we need the oscillation of $f$ on each subinterval to be uniformly small, not just pointwise small.

\begin{defbox}{Uniform Continuity}
$f$ is \textbf{uniformly continuous} on $S$ if the $\del$ in the continuity definition can be chosen independently of $x$:
\[ \forall \eps > 0,\; \exists \del > 0,\; \forall x, y \in S,\; |x - y| < \del \Rightarrow |f(x) - f(y)| < \eps. \]
\end{defbox}
\color{black}

\begin{thmbox}{Continuous Functions on Closed Intervals Are Uniformly Continuous}
If $f$ is continuous on $[a, b]$, it is uniformly continuous on $[a, b]$.
\end{thmbox}
\color{black}

\begin{proof}
Suppose not. Then there exists $\eps > 0$ and sequences $x_n, y_n$ in $[a,b]$ with $|x_n - y_n| < 1/n$ but $|f(x_n) - f(y_n)| \geq \eps$. By Bolzano--Weierstrass, $(x_n)$ has a convergent subsequence $x_{n_k} \to c \in [a,b]$. Then $y_{n_k} \to c$ as well. By continuity at $c$, $f(x_{n_k}) \to f(c)$ and $f(y_{n_k}) \to f(c)$, so $|f(x_{n_k}) - f(y_{n_k})| \to 0$---contradicting $\geq \eps$.
\end{proof}

\section{Pathological Examples}

Not every function is continuous. Two canonical examples:

\begin{example}[Dirichlet's Function]
$f(x) = \begin{cases} 1 & x \in \Q \\ 0 & x \notin \Q \end{cases}$ is continuous nowhere. For any $a \in \R$ and $\del > 0$, the interval $(a-\del, a+\del)$ contains both rationals and irrationals, so $f$ oscillates between $0$ and $1$.
\end{example}

\begin{example}[The Ruler Function]
$f(x) = \begin{cases} 1/q & x = p/q \text{ in lowest terms} \\ 0 & x \notin \Q \end{cases}$ is continuous at every irrational, discontinuous at every rational.
\end{example}

\begin{exbox}
\begin{enumerate}
  \item Prove that $f(x) = \sqrt{x}$ is uniformly continuous on $[0, \infty)$.
  \item Prove the Extreme Value Theorem: if $f$ is continuous on $[a,b]$, then $f$ attains a maximum and minimum value.
  \item Show that $f(x) = 1/x$ is continuous but not uniformly continuous on $(0, 1)$.
  \item Use the IVT to prove that every odd-degree polynomial with real coefficients has at least one real root.
  \item Formalize continuity of $f(x) = x^3$ at an arbitrary point $a$ in Lean~4.
  \item Prove the fixed-point theorem: if $f : [0,1] \to [0,1]$ is continuous, then there exists $c$ with $f(c) = c$.
\end{enumerate}
\end{exbox}
\color{black}

%%-----------------------------------------------------------------------------
\chapter{Topology of \texorpdfstring{$\R$}{R}}
%%-----------------------------------------------------------------------------
\color{black}

The theorems of the previous chapter (IVT, EVT, uniform continuity) all relied on the completeness of $\R$ and properties of $[a,b]$ as a domain. This chapter makes those structural properties precise via the language of point-set topology.

\section{Open and Closed Sets}

The theorems of the previous chapter relied on specific properties of $[a,b]$: it is bounded, it contains its endpoints, and it is connected.  Point-set topology gives us precise language for these properties.  An open set is one where every point has breathing room --- a small neighbourhood entirely inside the set.  A closed set contains all the limit points it could ever accumulate.  These two notions are dual, and together they provide the vocabulary for stating the deepest theorems of analysis in their most general form.

\begin{defbox}{Open and Closed Subsets of $\R$}
A set $U \subseteq \R$ is \textbf{open} if every point has a neighborhood contained in $U$:
\[ \forall x \in U,\; \exists \eps > 0,\; (x - \eps, x + \eps) \subseteq U. \]
A set $F \subseteq \R$ is \textbf{closed} if its complement $\R \setminus F$ is open, equivalently if $F$ contains all its limit points.
\end{defbox}
\color{black}

\begin{propbox}{Basic Topological Properties}
\begin{enumerate}[noitemsep]
  \item $\emptyset$ and $\R$ are both open and closed.
  \item Arbitrary unions of open sets are open.
  \item Finite intersections of open sets are open.
  \item Arbitrary intersections of closed sets are closed.
  \item Finite unions of closed sets are closed.
\end{enumerate}
\end{propbox}
\color{black}

\section{Compactness}

Compactness is the topological abstraction of ``closed and bounded,'' and it is the property that makes analysis work on finite-length intervals.  The key feature of compact sets is that every open cover has a \emph{finite} subcover: you cannot cover $[a,b]$ with infinitely many open sets without finitely many of them already doing the job.  This finiteness is what lets us extract uniform bounds and guarantees from pointwise ones --- it is the engine behind the Extreme Value Theorem, uniform continuity, and the completeness of the Riemann integral.

\begin{defbox}{Compactness}
A set $K \subseteq \R$ is \textbf{compact} if every open cover of $K$ has a finite subcover:
if $K \subseteq \bigcup_{\alpha} U_\alpha$ (each $U_\alpha$ open), then $K \subseteq U_{\alpha_1} \cup \cdots \cup U_{\alpha_n}$ for finitely many $\alpha_i$.

Equivalently (in $\R$): $K$ is compact iff $K$ is \textbf{closed and bounded} (Heine--Borel theorem).
\end{defbox}
\color{black}

\begin{thmbox}{Heine--Borel Theorem}
A subset $K \subseteq \R$ is compact if and only if it is closed and bounded.
\end{thmbox}
\color{black}

\begin{proof}
$(\Rightarrow)$: Compact sets in a metric space are closed and bounded. Boundedness: the cover $\{(-n,n)\}_{n\geq 1}$ has a finite subcover. Closure: if $x$ is a limit point not in $K$, the sets $\R \setminus \{y : |y-x| \leq 1/n\}$ cover $K$ but no finite subcollection does.

$(\Leftarrow)$: Closed and bounded in $\R$ implies compact. The key lemma (proved by the bisection method, using completeness) is that $[a,b]$ is compact; a closed subset of a compact set is compact.
\end{proof}

\begin{thmbox}{Extreme Value Theorem}
If $f : K \to \R$ is continuous and $K$ is compact, then $f$ attains its maximum and minimum on $K$.
\end{thmbox}
\color{black}

\begin{proof}
The image $f(K)$ is compact (continuous images of compact sets are compact). A compact subset of $\R$ is closed and bounded, hence has a maximum and minimum, both achieved.
\end{proof}

\section{Connectedness}

A connected set cannot be split into two separated open pieces.  Intuitively, it is ``in one piece.''  In $\R$, the connected sets are exactly the intervals --- a fact that is not obvious from the definition but is a theorem.  Connectedness is what powers the Intermediate Value Theorem: the image of a connected set under a continuous function is connected, so it must itself be an interval, and therefore contains all values between its endpoints.

\begin{defbox}{Connectedness}
A set $S \subseteq \R$ is \textbf{connected} if it cannot be written as a disjoint union $S = A \cup B$ of two nonempty open sets. The connected subsets of $\R$ are exactly the intervals (including rays and $\R$ itself).
\end{defbox}
\color{black}

\begin{thmbox}{Continuous Images of Connected Sets Are Connected}
If $f$ is continuous and $S$ is connected, then $f(S)$ is connected.
\end{thmbox}
\color{black}

\begin{proof}
Suppose $f(S) = A \cup B$ (disjoint open sets). Then $S = f^{-1}(A) \cup f^{-1}(B)$, both open by continuity of $f$. Since $S$ is connected, one must be empty.
\end{proof}

\begin{corbox}{Intermediate Value Theorem (Topological Proof)}
If $f : [a,b] \to \R$ is continuous and $v$ is between $f(a)$ and $f(b)$, then $v \in f([a,b])$.
\end{corbox}
\color{black}

\begin{proof}
$[a,b]$ is connected, so $f([a,b])$ is connected, hence an interval. Since $f(a)$ and $f(b)$ are in this interval and $v$ is between them, $v$ is in the interval.
\end{proof}

\begin{exbox}
\begin{enumerate}
  \item Prove that $(a,b)$ is open and $[a,b]$ is closed using the definitions.
  \item Prove that an infinite intersection of open sets need not be open (example: $\bigcap_{n=1}^\infty (-1/n, 1/n) = \{0\}$).
  \item Prove that compact subsets of $\R$ are closed and bounded.
  \item Prove: a continuous image of a compact set is compact.
  \item Show that $[0,1]$ and $(0,1)$ are not homeomorphic (hint: consider what happens when you remove a point).
  \item Prove that every continuous $f : [a,b] \to [a,b]$ has a fixed point (Brouwer in 1D).
\end{enumerate}
\end{exbox}
\color{black}

%%=============================================================================
\part{Differential Calculus}
%%=============================================================================
\color{black}

%%-----------------------------------------------------------------------------
\chapter{The Derivative}
%%-----------------------------------------------------------------------------
\color{black}

\begin{remarkbox}
\textbf{Historical Note: Newton, Leibniz, and the Priority Dispute}

The invention of calculus is one of the most famous priority disputes in the history of science. \textbf{Isaac Newton} developed his method of ``fluxions'' in 1665--66, during the plague years when Cambridge University was closed and he retreated to his family farm in Woolsthorpe. He did not publish his results until decades later. \textbf{Gottfried Wilhelm Leibniz} independently developed his calculus in 1675--76, invented the notation $dy/dx$ and $\int$, and published his results in 1684 and 1686 --- decades before Newton's publication.

Newton's priority over Leibniz in \emph{discovery} is not disputed: his notebooks are dated. But Leibniz's priority in \emph{publication} and the superiority of his \emph{notation} are equally clear. The dispute, which consumed both men's later years and poisoned English and Continental mathematics for a century, was fuelled by nationalist pride and the accusation (almost certainly false) that Leibniz had seen Newton's unpublished manuscripts during a visit to London in 1676. The Royal Society investigated in 1712 and, in a verdict widely regarded as dishonest, found for Newton --- largely because Newton, as President of the Royal Society, effectively controlled the investigation.

The mathematical cost was severe: British mathematicians clung to Newton's cumbersome ``dot'' notation ($\dot{x}$, $\ddot{x}$) out of loyalty, while the Continent used Leibniz's far more flexible $dy/dx$ notation. This notational disadvantage held British mathematics back for over a century. The Cambridge Analytical Society, founded in 1812 by Babbage, Herschel, and Peacock, explicitly aimed to replace the ``d-ism of Leibniz'' for the ``dot-age of Newton'' --- their pun, not ours.

The notation we use in this chapter is Lagrange's: $f'(x)$ for the derivative. Lagrange introduced this in his 1797 \emph{Th\'{e}orie des fonctions analytiques}, where he also introduced the term ``derivative'' (\emph{fonction d\'{e}riv\'{e}e}). Lagrange's approach was to define the derivative via Taylor series rather than limits, which is logically backwards (you need the derivative to define the Taylor series) but has the virtue of being computationally clean.
\end{remarkbox}
\color{black}

The derivative formalizes the notion of instantaneous rate of change. Geometrically, it is the slope of the tangent line to a curve at a point; analytically, it is a limit of difference quotients.

\section{The Difference Quotient and the Derivative}

The derivative formalises instantaneous rate of change.  The idea is to approximate the rate of change on a small interval $[x, x+h]$ by the slope $(f(x+h)-f(x))/h$ of the secant line, then take the limit as $h \to 0$.  If this limit exists, the function has a well-defined instantaneous slope at $x$ --- its tangent line slope.

The ``local linearity'' formulation $f(x+h) \approx f(x) + f'(x)h$ is in some ways more fundamental: it says differentiability is equivalent to the existence of a good linear approximation.  This is the form that generalises cleanly to $\R^n$, where the derivative becomes a linear map rather than a number.

\color{black}\begin{center}
\begin{tikzpicture}[xscale=1.6, yscale=1.3]
  \draw[->] (-0.2,0) -- (3.8,0) node[right, font=\small] {\textcolor{black}{$x$}};
  \draw[->] (0,-0.2) -- (0,3.0) node[above, font=\small] {\textcolor{black}{$y$}};

  % f(x) = 0.4x^2 + 0.3
  \draw[very thick, color=thmblue, domain=0.15:3.6, samples=60]
    plot (\x, {0.4*\x*\x + 0.3});

  \def\xp{1.5}   % point x
  \def\yp{0.4*\xp*\xp + 0.3}  % = 0.9 + 0.3 = 1.2
  \def\xph{2.5}  % x + h
  \def\yph{0.4*\xph*\xph + 0.3} % = 2.5 + 0.3 = 2.8

  % Secant line through (xp, yp) and (xph, yph)
  % slope = (yph - yp)/(xph - xp) = (2.8 - 1.2)/1 = 1.6
  % y - 1.2 = 1.6(x - 1.5)  => y = 1.6x - 2.4 + 1.2 = 1.6x - 1.2
  \draw[color=exborder, thick, domain=0.8:3.2]
    plot (\x, {1.6*\x - 1.2});
  \node[font=\scriptsize, rotate=50] at (2.8, 3.1)
    {\textcolor{exborder}{secant slope $= \frac{f(x+h)-f(x)}{h}$}};

  % Tangent line at xp
  \draw[color=defgreen, thick, domain=0.6:2.8]
    plot (\x, {1.2*\x - 0.6});
  \node[font=\scriptsize, rotate=40] at (0.7, 0.35)
    {\textcolor{defgreen}{tangent slope $= f'(x)$}};

  % Points
  \fill[thmblue] (\xp, {0.4*\xp*\xp + 0.3}) circle (2.5pt)
    node[above left, font=\scriptsize] {\textcolor{black}{$(x,f(x))$}};
  \fill[exborder] (\xph, {0.4*\xph*\xph + 0.3}) circle (2.5pt)
    node[above left, font=\scriptsize] {\textcolor{black}{$(x{+}h,\,f(x{+}h))$}};

  % h and rise labels
  \draw[|<->|, color=gray, thin] (\xp, -0.12) -- (\xph, -0.12)
    node[midway, below, font=\scriptsize] {\textcolor{gray}{$h$}};
  \draw[|<->|, color=gray, thin] (\xph+0.08, {0.4*\xp*\xp+0.3}) -- (\xph+0.08, {0.4*\xph*\xph+0.3})
    node[midway, right=3pt, font=\scriptsize] {\textcolor{gray}{$f(x{+}h)-f(x)$}};

  % Dashed drop lines
  \draw[dashed, thin, gray] (\xp,  0) -- (\xp,  {0.4*\xp*\xp+0.3});
  \draw[dashed, thin, gray] (\xph, 0) -- (\xph, {0.4*\xph*\xph+0.3});
  \draw[dashed, thin, gray] (\xph, {0.4*\xp*\xp+0.3}) -- (\xph, {0.4*\xph*\xph+0.3});
  \draw[dashed, thin, gray] (\xp,  {0.4*\xp*\xp+0.3}) -- (\xph, {0.4*\xp*\xp+0.3});

  % axis ticks
  \draw (\xp,  0.03) -- (\xp,  -0.03) node[below, font=\scriptsize] {\textcolor{black}{$x$}};
  \draw (\xph, 0.03) -- (\xph, -0.03) node[below, font=\scriptsize] {\textcolor{black}{$x{+}h$}};

  % Arrow showing h -> 0
  \node[font=\scriptsize\itshape, align=center] at (3.2, 0.5)
    {\textcolor{thmblue!80}{as $h\to 0$:}\\
     \textcolor{thmblue!80}{secant $\to$ tangent}};
  \draw[->, thmblue!50, thin] (2.9, 0.68) -- (2.3, 1.1);
\end{tikzpicture}
\end{center}

\begin{defbox}{The Derivative at a Point}
Let $f : (a,b) \to \R$ and $x \in (a,b)$. The \textbf{derivative} of $f$ at $x$ is:
\[ f'(x) := \lim_{h \to 0} \frac{f(x+h) - f(x)}{h}, \]
provided this limit exists. If it does, $f$ is \textbf{differentiable at} $x$.

Equivalently, $f'(x) = L$ if $\lim_{h\to 0} \frac{f(x+h) - f(x) - Lh}{h} = 0$, i.e., $f$ is ``locally linear'' with slope $L$.
\end{defbox}
\color{black}

This ``local linearity'' formulation will generalize cleanly to $\R^n$. We record it as a separate characterization:

\begin{propbox}{Local Linearity Characterization}
$f'(x) = L$ if and only if $f(x + h) = f(x) + Lh + o(h)$ as $h \to 0$, where $o(h)$ means the error divided by $h$ goes to zero.
\end{propbox}
\color{black}

\begin{thmbox}{Differentiability Implies Continuity}
If $f$ is differentiable at $x$, then $f$ is continuous at $x$.
\end{thmbox}
\color{black}

\begin{proof}
$f(x+h) - f(x) = \frac{f(x+h)-f(x)}{h} \cdot h \to f'(x) \cdot 0 = 0$ as $h \to 0$.
\end{proof}

\begin{remarkbox}
The converse fails: $f(x) = |x|$ is continuous at $0$ but not differentiable there (left derivative $= -1$, right derivative $= 1$). Even more dramatically, Weierstrass constructed a function continuous everywhere but differentiable nowhere.
\end{remarkbox}
\color{black}

\begin{leanbox}
\begin{lstlisting}
def HasDerivAt (f : Real -> Real) (x L : Real) : Prop :=
  forall eps : Real, eps > 0 ->
    exists delta : Real, delta > 0 /\
      forall h : Real, |h| < delta -> h ≠ 0 ->
        |(f (x + h) - f x) / h - L| < eps

-- The derivative of a constant is zero
theorem const_deriv (c x : Real) : HasDerivAt (fun _ => c) x 0 := by
  intro eps heps
  exact ⟨1, one_pos, fun h _ _ => by
    simp; rw [div_self (ne_of_gt (by linarith))]; simp; exact heps⟩

-- The derivative of x is 1
theorem id_deriv (x : Real) : HasDerivAt id x 1 := by
  intro eps heps
  exact ⟨1, one_pos, fun h _ _ => by
    simp [id]; rw [add_sub_cancel_left, div_self (by assumption)]; simp; linarith⟩

-- Differentiability implies continuity
theorem differentiable_implies_continuous (f : Real -> Real) (x L : Real)
    (hf : HasDerivAt f x L) : ContinuousAt f x := by
  intro eps heps
  obtain ⟨delta, hd, hfunc⟩ := hf (eps / (|L| + 1))
    (div_pos heps (by linarith [abs_nonneg L]))
  refine ⟨min delta 1, by positivity, fun y hy => ?_⟩
  sorry -- full proof uses: f(x+h) - f(x) = h * ((f(x+h)-f(x))/h) -> 0
\end{lstlisting}
\end{leanbox}
\color{black}

\section{Derivatives of Elementary Functions}

We compute the derivatives of the basic functions, which will become the foundation of the differentiation engine in Part~VII.

\begin{thmbox}{Derivatives of Power Functions}
For $n \in \N$, $(x^n)' = nx^{n-1}$, with the convention $x^0 = 1$.
\end{thmbox}
\color{black}

\begin{proof}
By the binomial theorem:
\[ \frac{(x+h)^n - x^n}{h} = \frac{\sum_{k=0}^n \binom{n}{k} x^k h^{n-k} - x^n}{h} = nx^{n-1} + \binom{n}{2}x^{n-2}h + \cdots \to nx^{n-1}. \]
\end{proof}

\begin{thmbox}{Derivatives of Trigonometric Functions}
$(\sin x)' = \cos x$ and $(\cos x)' = -\sin x$.
\end{thmbox}
\color{black}

\begin{proof}
Using the angle-addition formula and the limits $\lim_{h\to 0} \sin(h)/h = 1$ and $\lim_{h\to 0} (1 - \cos h)/h = 0$:
\[ \frac{\sin(x+h) - \sin x}{h} = \cos x \cdot \frac{\sin h}{h} - \sin x \cdot \frac{1 - \cos h}{h} \to \cos x. \]
\end{proof}

\begin{thmbox}{Derivative of $e^x$}
$(e^x)' = e^x$.
\end{thmbox}
\color{black}

\begin{proof}
$(e^{x+h} - e^x)/h = e^x(e^h - 1)/h \to e^x \cdot 1$, where $\lim_{h\to 0}(e^h-1)/h = 1$ follows from the power series $e^h = 1 + h + h^2/2! + \cdots$.
\end{proof}

\begin{exbox}
\begin{enumerate}
  \item Prove $(\ln x)' = 1/x$ for $x > 0$.
  \item Show directly from the definition that $(|x|)' = \text{sgn}(x)$ for $x \neq 0$ and that $|x|$ is not differentiable at $0$.
  \item Prove $(x^n)' = nx^{n-1}$ for negative integers $n$ (where $x \neq 0$).
  \item Prove $(\tan x)' = \sec^2 x$ using $\tan x = \sin x / \cos x$ and the quotient rule from the next chapter.
  \item Formalize \li{HasDerivAt (fun x => x * x * x) a (3 * a * a)} in Lean~4.
\end{enumerate}
\end{exbox}
\color{black}

%%-----------------------------------------------------------------------------
\chapter{Differentiation Rules}
%%-----------------------------------------------------------------------------
\color{black}

\section{The Four Fundamental Rules}

Every differentiation rule is a theorem --- a consequence of the limit definition.  We do not assume products differentiate by the product rule; we prove it.  These four rules (sum, product, quotient, chain) form a complete toolkit: any expression built from elementary functions by these four operations can be differentiated algorithmically by pattern-matching against the rules.  This is precisely what the differentiation engine in Part~VII will do.

\begin{thmbox}{Sum Rule}
If $f$ and $g$ are differentiable at $x$, then $(f+g)'(x) = f'(x) + g'(x)$.
\end{thmbox}
\color{black}

\begin{proof}
$\frac{(f(x+h)+g(x+h)) - (f(x)+g(x))}{h} = \frac{f(x+h)-f(x)}{h} + \frac{g(x+h)-g(x)}{h} \to f'(x) + g'(x)$.
\end{proof}

\begin{thmbox}{Product Rule}
$(fg)'(x) = f'(x)g(x) + f(x)g'(x)$.
\end{thmbox}
\color{black}

\begin{proof}
\[ \frac{f(x+h)g(x+h) - f(x)g(x)}{h} = \frac{f(x+h)-f(x)}{h} \cdot g(x+h) + f(x) \cdot \frac{g(x+h)-g(x)}{h}. \]
As $h \to 0$: first factor $\to f'(x)$, $g(x+h) \to g(x)$ (by continuity), second $\to f(x)g'(x)$.
\end{proof}

\begin{thmbox}{Quotient Rule}
If $g(x) \neq 0$, then $(f/g)'(x) = \frac{f'(x)g(x) - f(x)g'(x)}{g(x)^2}$.
\end{thmbox}
\color{black}

\begin{proof}
Write $f/g = f \cdot (1/g)$ and apply the product rule, after computing $(1/g)'(x) = -g'(x)/g(x)^2$.
\end{proof}

\begin{thmbox}{Chain Rule}
If $g$ is differentiable at $x$ and $f$ is differentiable at $g(x)$, then:
\[ (f \circ g)'(x) = f'(g(x)) \cdot g'(x). \]
\end{thmbox}
\color{black}

\begin{proof}
The naive proof (multiplying and dividing by $g(x+h) - g(x)$) fails when $g(x+h) = g(x)$. The correct proof uses the local linearity characterization: write $f(g(x)+k) = f(g(x)) + f'(g(x))k + \eps(k)k$ where $\eps(k) \to 0$. Setting $k = g(x+h) - g(x) = g'(x)h + \eta(h)h$:
\[ f(g(x+h)) - f(g(x)) = f'(g(x))(g'(x)h + \eta h) + \eps(k)(g'(x)h + \eta h). \]
Dividing by $h$ and taking $h \to 0$ gives $f'(g(x))g'(x)$.
\end{proof}

\begin{leanbox}
\begin{lstlisting}
-- Product rule
theorem product_rule (f g : Real -> Real) (x Lf Lg : Real)
    (hf : HasDerivAt f x Lf) (hg : HasDerivAt g x Lg) :
    HasDerivAt (fun x => f x * g x) x (Lf * g x + f x * Lg) := by
  intro eps heps
  -- Need: |(f(x+h)g(x+h) - f(x)g(x))/h - (Lf * g(x) + f(x) * Lg)| < eps
  -- Use: f(x+h)g(x+h) - f(x)g(x) = (f(x+h)-f(x))g(x+h) + f(x)(g(x+h)-g(x))
  sorry  -- Full proof requires bounding g near x

-- Chain rule
theorem chain_rule (f g : Real -> Real) (x Lf Lg : Real)
    (hg : HasDerivAt g x Lg) (hf : HasDerivAt f (g x) Lf) :
    HasDerivAt (f ∘ g) x (Lf * Lg) := by
  sorry  -- Requires the local linearity formulation
\end{lstlisting}
\end{leanbox}
\color{black}

\section{Implicit and Logarithmic Differentiation}

Sometimes a function is defined implicitly by a relation $F(x, y) = 0$ rather than by an explicit formula $y = f(x)$.

\begin{example}[Implicit Differentiation]
The unit circle $x^2 + y^2 = 1$ defines $y$ implicitly as a function of $x$ near any point with $y \neq 0$. Differentiating both sides with respect to $x$: $2x + 2y\,dy/dx = 0$, so $dy/dx = -x/y$.
\end{example}

\begin{example}[Logarithmic Differentiation]
To differentiate $f(x) = x^x$ for $x > 0$: let $y = x^x$, so $\ln y = x \ln x$. Differentiating: $y'/y = \ln x + 1$, so $y' = x^x(\ln x + 1)$.
\end{example}

\begin{thmbox}{Inverse Function Theorem (1D)}
If $f$ is differentiable and strictly monotone on $(a,b)$ with $f'(c) \neq 0$, then $f^{-1}$ is differentiable at $f(c)$ with:
\[ (f^{-1})'(f(c)) = \frac{1}{f'(c)}. \]
\end{thmbox}
\color{black}

\begin{proof}
Let $y = f(c)$, $k = f(c+h) - f(c)$. As $h \to 0$, $k \to 0$ (by continuity). Then:
\[ (f^{-1})'(y) = \lim_{k \to 0} \frac{f^{-1}(y+k) - f^{-1}(y)}{k} = \lim_{h \to 0} \frac{h}{f(c+h) - f(c)} = \frac{1}{f'(c)}. \]
\end{proof}

\begin{exbox}
\begin{enumerate}
  \item Use the inverse function theorem to prove $(\arcsin x)' = 1/\sqrt{1-x^2}$.
  \item Prove $(a^x)' = a^x \ln a$ for $a > 0$ using logarithmic differentiation.
  \item Prove the general power rule $(x^r)' = rx^{r-1}$ for real $r$ (write $x^r = e^{r \ln x}$).
  \item Differentiate $y = (\sin x)^{\cos x}$ using logarithmic differentiation.
  \item Formalize the sum rule in Lean~4: if \li{HasDerivAt f x Lf} and \li{HasDerivAt g x Lg}, then \li{HasDerivAt (fun x => f x + g x) x (Lf + Lg)}.
\end{enumerate}
\end{exbox}
\color{black}

%%-----------------------------------------------------------------------------
\chapter{Mean Value Theorems}
%%-----------------------------------------------------------------------------
\color{black}

\section{Rolle's Theorem and the MVT}

The Mean Value Theorem is one of the most useful theorems in calculus: it says that the average rate of change over an interval $[a,b]$ is achieved as an instantaneous rate at some interior point.  Geometrically, the secant line from $(a, f(a))$ to $(b, f(b))$ has the same slope as some tangent line.  Rolle's theorem is the special case where the endpoints have equal values, and the MVT follows from it by a linear correction.

\begin{thmbox}{Rolle's Theorem}
If $f : [a,b] \to \R$ is continuous, differentiable on $(a,b)$, and $f(a) = f(b)$, then there exists $c \in (a,b)$ with $f'(c) = 0$.
\end{thmbox}
\color{black}

\begin{proof}
By EVT, $f$ attains its maximum $M$ and minimum $m$ on $[a,b]$. If $M = m$, $f$ is constant and $f' \equiv 0$. Otherwise, at least one extremum is attained at an interior point $c \in (a,b)$. At an interior maximum, $f'(c) \geq 0$ (from the left) and $f'(c) \leq 0$ (from the right), so $f'(c) = 0$.
\end{proof}

\begin{thmbox}{Mean Value Theorem (MVT)}
If $f : [a,b] \to \R$ is continuous and differentiable on $(a,b)$, then there exists $c \in (a,b)$ with:
\[ f'(c) = \frac{f(b) - f(a)}{b - a}. \]
\end{thmbox}
\color{black}

\begin{proof}
Define $g(x) = f(x) - \frac{f(b)-f(a)}{b-a}(x-a)$. Then $g(a) = f(a) = g(b)$. By Rolle's theorem, $g'(c) = 0$ for some $c$, i.e., $f'(c) = \frac{f(b)-f(a)}{b-a}$.
\end{proof}

\begin{corbox}{Monotonicity via the Derivative}
If $f'(x) > 0$ on $(a,b)$, then $f$ is strictly increasing. If $f'(x) < 0$, strictly decreasing.
If $f'(x) = 0$ everywhere on $(a,b)$, then $f$ is constant.
\end{corbox}
\color{black}

\begin{proof}
For $a < x < y < b$: by MVT, $f(y) - f(x) = f'(c)(y-x)$ for some $c \in (x,y)$. The sign of $f'(c)$ determines whether $f(y) > f(x)$ or $f(y) < f(x)$.
\end{proof}

\section{L'H\^opital's Rule}

L'H\^opital's rule resolves indeterminate forms like $0/0$ or $\infty/\infty$ by replacing the ratio of functions with the ratio of their derivatives.  The intuition: near the point where both numerator and denominator vanish, the dominant behaviour is determined by first-order Taylor terms, whose ratio is the ratio of derivatives.  The rule is powerful but requires care: it applies only to genuine indeterminate forms, and $\lim f'/g'$ must actually exist.

\begin{thmbox}{L'H\^opital's Rule}
Suppose $\lim_{x \to a} f(x) = \lim_{x \to a} g(x) = 0$ (or both $\pm\infty$), and $g'(x) \neq 0$ near $a$. If $\lim_{x\to a} f'(x)/g'(x) = L$, then $\lim_{x \to a} f(x)/g(x) = L$.
\end{thmbox}
\color{black}

\begin{proof}
In the $0/0$ case: define $f(a) = g(a) = 0$. For $x$ near $a$, by Cauchy's MVT there exists $c$ between $a$ and $x$ with $\frac{f(x)}{g(x)} = \frac{f(x)-f(a)}{g(x)-g(a)} = \frac{f'(c)}{g'(c)}$. As $x \to a$, $c \to a$, so the ratio approaches $L$.
\end{proof}

\begin{exbox}
\begin{enumerate}
  \item Use L'H\^opital to evaluate $\lim_{x \to 0} (\sin x)/x$ and $\lim_{x \to 0} (e^x - 1 - x)/x^2$.
  \item Prove: if $f'(x) \geq 0$ on $(a,b)$ and $f'$ is not identically zero, then $f$ is strictly increasing.
  \item Use the MVT to prove the AM-GM inequality: for $a, b > 0$, $\sqrt{ab} \leq (a+b)/2$.
  \item Prove Cauchy's MVT: if $f, g$ are continuous on $[a,b]$, differentiable on $(a,b)$, and $g'(x) \neq 0$ on $(a,b)$, then there exists $c$ with $f'(c)/g'(c) = (f(b)-f(a))/(g(b)-g(a))$.
\end{enumerate}
\end{exbox}
\color{black}

%%-----------------------------------------------------------------------------
\chapter{Higher Derivatives and Taylor's Theorem}
%%-----------------------------------------------------------------------------
\color{black}

\begin{remarkbox}
\textbf{Historical Note: Taylor, Maclaurin, and the Power Series Debate}

Brook Taylor published his series in 1715, claiming (as was common at the time) more generality than he actually proved. The series had been known earlier: James Gregory had used it in 1671, and Leibniz had the expansion for $\arctan$ in 1673. What Taylor contributed was a systematic account treating the series as a general tool rather than a collection of special cases.

Colin Maclaurin, whose name is attached to the special case $a = 0$, never claimed to have discovered it --- he explicitly credited Taylor. The Maclaurin series appears in his 1742 \emph{Treatise of Fluxions}, written partly to defend Newton's calculus against Berkeley's philosophical attacks, and the special case became associated with his name through the accident of emphasis rather than discovery.

The question of \emph{when} a smooth function equals its Taylor series --- the issue of \emph{analyticity} versus mere smoothness --- was not clearly understood until the nineteenth century. Cauchy constructed the now-standard counterexample: $f(x) = e^{-1/x^2}$ for $x \neq 0$ and $f(0) = 0$ is infinitely differentiable everywhere, yet its Taylor series at $0$ is identically zero, so the series converges but not to $f$. This function is smooth ($C^\infty$) but not analytic at $0$.

The distinction between smooth and analytic functions is fundamental in complex analysis: every holomorphic (complex-differentiable) function is automatically analytic, which is one of the remarkable rigidity properties of complex analysis. In real analysis, smoothness is far weaker than analyticity, and the gap between them is inhabited by important functions like bump functions and cutoff functions used in differential geometry and partial differential equations.

The Lagrange remainder formula we prove --- bounding $|R_n(x)|$ by $M|x-a|^{n+1}/(n+1)!$ --- was Lagrange's main contribution to Taylor series theory. It turned Taylor series from a formal tool into a theorem with an error bound, making numerical approximations rigorous.
\end{remarkbox}
\color{black}

\section{Higher-Order Derivatives}

If the derivative $f'$ is itself differentiable, we get the second derivative $f'' = (f')'$, measuring how the rate of change is itself changing.  For a position function, $f'$ is velocity and $f''$ is acceleration.  Iterating gives higher-order derivatives.  The smoothness class $C^k$ records how many times a function can be differentiated and still have a continuous result.  $C^\infty$ (infinitely differentiable) is the natural domain for Taylor series.

\begin{defbox}{Higher-Order Derivatives}
The $n$-th derivative $f^{(n)}$ is defined inductively: $f^{(0)} = f$ and $f^{(n)} = (f^{(n-1)})'$.

$f$ is of class $C^k$ on an interval if $f^{(k)}$ exists and is continuous. $f$ is \textbf{smooth} ($C^\infty$) if all derivatives exist.
\end{defbox}
\color{black}

\section{Taylor's Theorem}

Taylor's theorem gives the best polynomial approximation to a smooth function near a point, together with a precise bound on the error.  It explains why tables of sines and logarithms work: $\sin x \approx x - x^3/6$ for small $x$ because the error is bounded by $|x|^5/120$.  The Lagrange remainder turns an approximation into a theorem with a quantitative error bound, and it will be our main tool for proving that power series converge to the functions they represent.

\begin{thmbox}{Taylor's Theorem with Lagrange Remainder}
If $f$ is $n+1$ times differentiable on $(a-r, a+r)$, then for each $x$ in this interval:
\[ f(x) = \sum_{k=0}^{n} \frac{f^{(k)}(a)}{k!}(x-a)^k + R_n(x), \]
where the \textbf{Lagrange remainder} satisfies $R_n(x) = \frac{f^{(n+1)}(c)}{(n+1)!}(x-a)^{n+1}$ for some $c$ between $a$ and $x$.
\end{thmbox}
\color{black}

\begin{proof}
Define $g(t) = f(x) - \sum_{k=0}^n \frac{f^{(k)}(t)}{k!}(x-t)^k$. Then $g(x) = 0$ and $g'(t) = -\frac{f^{(n+1)}(t)}{n!}(x-t)^n$. Apply the MVT to $g$ on $[a, x]$ using the auxiliary function $h(t) = (x-t)^{n+1}$ to obtain Cauchy's MVT form of the remainder.
\end{proof}

\begin{example}
The Taylor polynomial for $\sin x$ around $0$ is:
\[ \sin x = x - \frac{x^3}{3!} + \frac{x^5}{5!} - \cdots + (-1)^n \frac{x^{2n+1}}{(2n+1)!} + R_{2n+1}(x). \]
The remainder satisfies $|R_{2n+1}(x)| \leq |x|^{2n+3}/(2n+3)!$, which $\to 0$ for any fixed $x$ as $n \to \infty$.
\end{example}

\begin{exbox}
\begin{enumerate}
  \item Compute the Taylor polynomial of degree 4 for $\cos x$ around $0$.
  \item Use Taylor's theorem to prove that $e^x \geq 1 + x$ for all $x \in \R$.
  \item Prove the second derivative test: if $f'(a) = 0$ and $f''(a) > 0$, then $a$ is a local minimum.
  \item Show that the function $f(x) = e^{-1/x^2}$ (for $x \neq 0$) and $f(0) = 0$ is $C^\infty$ but its Taylor series at $0$ is identically zero, failing to converge to $f$. This shows smooth $\neq$ analytic.
  \item Prove the general binomial theorem: $(1+x)^\alpha = \sum_{k=0}^\infty \binom{\alpha}{k} x^k$ for $|x| < 1$.
\end{enumerate}
\end{exbox}
\color{black}

%%-----------------------------------------------------------------------------
\chapter{Applications of Differentiation}
%%-----------------------------------------------------------------------------
\color{black}

\section{Optimization}

One of the most immediate applications of the derivative is finding maxima and minima.  At a local extremum in the interior of the domain, the tangent line must be horizontal --- so the derivative must be zero.  This reduces the infinite search for a maximum to a finite computation: find all critical points, classify them, and compare values.  The second derivative test provides the classification by measuring whether the function curves upward (minimum) or downward (maximum) at the critical point.

\begin{defbox}{Critical Points}
A point $c$ in the domain of $f$ is a \textbf{critical point} if $f'(c) = 0$ or $f'(c)$ does not exist.
\end{defbox}
\color{black}

\begin{thmbox}{First Derivative Test}
Suppose $f'(c) = 0$.
\begin{itemize}[noitemsep]
  \item If $f'$ changes from positive to negative at $c$, then $c$ is a local maximum.
  \item If $f'$ changes from negative to positive at $c$, then $c$ is a local minimum.
  \item If $f'$ does not change sign, $c$ is neither.
\end{itemize}
\end{thmbox}
\color{black}

\begin{thmbox}{Second Derivative Test}
If $f'(c) = 0$ and $f''(c)$ exists:
\begin{itemize}[noitemsep]
  \item $f''(c) > 0 \Rightarrow$ local minimum at $c$.
  \item $f''(c) < 0 \Rightarrow$ local maximum at $c$.
  \item $f''(c) = 0 \Rightarrow$ inconclusive.
\end{itemize}
\end{thmbox}
\color{black}

\section{Newton's Method}

Newton's method finds roots of $f(x) = 0$ using the iteration:
\[ x_{n+1} = x_n - \frac{f(x_n)}{f'(x_n)}. \]

\begin{thmbox}{Convergence of Newton's Method}
If $f$ is twice continuously differentiable near a simple root $r$ (with $f(r) = 0$, $f'(r) \neq 0$) and $x_0$ is sufficiently close to $r$, then Newton's method converges quadratically:
\[ |x_{n+1} - r| \leq C|x_n - r|^2 \]
for some constant $C$.
\end{thmbox}
\color{black}

\begin{proof}
By Taylor's theorem: $0 = f(r) = f(x_n) + f'(x_n)(r - x_n) + \frac{f''(\xi)}{2}(r-x_n)^2$ for some $\xi$ between $r$ and $x_n$. Dividing by $f'(x_n)$: $r - x_{n+1} = r - x_n + \frac{f(x_n)}{f'(x_n)} = -\frac{f''(\xi)}{2f'(x_n)}(r-x_n)^2$.
\end{proof}

\begin{leanbox}
\begin{lstlisting}
-- Newton's method as a recursive function
def newton (f df : Real -> Real) : Nat -> Real -> Real
  | 0,     x => x
  | n + 1, x => newton f df n (x - f x / df x)

-- Example: computing sqrt(2) via Newton's method for f(x) = x^2 - 2
def sqrtNewton (n : Nat) (x0 : Real) : Real :=
  newton (fun x => x * x - 2) (fun x => 2 * x) n x0

-- Starting from x0 = 1.5:
-- 1.5 -> 1.4166... -> 1.41421356... (converges to sqrt(2))
#eval sqrtNewton 5 1.5  -- (using Float approximation)
\end{lstlisting}
\end{leanbox}
\color{black}

\begin{exbox}
\begin{enumerate}
  \item Find all local extrema of $f(x) = x^4 - 4x^3 + 4x^2$.
  \item Prove that $e^x > 1 + x + x^2/2$ for all $x > 0$ using the MVT applied twice.
  \item Apply Newton's method to compute $\sqrt[3]{7}$ to 5 decimal places by hand (first two iterations).
  \item Prove formally that a cubic polynomial always has at least one real critical point.
  \item Prove the AM-GM inequality for $n$ numbers using calculus (minimize $x_1 + \cdots + x_n$ subject to $x_1 \cdots x_n = 1$).
\end{enumerate}
\end{exbox}
\color{black}

%%=============================================================================
\part{Integral Calculus}
%%=============================================================================
\color{black}

%%-----------------------------------------------------------------------------
\chapter{The Riemann Integral}
%%-----------------------------------------------------------------------------
\color{black}

\begin{remarkbox}
\textbf{Historical Note: From Cavalieri to Lebesgue}

Integration has a longer history than differentiation. \textbf{Archimedes} computed areas by the ``method of exhaustion'' --- approximating curved regions by polygons and taking limits --- in the third century BCE. This is recognisably the idea behind Riemann sums, though Archimedes had no notion of a function or a limit in the modern sense.

The connection between integration (area) and differentiation (slope) was recognised before either was fully formalised. \textbf{Barrow} (Newton's teacher) proved a geometric version of the Fundamental Theorem around 1665. Newton and Leibniz made it the cornerstone of calculus, though they understood integration primarily as ``anti-differentiation'' --- working backwards from derivatives to areas --- rather than as a limit of sums.

The rigorous limit-of-sums definition was given by \textbf{Riemann} in his 1854 Habilitation thesis, submitted to qualify as a lecturer at G\"{o}ttingen. Riemann's definition --- tagged partitions, Riemann sums, the limit as the mesh goes to zero --- is the one most calculus textbooks still use. In the same paper, Riemann showed that a function is integrable (by his definition) if and only if its set of discontinuities has ``measure zero'' in an informal sense. He exhibited functions with discontinuities dense in $[0,1]$ that are nonetheless Riemann integrable, and functions that are not.

The Darboux reformulation --- upper and lower sums, with integrability equivalent to their infimum and supremum coinciding --- was given by \textbf{Darboux} in 1875. It is logically equivalent to Riemann's definition but much easier to work with for proving theorems. We follow Darboux.

The limitations of Riemann's integral became clear in the early twentieth century. Dirichlet's function is not Riemann integrable. The limit of a uniformly bounded sequence of Riemann-integrable functions need not be Riemann integrable. \textbf{Lebesgue} (1901--1902) introduced a completely new approach: instead of partitioning the \emph{domain}, partition the \emph{range}, and measure the preimage of each range interval. This gives an integral that handles Dirichlet's function easily (integral $= 0$), allows far more general limit theorems (the Dominated Convergence Theorem), and extends naturally to abstract measure spaces far beyond $\R$. Lebesgue integration is the foundation of modern probability theory, functional analysis, and the $L^p$ spaces used in partial differential equations.

For this book we develop Riemann integration, which suffices for all the calculus of Parts III--V. The Lebesgue integral would require a separate book on measure theory, though the connection appears briefly in the exercises.
\end{remarkbox}
\color{black}

We now develop a rigorous theory of integration, answering the question: what does it mean to sum up infinitely many infinitesimally thin rectangles?

\section{Partitions and Riemann Sums}

The integral captures the notion of ``signed area under a curve.''  Our approach is via Darboux sums, which bound the integral from above and below using the supremum and infimum of the function on each subinterval.  A function is integrable precisely when the upper and lower bounds can be made to agree arbitrarily closely by refining the partition.  This is more technically convenient than the tagged-partition approach of Riemann, and it produces the same integral for all functions that matter in practice.

\color{black}\begin{center}
\begin{tikzpicture}[xscale=1.3, yscale=1.1]
  \draw[->] (-0.2,0) -- (5.5,0) node[right, font=\small] {\textcolor{black}{$x$}};
  \draw[->] (0,-0.2) -- (0,3.2) node[above, font=\small] {\textcolor{black}{$y$}};

  % Upper sums (in light blue)
  \foreach \i/\xl/\xr/\Mh in {
    1/0.3/1.3/2.35,
    2/1.3/2.3/2.55,
    3/2.3/3.3/2.1,
    4/3.3/4.3/1.35,
    5/4.3/5.1/1.1
  }{
    \fill[thmblue!20] (\xl,0) rectangle (\xr,\Mh);
    \draw[thmblue!60, thin] (\xl,0) rectangle (\xr,\Mh);
  }

  % Lower sums (in teal)
  \foreach \i/\xl/\xr/\mh in {
    1/0.3/1.3/1.55,
    2/1.3/2.3/2.05,
    3/2.3/3.3/1.45,
    4/3.3/4.3/0.95,
    5/4.3/5.1/0.75
  }{
    \fill[tealaccent!25] (\xl,0) rectangle (\xr,\mh);
    \draw[tealaccent!70, thin] (\xl,0) rectangle (\xr,\mh);
  }

  % The function curve on top
  \draw[very thick, color=navyblue, domain=0.3:5.1, samples=80]
    plot (\x, {2.0 + 0.6*sin(180*(\x-0.3)/2.4) - 0.1*(\x-2.7)*(\x-2.7)});

  % a and b labels
  \draw (0.3, 0.05) -- (0.3,-0.05) node[below, font=\small] {\textcolor{black}{$a$}};
  \draw (5.1, 0.05) -- (5.1,-0.05) node[below, font=\small] {\textcolor{black}{$b$}};

  % Legend
  \fill[thmblue!25] (0.35, 2.75) rectangle (0.75, 3.05);
  \draw[thmblue!60] (0.35, 2.75) rectangle (0.75, 3.05);
  \node[right, font=\scriptsize] at (0.8, 2.9) {\textcolor{black}{Upper sum $U(f,P)$}};
  \fill[tealaccent!30] (0.35, 2.35) rectangle (0.75, 2.65);
  \draw[tealaccent!70] (0.35, 2.35) rectangle (0.75, 2.65);
  \node[right, font=\scriptsize] at (0.8, 2.5) {\textcolor{black}{Lower sum $L(f,P)$}};
  \draw[very thick, navyblue] (0.35, 2.1) -- (0.75, 2.1);
  \node[right, font=\scriptsize] at (0.8, 2.1) {\textcolor{black}{$f(x)$}};
\end{tikzpicture}
\end{center}

\begin{defbox}{Partitions and Riemann Sums}
A \textbf{partition} of $[a,b]$ is a finite sequence $P = \{a = x_0 < x_1 < \cdots < x_n = b\}$.

The \textbf{mesh} of $P$ is $\|P\| = \max_i (x_i - x_{i-1})$.

For $f : [a,b] \to \R$, the \textbf{upper Darboux sum} and \textbf{lower Darboux sum} are:
\[ U(f, P) = \sum_{i=1}^n M_i(x_i - x_{i-1}), \quad L(f, P) = \sum_{i=1}^n m_i(x_i - x_{i-1}), \]
where $M_i = \sup_{x \in [x_{i-1},x_i]} f(x)$ and $m_i = \inf_{x \in [x_{i-1},x_i]} f(x)$.
\end{defbox}
\color{black}

\begin{defbox}{Riemann Integrability}
$f$ is \textbf{Riemann integrable} on $[a,b]$ if:
\[ \sup_P L(f,P) = \inf_P U(f,P). \]
This common value is the \textbf{Riemann integral} $\int_a^b f(x)\,dx$.
\end{defbox}
\color{black}

\begin{thmbox}{Darboux's Criterion}
$f$ is integrable on $[a,b]$ if and only if for every $\eps > 0$ there exists a partition $P$ with $U(f,P) - L(f,P) < \eps$.
\end{thmbox}
\color{black}

\begin{thmbox}{Continuous Functions Are Integrable}
If $f : [a,b] \to \R$ is continuous, then $f$ is Riemann integrable.
\end{thmbox}
\color{black}

\begin{proof}
By EVT, $f$ is bounded. Since $[a,b]$ is compact, $f$ is uniformly continuous. Given $\eps > 0$, choose $\del$ so $|x-y| < \del \Rightarrow |f(x)-f(y)| < \eps/(b-a)$. Take a partition with mesh $< \del$. Then $M_i - m_i < \eps/(b-a)$ on each subinterval, so $U - L < \eps$.
\end{proof}

\begin{thmbox}{Monotone Functions Are Integrable}
If $f : [a,b] \to \R$ is monotone, then $f$ is Riemann integrable.
\end{thmbox}
\color{black}

\begin{proof}
For increasing $f$ and partition $P$ with equal subintervals of width $(b-a)/n$: $M_i = f(x_i)$, $m_i = f(x_{i-1})$, so $U - L = \frac{b-a}{n}(f(b)-f(a)) \to 0$.
\end{proof}

\begin{leanbox}
\begin{lstlisting}
-- A partition of [a, b]
structure Partition (a b : Real) where
  points : List Real
  sorted : points.Sorted (· ≤ ·)
  starts : points.head? = some a
  ends   : points.getLast? = some b

-- Subintervals from a partition
def subintervals (P : Partition a b) : List (Real × Real) :=
  P.points.zip P.points.tail

-- Upper and lower Darboux sums (sketch)
def upperSum (f : Real -> Real) (P : Partition a b) : Real :=
  (subintervals P).foldl (fun acc (xi_prev, xi) =>
    let Mi := sorry  -- sup of f on [xi_prev, xi]
    acc + Mi * (xi - xi_prev)) 0

-- Riemann integrability
def IsIntegrable (f : Real -> Real) (a b : Real) : Prop :=
  exists I : Real,
    forall eps : Real, eps > 0 ->
      exists P : Partition a b,
        upperSum f P - I < eps /\ I - sorry < eps
\end{lstlisting}
\end{leanbox}
\color{black}

\section{Properties of the Integral}

The four basic properties below --- linearity, monotonicity, additivity, and boundedness --- are used in nearly every integral calculation.  Linearity says integration is a linear operation: you can pull out constants and split sums.  Monotonicity says the integral respects the ordering of functions.  Additivity lets us split an integral at any intermediate point, which is essential for handling piecewise-defined functions and improper integrals.

\begin{thmbox}{Basic Properties of the Riemann Integral}
If $f, g$ are integrable on $[a,b]$:
\begin{enumerate}[noitemsep]
  \item \textbf{Linearity:} $\int_a^b (cf + g) = c\int_a^b f + \int_a^b g$.
  \item \textbf{Monotonicity:} $f \leq g \Rightarrow \int_a^b f \leq \int_a^b g$.
  \item \textbf{Additivity:} $\int_a^b f = \int_a^c f + \int_c^b f$ for $a \leq c \leq b$.
  \item \textbf{Boundedness:} $\left|\int_a^b f\right| \leq \int_a^b |f| \leq M(b-a)$ where $M = \sup|f|$.
\end{enumerate}
\end{thmbox}
\color{black}

\begin{exbox}
\begin{enumerate}
  \item Compute $\int_0^1 x^2\,dx$ directly from the definition, using partitions $P_n = \{0, 1/n, 2/n, \ldots, 1\}$.
  \item Prove that if $f$ is integrable and $f = g$ everywhere except finitely many points, then $g$ is integrable with $\int f = \int g$.
  \item Show that Dirichlet's function $f(x) = 1$ for $x \in \Q$ and $0$ otherwise is not Riemann integrable on $[0,1]$.
  \item Prove the triangle inequality for integrals: $\left|\int_a^b f\right| \leq \int_a^b |f|$.
  \item Formalize the definition of \li{IsIntegrable} in Lean~4 using Lean's \li{Real} type.
\end{enumerate}
\end{exbox}
\color{black}

%%-----------------------------------------------------------------------------
\chapter{The Fundamental Theorem of Calculus}
%%-----------------------------------------------------------------------------
\color{black}

\begin{remarkbox}
\textbf{Historical Note: The Fundamental Theorem and Its Discovery}

The Fundamental Theorem of Calculus is the hinge on which the entire subject turns, yet its precise statement and proof were slow to emerge. \textbf{Isaac Barrow} proved a geometric version in his \emph{Lectiones Geometricae} (1670): if you know the area function $A(t) = \int_a^t f(x)\,dx$, then its tangent slope is $f(t)$. This is Part~1 of our theorem.

Newton understood Part~2 --- that the definite integral can be evaluated by finding an antiderivative --- and used it constantly in his \emph{Method of Fluxions} (written 1671, published 1736). For Newton, this was practically the \emph{definition} of the integral: integration \emph{was} anti-differentiation, and the connection to areas was a consequence. Leibniz, by contrast, thought of $\int$ as a summing operation from the start, and the connection to derivatives was the theorem.

The first completely rigorous proof was given by \textbf{Cauchy} in his 1823 \emph{R\'{e}sum\'{e}}, using the definition of the integral as a limit of sums that he had just introduced. Cauchy proved Part~1 by estimating the difference quotient $(F(x+h) - F(x))/h$ using the intermediate value theorem for $f$, essentially the argument we give. This required continuity of $f$, a hypothesis that Cauchy was careful to state.

The generalisation to discontinuous functions --- asking exactly which $f$ can be integrated and which integrated functions can be differentiated --- occupied analysis for the rest of the nineteenth century, culminating in Lebesgue's theory. The most general version of the FTC (in the Lebesgue setting) states that an absolutely continuous function is differentiable almost everywhere, and the derivative, integrated, returns the original function. This precise formulation required the full machinery of measure theory.
\end{remarkbox}
\color{black}

The Fundamental Theorem is the central result of calculus: it reveals that differentiation and integration are inverse operations, connecting the two branches of the subject.

\section{FTC Part 1: The Integral as an Antiderivative}

The first part of the Fundamental Theorem reveals that differentiation and integration are inverse operations: if you integrate a function $f$ from a fixed base point $a$ to a variable upper limit $x$, the resulting function $F(x) = \int_a^x f(t)\,dt$ is an antiderivative of $f$.  This is a deep fact, not a tautology.  The integral is defined as a limit of sums; the derivative is a limit of difference quotients.  That these two limiting operations undo each other is the central miracle of calculus.

\begin{defbox}{The Integral Function}
For $f$ integrable on $[a,b]$, define the \textbf{integral function}:
\[ F(x) := \int_a^x f(t)\,dt, \quad x \in [a,b]. \]
\end{defbox}
\color{black}

\begin{thmbox}{Fundamental Theorem of Calculus, Part 1}
If $f$ is continuous on $[a,b]$, then $F$ is differentiable on $(a,b)$ and $F'(x) = f(x)$.
\end{thmbox}
\color{black}

\begin{proof}
For $h \neq 0$:
\[ \frac{F(x+h) - F(x)}{h} = \frac{1}{h}\int_x^{x+h} f(t)\,dt. \]
Since $f$ is continuous at $x$, for any $\eps > 0$ there exists $\del > 0$ so $|f(t) - f(x)| < \eps$ when $|t - x| < \del$. For $|h| < \del$:
\[ \left|\frac{1}{h}\int_x^{x+h} f(t)\,dt - f(x)\right| = \left|\frac{1}{h}\int_x^{x+h}(f(t) - f(x))\,dt\right| \leq \frac{|h|\eps}{|h|} = \eps. \]
\end{proof}

\section{FTC Part 2: The Evaluation Theorem}

Part 2 is what makes integration computable in practice.  Rather than computing a limit of Riemann sums (which is laborious even for simple functions), we simply find any antiderivative $G$ and evaluate $G(b) - G(a)$.  The two parts together mean: \emph{to integrate, anti-differentiate}.  This transforms the problem of integration into the problem of recognising derivatives --- which is exactly what the integration engine in Part~VII formalises as a pattern-matching algorithm.

\begin{thmbox}{Fundamental Theorem of Calculus, Part 2}
If $f$ is continuous on $[a,b]$ and $G$ is any antiderivative of $f$ (i.e., $G' = f$), then:
\[ \int_a^b f(x)\,dx = G(b) - G(a). \]
\end{thmbox}
\color{black}

\begin{proof}
Let $F(x) = \int_a^x f(t)\,dt$. By Part 1, $F' = f = G'$, so $G - F$ is constant (by MVT, $(G-F)'=0 \Rightarrow G-F=c$). Then $\int_a^b f = F(b) = G(b) - G(a) + F(b) - G(b) + G(a) = G(b) - G(a)$ (using $F(a) = 0$).
\end{proof}

The power of FTC Part 2 cannot be overstated: instead of computing limits of Riemann sums, we need only find an antiderivative. This is why the integration techniques of Chapter~17 matter.

\begin{leanbox}
\begin{lstlisting}
-- FTC Part 2 in Lean: if G' = f and f is continuous, then integral = G(b) - G(a)
-- We state it in terms of our HasDerivAt definition
theorem ftc_part2 (f G : Real -> Real) (a b : Real)
    (hcont : forall x, a <= x -> x <= b -> ContinuousAt f x)
    (hderiv : forall x, a < x -> x < b -> HasDerivAt G x (f x))
    (hint : IsIntegrable f a b) :
    integral f a b = G b - G a := by
  sorry  -- Full proof requires the full Riemann integral machinery
  -- Sketch: F(x) := integral f a x satisfies F' = f by FTC Part 1
  -- G' = f = F', so (G - F)' = 0, so G - F = constant C
  -- At x = a: G(a) - F(a) = G(a) - 0 = G(a), so C = G(a)
  -- At x = b: integral = F(b) = G(b) - C = G(b) - G(a)
\end{lstlisting}
\end{leanbox}
\color{black}

\begin{exbox}
\begin{enumerate}
  \item Use FTC Part 2 to evaluate $\int_0^\pi \sin x\,dx$, $\int_1^e (1/x)\,dx$, and $\int_0^1 e^x\,dx$.
  \item Prove that $\frac{d}{dx}\int_a^{g(x)} f(t)\,dt = f(g(x)) \cdot g'(x)$ (chain rule + FTC Part 1).
  \item Prove: if $F' = f$ everywhere and $f \geq 0$, then $F$ is non-decreasing.
  \item Show: FTC Part 1 can fail if $f$ is not continuous. (Hint: consider $f = 1_\Q$ on $[0,1]$.)
  \item Formalize FTC Part 1 in Lean~4 assuming the necessary lemmas.
\end{enumerate}
\end{exbox}
\color{black}

%%-----------------------------------------------------------------------------
\chapter{Techniques of Integration}
%%-----------------------------------------------------------------------------
\color{black}

Each integration technique is a theorem---a rule that transforms one integral into another, simpler one. This chapter is also the direct mathematical foundation for the integration engine in Part~VII.

\section{Substitution (Change of Variables)}

$U$-substitution is the integration counterpart of the chain rule: just as $(f \circ g)' = (f' \circ g) \cdot g'$, integrating $f(g(x)) g'(x)$ reduces to integrating $f(u)$ after the substitution $u = g(x)$.  Geometrically, substitution rescales the axis: the factor $g'(x)$ accounts for how the substitution stretches or compresses the $x$-axis.  In formal terms, it is the change-of-variables theorem applied to a one-dimensional integral.

\begin{thmbox}{Integration by Substitution}
If $g : [a,b] \to \R$ is differentiable with continuous $g'$, and $f$ is continuous on the range of $g$, then:
\[ \int_a^b f(g(x))\,g'(x)\,dx = \int_{g(a)}^{g(b)} f(u)\,du. \]
\end{thmbox}
\color{black}

\begin{proof}
Let $F$ be an antiderivative of $f$. By the chain rule, $\frac{d}{dx}F(g(x)) = f(g(x))g'(x)$. By FTC Part 2: LHS $= F(g(b)) - F(g(a)) =$ RHS.
\end{proof}

\begin{example}
$\int_0^1 2x e^{x^2}\,dx$. Let $u = x^2$, $du = 2x\,dx$. When $x=0$, $u=0$; when $x=1$, $u=1$:
$\int_0^1 e^u\,du = e^1 - e^0 = e - 1$.
\end{example}

\section{Integration by Parts}

Integration by parts is the integration counterpart of the product rule.  It trades one integral for another: $\int u\,dv = uv - \int v\,du$.  The art is in choosing $u$ and $dv$ so the new integral $\int v\,du$ is simpler than the original.  The LIATE heuristic (Logarithm, Inverse trig, Algebraic, Trig, Exponential) ranks which factor should be $u$ --- it encodes centuries of experience about which choices tend to reduce complexity.

\begin{thmbox}{Integration by Parts}
If $f$ and $g$ are differentiable with continuous derivatives on $[a,b]$:
\[ \int_a^b f(x)g'(x)\,dx = [f(x)g(x)]_a^b - \int_a^b f'(x)g(x)\,dx. \]
\end{thmbox}
\color{black}

\begin{proof}
By the product rule: $(fg)' = f'g + fg'$. Integrate both sides from $a$ to $b$ and apply FTC.
\end{proof}

\begin{example}
$\int x e^x\,dx$. Let $f = x$, $g' = e^x$, so $f' = 1$, $g = e^x$:
$= xe^x - \int e^x\,dx = xe^x - e^x + C = (x-1)e^x + C$.
\end{example}

\section{Partial Fraction Decomposition}

Rational functions $p(x)/q(x)$ can be integrated by decomposing into partial fractions.

\begin{thmbox}{Partial Fraction Decomposition}
Every rational function $p(x)/q(x)$ (with $\deg p < \deg q$) over $\R$ can be written as a sum of terms of the form $\frac{A}{(x-a)^k}$ and $\frac{Bx + C}{(x^2 + bx + c)^k}$ (for irreducible quadratics $x^2 + bx + c$).
\end{thmbox}
\color{black}

The proof uses the factorization of $q(x)$ over $\R$ and the algebraic uniqueness of the decomposition.

\begin{example}
$\int \frac{1}{x^2-1}\,dx = \int \frac{1}{(x-1)(x+1)}\,dx = \int \left(\frac{1/2}{x-1} - \frac{1/2}{x+1}\right)dx = \frac{1}{2}\ln|x-1| - \frac{1}{2}\ln|x+1| + C$.
\end{example}

\section{Trigonometric Substitutions}

For integrands involving $\sqrt{a^2-x^2}$, $\sqrt{a^2+x^2}$, or $\sqrt{x^2-a^2}$, use the substitutions $x = a\sin\theta$, $x = a\tan\theta$, $x = a\sec\theta$ respectively to remove the square root.

\section{Trigonometric Integrals}

Integrals of the form $\int \sin^m x \cos^n x\,dx$ are handled by:
\begin{itemize}[noitemsep]
  \item If $m$ or $n$ is odd: save one factor and use $\sin^2 + \cos^2 = 1$ to convert.
  \item If both even: use double-angle formulas $\sin^2 x = (1-\cos 2x)/2$ and $\cos^2 x = (1+\cos 2x)/2$.
\end{itemize}

\begin{exbox}
\begin{enumerate}
  \item Evaluate $\int_0^{\pi/2} \sin^3 x \cos^2 x\,dx$ using the odd-power strategy.
  \item Evaluate $\int \frac{x^2+1}{x^3-x}\,dx$ using partial fractions.
  \item Evaluate $\int \sqrt{1-x^2}\,dx$ using trigonometric substitution.
  \item Derive the reduction formula: $\int \sin^n x\,dx = -\frac{\sin^{n-1}x\cos x}{n} + \frac{n-1}{n}\int\sin^{n-2}x\,dx$.
  \item Use integration by parts repeatedly to show $\int x^n e^x\,dx$ satisfies a recursive formula.
  \item Show $\int_0^\infty x^{n-1}e^{-x}\,dx = (n-1)!$ using integration by parts (this is $\Gamma(n)$).
\end{enumerate}
\end{exbox}
\color{black}

%%-----------------------------------------------------------------------------
\chapter{Improper Integrals}
%%-----------------------------------------------------------------------------
\color{black}

\begin{defbox}{Improper Integrals of Type I and II}
\textbf{Type I} (infinite limits):
$\int_a^\infty f(x)\,dx := \lim_{b\to\infty} \int_a^b f(x)\,dx$ (if the limit exists).

\textbf{Type II} (unbounded integrand):
$\int_a^b f(x)\,dx := \lim_{c\to a^+}\int_c^b f(x)\,dx$ (if $f$ is unbounded near $a$).

An improper integral \textbf{converges} if the limit is finite; otherwise it \textbf{diverges}.
\end{defbox}
\color{black}

\begin{thmbox}{Comparison Test for Improper Integrals}
If $0 \leq f(x) \leq g(x)$ for $x \geq a$:
\begin{itemize}[noitemsep]
  \item If $\int_a^\infty g$ converges, so does $\int_a^\infty f$.
  \item If $\int_a^\infty f$ diverges, so does $\int_a^\infty g$.
\end{itemize}
\end{thmbox}
\color{black}

\begin{thmbox}{The $p$-integral}
$\int_1^\infty \frac{1}{x^p}\,dx$ converges iff $p > 1$. Similarly, $\int_0^1 \frac{1}{x^p}\,dx$ converges iff $p < 1$.
\end{thmbox}
\color{black}

\begin{defbox}{The Gamma Function}
For $n > 0$: $\Gamma(n) := \int_0^\infty x^{n-1} e^{-x}\,dx$.
\end{defbox}
\color{black}

\begin{thmbox}{$\Gamma(n) = (n-1)!$ for $n \in \N$}
\textbf{Proof.}
$\Gamma(1) = \int_0^\infty e^{-x}\,dx = \bigl[-e^{-x}\bigr]_0^\infty = 1$.
Integration by parts with $u=x^n$, $dv=e^{-x}\,dx$ gives
$\Gamma(n+1) = n\,\Gamma(n)$.
By induction $\Gamma(n) = (n-1)!$ for all $n \in \N$. \qed
\end{thmbox}
\color{black}

\begin{exbox}
\begin{enumerate}
  \item Determine convergence: $\int_1^\infty e^{-x^2}\,dx$, $\int_0^1 \ln x\,dx$, $\int_0^\infty \frac{\sin x}{x}\,dx$.
  \item Prove $\int_0^\infty e^{-x^2}\,dx = \sqrt{\pi}/2$ (Gaussian integral, using the double-integral technique).
  \item Prove the functional equation $\Gamma(n+1) = n\Gamma(n)$ for all $n > 0$.
  \item Show $\int_{-\infty}^\infty \frac{1}{1+x^2}\,dx = \pi$ and interpret it probabilistically.
\end{enumerate}
\end{exbox}
\color{black}

%%=============================================================================
\part{Sequences and Series}
%%=============================================================================
\color{black}

%%-----------------------------------------------------------------------------
\chapter{Infinite Series}
%%-----------------------------------------------------------------------------
\color{black}

\begin{remarkbox}
\textbf{Historical Note: The Paradoxes of the Infinite and the Birth of Convergence Theory}

Infinite series were used long before they were understood. \textbf{Archimedes} summed the geometric series $\sum (1/4)^n = 4/3$ to compute the area of a parabolic segment. Medieval scholars debated Zeno's paradox of Achilles and the tortoise --- which is really just the statement that $\sum_{n=1}^\infty (1/2)^n = 1$ --- without having the tools to resolve it rigorously.

The danger of infinite series became apparent when \textbf{Grandi} (1703) noted the series $1 - 1 + 1 - 1 + \cdots$. Grouping as $(1-1) + (1-1) + \cdots$ gives $0$; grouping as $1 + (-1+1) + (-1+1) + \cdots$ gives $1$. Leibniz argued, by probability, the ``true'' sum was $1/2$ (the average of $0$ and $1$). Euler, who was cavalier about convergence throughout his career, evaluated the series $1 + 2 + 3 + 4 + \cdots = -1/12$ using his divergent series techniques. (This is not merely wrong: in the context of \emph{analytic continuation} of the Riemann zeta function, it is a meaningful statement. But treating it as a convergent series in the ordinary sense is incorrect.)

\textbf{Cauchy} imposed order in 1821 by insisting on a precise definition of convergence via partial sums, as we give below. He proved that a series converges if and only if its partial sums form a Cauchy sequence (the Cauchy criterion for series), and derived the ratio and root tests. \textbf{Abel} (1826) proved the theorem bearing his name: if a power series $\sum a_n x^n$ converges at $x = R$, then it converges uniformly on $[0, R]$. Abel used this to correct Cauchy's erroneous claim that a convergent series of continuous functions is always continuous --- continuity at the endpoint $R$ requires uniform convergence, not just pointwise.

The distinction between absolute and conditional convergence was clarified by \textbf{Dirichlet} (1837), who proved that a conditionally convergent series can be rearranged to converge to any value --- the \textbf{Riemann rearrangement theorem} we prove in this chapter (Riemann gave a complete proof in 1854). This is one of the most counterintuitive results in analysis: the infinite analogue of addition is not commutative for conditionally convergent series.
\end{remarkbox}
\color{black}

\begin{defbox}{Convergence of a Series}
Given $(a_n)_{n=1}^\infty$, define the partial sums $S_n = \sum_{k=1}^n a_k$.
The series $\sum_{n=1}^\infty a_n$ \textbf{converges} to $S$ if $S_n \to S$; otherwise it \textbf{diverges}.
\end{defbox}
\color{black}

\begin{thmbox}{Divergence Test}
If $\sum a_n$ converges, then $a_n \to 0$.
Equivalently, if $a_n \not\to 0$, the series diverges.
\end{thmbox}
\color{black}

\begin{thmbox}{Geometric Series}
$\sum_{n=0}^\infty r^n = \frac{1}{1-r}$ for $|r| < 1$, and diverges for $|r| \geq 1$.
\end{thmbox}
\color{black}

\begin{proof}
$S_n = 1 + r + \cdots + r^n = \frac{1-r^{n+1}}{1-r}$. For $|r| < 1$, $r^{n+1} \to 0$, so $S_n \to \frac{1}{1-r}$.
\end{proof}

\section{Convergence Tests}

Deciding whether an infinite series converges is often harder than evaluating it.  The convergence tests below are organised by the type of information they extract: the ratio and root tests measure the asymptotic growth rate of individual terms; the integral test connects series to improper integrals; the alternating series test exploits sign cancellation.  No single test works for everything --- part of the skill of analysis is recognising which test applies to which series.

\begin{thmbox}{Tests for Convergence}
For a series $\sum a_n$ with $a_n \geq 0$:
\begin{enumerate}[noitemsep]
  \item \textbf{Comparison:} If $a_n \leq b_n$ and $\sum b_n$ converges, so does $\sum a_n$.
  \item \textbf{Ratio Test:} If $\lim |a_{n+1}/a_n| = L$: converges if $L < 1$, diverges if $L > 1$.
  \item \textbf{Root Test:} If $\lim |a_n|^{1/n} = L$: same conclusion.
  \item \textbf{Integral Test:} If $f(n) = a_n$ with $f$ positive, decreasing: $\sum a_n$ and $\int_1^\infty f$ either both converge or both diverge.
  \item \textbf{Alternating Series Test:} If $b_n \searrow 0$, then $\sum (-1)^n b_n$ converges.
\end{enumerate}
\end{thmbox}
\color{black}

\begin{defbox}{Absolute vs.\ Conditional Convergence}
$\sum a_n$ \textbf{converges absolutely} if $\sum |a_n|$ converges. Absolute convergence implies convergence.
$\sum a_n$ \textbf{converges conditionally} if it converges but not absolutely.
\end{defbox}
\color{black}

\begin{thmbox}{Riemann Rearrangement Theorem}
A conditionally convergent series can be rearranged to converge to any real number $L$, or to diverge.
\end{thmbox}
\color{black}

\begin{exbox}
\begin{enumerate}
  \item Determine convergence of $\sum_{n=1}^\infty \frac{n^2}{2^n}$, $\sum \frac{(-1)^n}{\sqrt{n}}$, $\sum \frac{1}{n \ln n}$.
  \item Prove the ratio test from first principles.
  \item Show that $\sum 1/n^2 = \pi^2/6$ (Basel problem) using the fact that $\pi^2/6 = \int_0^1\int_0^1 \frac{1}{1-xy}\,dx\,dy$.
  \item Construct an explicit rearrangement of $\sum (-1)^n/n$ that converges to $2$.
  \item Prove: if $\sum a_n$ converges absolutely and $(b_n)$ is bounded, then $\sum a_n b_n$ converges.
\end{enumerate}
\end{exbox}
\color{black}

%%-----------------------------------------------------------------------------
\chapter{Power Series}
%%-----------------------------------------------------------------------------
\color{black}

\begin{defbox}{Power Series and Radius of Convergence}
A \textbf{power series} centered at $a$ is $\sum_{n=0}^\infty c_n(x-a)^n$.

Its \textbf{radius of convergence} is $R = 1/\limsup_{n\to\infty} |c_n|^{1/n}$ (the Cauchy--Hadamard formula). The series converges absolutely for $|x-a| < R$ and diverges for $|x-a| > R$.
\end{defbox}
\color{black}

\begin{thmbox}{Term-by-Term Differentiation and Integration}
If $f(x) = \sum_{n=0}^\infty c_n(x-a)^n$ has radius of convergence $R > 0$, then for $|x-a| < R$:
\begin{align*}
  f'(x) &= \sum_{n=1}^\infty n c_n(x-a)^{n-1}, \quad \text{also with radius of convergence } R. \\
  \int_a^x f(t)\,dt &= \sum_{n=0}^\infty \frac{c_n}{n+1}(x-a)^{n+1}.
\end{align*}
\end{thmbox}
\color{black}

\begin{proof}
The key lemma is that the derived series has the same radius of convergence (since $\lim n^{1/n} = 1$). The equality $f' = \sum nc_n(x-a)^{n-1}$ is proved by showing uniform convergence on compact subsets and passing the limit through the derivative.
\end{proof}

\begin{thmbox}{Uniqueness of Power Series}
If $f(x) = \sum_{n=0}^\infty c_n(x-a)^n$ for $|x-a| < R$, then $c_n = f^{(n)}(a)/n!$.
\end{thmbox}
\color{black}

\begin{exbox}
\begin{enumerate}
  \item Find the radius of convergence of $\sum n! x^n$, $\sum x^n/n!$, $\sum x^n/n$.
  \item Use the geometric series $1/(1-x) = \sum x^n$ and term-by-term integration to derive the series for $\ln(1+x)$.
  \item Derive the series for $\arctan x$ from $1/(1+x^2) = \sum (-1)^n x^{2n}$.
  \item Prove that $\sum x^n/n!$ converges for all $x \in \R$ and defines a differentiable function equal to its own derivative.
\end{enumerate}
\end{exbox}
\color{black}

%%-----------------------------------------------------------------------------
\chapter{Taylor and Maclaurin Series}
%%-----------------------------------------------------------------------------
\color{black}

\begin{defbox}{Taylor Series}
If $f$ is infinitely differentiable at $a$, its \textbf{Taylor series} is:
\[ \sum_{n=0}^\infty \frac{f^{(n)}(a)}{n!}(x-a)^n. \]
A function is \textbf{analytic} at $a$ if it equals its Taylor series in some neighborhood of $a$.
\end{defbox}
\color{black}

The standard Taylor series (all centered at $0$, valid for the indicated $x$):

\begin{align*}
  e^x &= \sum_{n=0}^\infty \frac{x^n}{n!} & &\text{all }x \\
  \sin x &= \sum_{n=0}^\infty \frac{(-1)^n x^{2n+1}}{(2n+1)!} & &\text{all }x \\
  \cos x &= \sum_{n=0}^\infty \frac{(-1)^n x^{2n}}{(2n)!} & &\text{all }x \\
  \ln(1+x) &= \sum_{n=1}^\infty \frac{(-1)^{n-1} x^n}{n} & &|x| < 1 \\
  (1+x)^\alpha &= \sum_{n=0}^\infty \binom{\alpha}{n} x^n & &|x| < 1
\end{align*}

\begin{remarkbox}
Not every $C^\infty$ function is analytic. The function $f(x) = e^{-1/x^2}$ (extended by $f(0) = 0$) is $C^\infty$ everywhere, yet its Taylor series at $0$ is identically $0 \neq f(x)$ for $x \neq 0$. Analytic functions form a strictly smaller class than smooth functions.
\end{remarkbox}
\color{black}

\begin{leanbox}
\begin{lstlisting}
-- We can define e^x via its Taylor series and prove properties
-- Taylor polynomials for e^x
def expTaylor (n : Nat) (x : Real) : Real :=
  Finset.sum (Finset.range (n + 1)) (fun k =>
    x ^ k / Nat.factorial k)

-- The k-th term
def expTerm (k : Nat) (x : Real) : Real :=
  x ^ k / Nat.factorial k

-- Ratio test: |a_{k+1}/a_k| = |x|/(k+1) -> 0 for any fixed x
theorem expTaylor_converges (x : Real) :
    SeqLimit (fun n => expTaylor n x) (Real.exp x) := by
  sorry  -- requires series convergence machinery from Ch. 19-20

-- e^x satisfies (e^x)' = e^x
theorem exp_deriv (x : Real) : HasDerivAt Real.exp x (Real.exp x) := by
  sorry  -- follows from term-by-term differentiation of the Taylor series
\end{lstlisting}
\end{leanbox}
\color{black}

\begin{exbox}
\begin{enumerate}
  \item Derive the Taylor series for $\arcsin x$ from $1/\sqrt{1-x^2} = \sum \binom{-1/2}{n}(-x^2)^n$.
  \item Use series to evaluate $\int_0^1 e^{-x^2}\,dx$ to 4 decimal places.
  \item Prove $e^{ix} = \cos x + i\sin x$ (Euler's formula) using the Taylor series.
  \item Show that the Taylor series for $\ln(1+x)$ at $x = 1$ converges to $\ln 2$ by the alternating series test.
  \item Prove the addition formula for $\exp$: $e^{x+y} = e^x \cdot e^y$, using series multiplication.
\end{enumerate}
\end{exbox}
\color{black}

%%=============================================================================
\part{Multivariable Calculus}
%%=============================================================================
\color{black}

%%-----------------------------------------------------------------------------
\chapter{\texorpdfstring{$\R^n$}{Rn} and Its Topology}
%%-----------------------------------------------------------------------------
\color{black}

\section{Vectors in \texorpdfstring{$\R^n$}{Rn}}

Multivariable calculus generalises everything from Part~III onward to functions of $n$ real variables.  The domain is no longer a line but $n$-dimensional Euclidean space $\R^n$.  The correct notion of distance is the Euclidean norm, and the correct notion of angle is the dot product.  Cauchy--Schwarz --- the single most useful inequality in analysis --- bounds the dot product by the product of norms, and its proof is an elegant one-liner using the non-negativity of a quadratic.  In Lean~4, we model $\R^n$ as the function type \li{Fin n -> Real}: a vector is just a function from $n$ indices to real values.

\begin{defbox}{Euclidean Space $\R^n$}
$\R^n$ is the set of $n$-tuples $(x_1, \ldots, x_n)$ with $x_i \in \R$. In Lean~4 we model it as \li{Fin n -> Real}.

The \textbf{Euclidean norm} is $\|x\| = \sqrt{x_1^2 + \cdots + x_n^2}$.
The \textbf{dot product} is $x \cdot y = \sum_{i=1}^n x_i y_i$.
The \textbf{Euclidean distance} is $d(x, y) = \|x - y\|$.
\end{defbox}
\color{black}

\begin{thmbox}{Cauchy--Schwarz Inequality}
For all $x, y \in \R^n$: $|x \cdot y| \leq \|x\|\,\|y\|$.
\end{thmbox}
\color{black}

\begin{proof}
For any $t \in \R$: $0 \leq \|x + ty\|^2 = \|x\|^2 + 2t(x \cdot y) + t^2\|y\|^2$. This quadratic in $t$ is non-negative, so its discriminant satisfies $4(x \cdot y)^2 - 4\|x\|^2\|y\|^2 \leq 0$.
\end{proof}

\begin{thmbox}{Triangle Inequality in $\R^n$}
$\|x + y\| \leq \|x\| + \|y\|$.
\end{thmbox}
\color{black}

\begin{proof}
$\|x+y\|^2 = \|x\|^2 + 2x \cdot y + \|y\|^2 \leq \|x\|^2 + 2\|x\|\|y\| + \|y\|^2 = (\|x\| + \|y\|)^2$.
\end{proof}

\begin{leanbox}
\begin{lstlisting}
-- R^n as functions from Fin n to Real
def Vec (n : Nat) := Fin n -> Real

def dotProduct {n : Nat} (v w : Vec n) : Real :=
  Finset.sum Finset.univ (fun i => v i * w i)

def normSq {n : Nat} (v : Vec n) : Real :=
  dotProduct v v

def euclidNorm {n : Nat} (v : Vec n) : Real :=
  Real.sqrt (normSq v)

-- Cauchy-Schwarz
theorem cauchy_schwarz {n : Nat} (v w : Vec n) :
    dotProduct v w ^ 2 <= normSq v * normSq w := by
  -- Classic proof via discriminant
  have h : forall t : Real,
      0 <= normSq (fun i => v i + t * w i) := fun t => by
    simp [normSq, dotProduct]; positivity
  -- The LHS is a non-negative quadratic in t: ||v||^2 + 2t(v.w) + t^2||w||^2
  -- Discriminant must be <= 0
  sorry
\end{lstlisting}
\end{leanbox}
\color{black}

\section{Open Sets, Convergence, and Compactness in \texorpdfstring{$\R^n$}{Rn}}

All the topological concepts from Chapter~9 carry over to $\R^n$ essentially unchanged, with the absolute value $|x - a|$ replaced by the Euclidean norm $\|x - a\|$.  Open balls, open sets, closed sets, convergence, and compactness are all defined by the same pattern.  The Heine--Borel theorem --- compact iff closed and bounded --- holds in $\R^n$ for exactly the same reason as in $\R$: completeness allows the bisection argument to work in each coordinate independently.

\begin{defbox}{Topology of $\R^n$}
An \textbf{open ball} in $\R^n$ is $B(a, r) = \{x \in \R^n : \|x - a\| < r\}$.
A set $U \subseteq \R^n$ is \textbf{open} if every point has a neighborhood ball inside $U$.
A set is \textbf{closed} if its complement is open.
A sequence $(x_k)$ in $\R^n$ \textbf{converges} to $L$ if $\|x_k - L\| \to 0$.
\end{defbox}
\color{black}

\begin{thmbox}{Heine--Borel in $\R^n$}
A subset of $\R^n$ is compact iff it is closed and bounded.
\end{thmbox}
\color{black}

The proof follows the same bisection argument as in $\R$, now bisecting an $n$-dimensional box.

\begin{exbox}
\begin{enumerate}
  \item Prove: $(x_k) \to L$ in $\R^n$ iff each component $(x_k)_i \to L_i$ in $\R$.
  \item Prove the triangle inequality in $\R^n$ from Cauchy--Schwarz.
  \item Show that the closed unit ball $\{x : \|x\| \leq 1\}$ is compact in $\R^n$.
  \item Define continuity of $f : \R^m \to \R^n$ at a point and prove compositions of continuous maps are continuous.
\end{enumerate}
\end{exbox}
\color{black}

%%-----------------------------------------------------------------------------
\chapter{Partial Derivatives and Differentiability}
%%-----------------------------------------------------------------------------
\color{black}

\section{Partial Derivatives}

For a function of $n$ variables, we can differentiate in each coordinate direction separately by holding all other variables constant.  The result is called a partial derivative.  Partial derivatives are easy to compute --- just apply the single-variable rules with all other variables treated as constants --- but they carry less information than the total derivative: a function can have all partial derivatives at a point yet still fail to be differentiable there, as the classic counterexample below shows.

\begin{defbox}{Partial Derivative}
For $f : \R^n \to \R$ and the $i$-th standard basis vector $e_i$:
\[ \frac{\partial f}{\partial x_i}(a) := \lim_{h \to 0} \frac{f(a + he_i) - f(a)}{h}. \]
This is the derivative of $f$ in the direction of the $i$-th coordinate axis, with all other coordinates fixed.
\end{defbox}
\color{black}

\section{Total Differentiability}

A function may have all partial derivatives yet still fail to be ``differentiable'' in a meaningful sense. The right notion is the \emph{total derivative}.

\begin{defbox}{Total Derivative (Fr\'echet Derivative)}
$f : \R^n \to \R^m$ is \textbf{differentiable} at $a$ if there exists a linear map $Df(a) : \R^n \to \R^m$ such that:
\[ \lim_{h \to 0} \frac{\|f(a+h) - f(a) - Df(a)(h)\|}{\|h\|} = 0. \]
$Df(a)$ is the \textbf{total derivative} (or Fr\'echet derivative). Its matrix in standard bases is the \textbf{Jacobian} $J_{ij} = \partial f_i/\partial x_j$.
\end{defbox}
\color{black}

\begin{thmbox}{Differentiability Implies Existence of Partials (But Not Vice Versa)}
If $f$ is differentiable at $a$, then all partial derivatives exist at $a$ and $Df(a)$ is the Jacobian matrix.
The converse fails: the function $f(x,y) = xy/(x^2+y^2)$ (with $f(0,0)=0$) has both partial derivatives at the origin but is not differentiable there.
\end{thmbox}
\color{black}

\begin{thmbox}{Continuous Partials $\Rightarrow$ Differentiability}
If all partial derivatives of $f$ exist and are continuous at $a$, then $f$ is differentiable at $a$.
\end{thmbox}
\color{black}

\begin{proof}
By the MVT applied in each coordinate direction. Write $f(a+h) - f(a) = \sum_i [f(a + h_1 e_1 + \cdots + h_i e_i) - f(a + h_1 e_1 + \cdots + h_{i-1} e_{i-1})]$; apply the MVT to each term; use continuity of the partials to show the error is $o(\|h\|)$.
\end{proof}

\begin{leanbox}
\begin{lstlisting}
-- Partial derivative: fix all coordinates except the i-th
def partialDeriv (f : Vec n -> Real) (i : Fin n) (a : Vec n) : Real :=
  let g := fun (t : Real) => f (fun j => if j = i then t else a j)
  -- This is just the ordinary derivative of g at a i
  Classical.choose (HasDerivAt g (a i)) -- sketch

-- Total (Frechet) derivative as a linear map
def HasTotalDerivAt (f : Vec n -> Real) (a : Vec n) (L : Vec n -> Real) : Prop :=
  (forall v w : Vec n, L (fun i => v i + w i) = L v + L w) /\   -- linearity
  (forall c : Real, forall v, L (fun i => c * v i) = c * L v) /\ -- homogeneity
  forall eps : Real, eps > 0 ->
    exists delta : Real, delta > 0 /\
      forall h : Vec n, euclidNorm h < delta -> euclidNorm h > 0 ->
        |f (fun i => a i + h i) - f a - L h| < eps * euclidNorm h
\end{lstlisting}
\end{leanbox}
\color{black}

\begin{exbox}
\begin{enumerate}
  \item Compute all partial derivatives of $f(x,y) = x^2 y + \sin(xy)$.
  \item Show that $f(x,y) = xy/\sqrt{x^2+y^2}$ (for $(x,y) \neq 0$) is not differentiable at the origin despite having partial derivatives there.
  \item Prove: if $f : \R^n \to \R$ is differentiable at $a$, then $f$ is continuous at $a$.
  \item Prove the multivariable chain rule: if $g : \R^m \to \R^n$ and $f : \R^n \to \R^p$ are differentiable, then $D(f \circ g)(a) = Df(g(a)) \circ Dg(a)$.
\end{enumerate}
\end{exbox}
\color{black}

%%-----------------------------------------------------------------------------
\chapter{The Gradient and Directional Derivatives}
%%-----------------------------------------------------------------------------
\color{black}

\begin{defbox}{Directional Derivative and Gradient}
The \textbf{directional derivative} of $f : \R^n \to \R$ at $a$ in direction $v$ (with $\|v\| = 1$) is:
\[ D_v f(a) := \lim_{t \to 0} \frac{f(a + tv) - f(a)}{t}. \]

The \textbf{gradient} is the vector of partial derivatives:
\[ \nabla f(a) = \left(\frac{\partial f}{\partial x_1}(a),\; \ldots,\; \frac{\partial f}{\partial x_n}(a)\right). \]
\end{defbox}
\color{black}

\begin{thmbox}{Gradient and Directional Derivatives}
If $f$ is differentiable at $a$, then $D_v f(a) = \nabla f(a) \cdot v$.
The gradient $\nabla f(a)$ points in the direction of steepest ascent.
\end{thmbox}
\color{black}

\begin{proof}
By the chain rule: $D_v f(a) = Df(a)(v) = \nabla f(a) \cdot v$. By Cauchy--Schwarz, $|\nabla f(a) \cdot v| \leq \|\nabla f(a)\|$, with equality when $v = \nabla f(a)/\|\nabla f(a)\|$.
\end{proof}

\begin{thmbox}{Chain Rule in $\R^n$}
If $\mathbf{r}(t) : \R \to \R^n$ is differentiable and $f : \R^n \to \R$ is differentiable, then:
\[ (f \circ \mathbf{r})'(t) = \nabla f(\mathbf{r}(t)) \cdot \mathbf{r}'(t). \]
\end{thmbox}
\color{black}

\begin{exbox}
\begin{enumerate}
  \item Find the gradient of $f(x,y,z) = x^2 + yz + e^{xz}$.
  \item Find the direction of steepest ascent for $f(x,y) = x^2 - y^2$ at $(1,1)$.
  \item Prove that if $f$ is constant on a surface $g(x) = c$, then $\nabla f(a)$ is proportional to $\nabla g(a)$ at any $a$ on the surface (this is the geometric basis for Lagrange multipliers).
  \item Compute the directional derivative of $f(x,y) = \sin(x+y)$ at $(0,0)$ in direction $(1/\sqrt{2}, 1/\sqrt{2})$.
\end{enumerate}
\end{exbox}
\color{black}

%%-----------------------------------------------------------------------------
\chapter{Higher-Order Partials}
%%-----------------------------------------------------------------------------
\color{black}

\begin{defbox}{Mixed Partial Derivatives and the Hessian}
The \textbf{Hessian matrix} of $f : \R^n \to \R$ at $a$ is $H_{ij}(a) = \frac{\partial^2 f}{\partial x_i \partial x_j}(a)$.
\end{defbox}
\color{black}

\begin{thmbox}{Clairaut's Theorem (Equality of Mixed Partials)}
If $\frac{\partial^2 f}{\partial x_i \partial x_j}$ and $\frac{\partial^2 f}{\partial x_j \partial x_i}$ are both continuous at $a$, then they are equal there.
\end{thmbox}
\color{black}

\begin{proof}
By a careful limit argument involving the mixed difference quotient $f(a+he_i+ke_j) - f(a+he_i) - f(a+ke_j) + f(a)$. Apply the MVT twice.
\end{proof}

The Hessian is symmetric (when mixed partials are continuous): $H = H^T$.


\begin{thmbox}{Second-Order Taylor Expansion in $\R^n$}
If $f:\R^n\to\R$ is twice continuously differentiable at $a$, then:
\[f(a+h) = f(a) + \nabla f(a)\cdot h + \tfrac{1}{2}h^T H(a)\, h + o(\|h\|^2),\]
where $H(a)$ is the Hessian matrix at $a$.
\end{thmbox}
\color{black}

\begin{proof}
Apply the scalar Taylor theorem to $g(t)=f(a+th)$ at $t=0$ up to second order,
then expand using the chain rule twice.
The first derivative gives $\nabla f(a)\cdot h$; the second derivative gives
$h^T H(a)\,h$.
\end{proof}

\begin{exbox}
\begin{enumerate}
  \item Compute the Hessian of $f(x,y)=x^3-3xy^2$ and verify Clairaut's theorem.
  \item Show that $f(x,y)=xy\sin(1/(x^2+y^2))$ (extended by $f(0,0)=0$) has
        equal mixed partials at the origin despite neither being continuous there.
  \item Use the second-order expansion to classify the critical points of
        $f(x,y)=x^4+y^4-4xy$.
  \item Prove: if $H(a)$ is positive definite, then $a$ is a strict local minimum.
\end{enumerate}
\end{exbox}
\color{black}

%%-----------------------------------------------------------------------------
\chapter{Multivariable Optimization}
%%-----------------------------------------------------------------------------
\color{black}

\section{Critical Points and the Second Derivative Test}

The multivariable second-order analysis mirrors the one-dimensional case, with the second derivative replaced by the Hessian matrix.  A critical point is where the gradient vanishes.  Classification depends on the definiteness of the Hessian: a positive definite Hessian means the function curves upward in every direction (local minimum), while an indefinite Hessian means the function goes up in some directions and down in others (saddle point).  The eigenvalues of the Hessian determine definiteness --- a bridge to the linear algebra underlying optimisation.

\begin{defbox}{Critical Points in $\R^n$}
A point $a\in\R^n$ is a \textbf{critical point} of $f:\R^n\to\R$ if
$\nabla f(a)=\mathbf{0}$, i.e.\ all partial derivatives vanish at $a$.
\end{defbox}
\color{black}

\begin{thmbox}{Second Derivative Test in $\R^n$}
Let $a$ be a critical point of $f$ and $H=H(a)$ the Hessian at $a$.
\begin{itemize}[noitemsep]
  \item If $H$ is positive definite ($h^T Hh>0$ for all $h\neq 0$): strict local minimum.
  \item If $H$ is negative definite: strict local maximum.
  \item If $H$ is indefinite (positive and negative eigenvalues): saddle point.
  \item If $H$ is semidefinite: inconclusive.
\end{itemize}
\end{thmbox}
\color{black}

For $n=2$ the test reduces to: let $D=f_{xx}f_{yy}-f_{xy}^2$.
If $D>0$ and $f_{xx}>0$: local min; $D>0$ and $f_{xx}<0$: local max; $D<0$: saddle.

\section{Lagrange Multipliers}

Lagrange multipliers handle optimisation under constraints.  The geometric insight is clean: at a constrained extremum, the gradient of the objective must be parallel to the gradient of the constraint (otherwise we could move along the constraint surface and improve the objective).  The scalar $\lambda$ that relates the two gradients is the Lagrange multiplier; its value measures how much the optimal value would change if we relaxed the constraint slightly.  This interpretation makes Lagrange multipliers central to economics (shadow prices) and physics (action principles).

\begin{thmbox}{Lagrange Multiplier Theorem}
Suppose $f,g:\R^n\to\R$ are continuously differentiable and $a$ is a local
extremum of $f$ subject to the constraint $g(x)=c$, with $\nabla g(a)\neq\mathbf{0}$.
Then there exists $\lambda\in\R$ (the \textbf{Lagrange multiplier}) such that
$\nabla f(a)=\lambda\,\nabla g(a)$.
\end{thmbox}
\color{black}

\begin{proof}
Near $a$, the constraint surface $g(x)=c$ is locally a smooth $(n-1)$-dimensional
manifold.  Any curve $\mathbf{r}(t)$ on this surface through $a$ satisfies
$\nabla g(a)\cdot\mathbf{r}'(0)=0$.  At an extremum, $\nabla f(a)\cdot\mathbf{r}'(0)=0$
for every such curve.  Hence $\nabla f(a)$ is orthogonal to every vector in
$T_a(\{g=c\})$, so it is parallel to $\nabla g(a)$.
\end{proof}

\begin{example}
Maximise $f(x,y)=xy$ on the ellipse $x^2/4+y^2=1$.
Setting $\nabla f=\lambda\nabla g$: $(y,x)=\lambda(x/2,2y)$.
From $y=\lambda x/2$ and $x=2\lambda y$: $x=2\lambda\cdot\lambda x/2=\lambda^2 x$,
so $\lambda=\pm 1$.  The maximum is $f=1/2$ at $(\sqrt{2},1/\sqrt{2})$.
\end{example}

\begin{exbox}
\begin{enumerate}
  \item Find all critical points and classify them for
        $f(x,y)=x^3+y^3-3x-3y$.
  \item Use Lagrange multipliers to find the point on the plane $x+2y+3z=6$
        closest to the origin.
  \item Prove that the maximum of $\sum x_i^2$ subject to $\sum x_i=1$ and
        $x_i\geq 0$ is $1$, attained at the vertices.
  \item Show that the AM-GM inequality $\sqrt[n]{x_1\cdots x_n}\leq(x_1+\cdots+x_n)/n$
        follows from maximising $\ln x_1+\cdots+\ln x_n$ on $\sum x_i=n$.
\end{enumerate}
\end{exbox}
\color{black}

%%-----------------------------------------------------------------------------
\chapter{Multiple Integrals}
%%-----------------------------------------------------------------------------
\color{black}

\section{Double Integrals and Iterated Integration}

Integrating a function of two variables over a region $D$ means summing up infinitely many infinitely thin volume slices.  Fubini's theorem --- the central result of this section --- says we can compute this double integral by iterating two single-variable integrals, integrating in one variable at a time.  The order does not matter for continuous functions.  This is a deep theorem: it is not obvious that integrating $y$ first then $x$ gives the same answer as $x$ first then $y$, and the theorem can fail for non-integrable functions.

\begin{defbox}{Double Integral over a Rectangle}
For $f:[a,b]\times[c,d]\to\R$ bounded, the \textbf{double integral}
$\iint_R f\,dA$ is defined as the limit of double Riemann sums
$\sum_{i,j}f(x_i^*,y_j^*)\Delta x_i\Delta y_j$ as the mesh of both
partitions tends to zero.
\end{defbox}
\color{black}

\begin{thmbox}{Fubini's Theorem}
If $f:[a,b]\times[c,d]\to\R$ is continuous, then:
\[\iint_R f\,dA = \int_a^b\!\!\int_c^d f(x,y)\,dy\,dx
                = \int_c^d\!\!\int_a^b f(x,y)\,dx\,dy.\]
Both iterated integrals equal the double integral; in particular they are equal
to each other.
\end{thmbox}
\color{black}

\begin{proof}
Define $F(x)=\int_c^d f(x,y)\,dy$.  By uniform continuity of $f$ on the compact
rectangle, $F$ is continuous in $x$.  Show that $\int_a^b F(x)\,dx$ equals the
double Riemann sums in the limit using the one-variable integrability of $f(\cdot,y)$
for each $y$, then exchange the limits.
\end{proof}

\section{Change of Variables in $\R^n$}

The change of variables theorem is the multidimensional counterpart of $u$-substitution.  When we transform coordinates via $\mathbf{x} = \mathbf{G}(\mathbf{u})$, the volume element transforms by the absolute Jacobian determinant $|\det D\mathbf{G}|$.  This determinant measures how much the transformation stretches or compresses volumes locally.  Polar coordinates are the canonical example: the factor $r$ in $dA = r\,dr\,d\theta$ accounts for the fact that rings at larger radius are wider for the same angular increment.

\begin{thmbox}{Change of Variables (General)}
Let $\mathbf{G}:U\to V$ be a $C^1$ bijection between open sets in $\R^n$,
with Jacobian $J_\mathbf{G}(\mathbf{u})=\det D\mathbf{G}(\mathbf{u})\neq 0$.
Then for any integrable $f:V\to\R$:
\[\int_V f(\mathbf{x})\,d\mathbf{x}
  = \int_U f(\mathbf{G}(\mathbf{u}))\,|J_\mathbf{G}(\mathbf{u})|\,d\mathbf{u}.\]
\end{thmbox}
\color{black}

\begin{proof}[Proof sketch]
Partition $U$ into small boxes.  $\mathbf{G}$ maps each box approximately to a
parallelotope whose volume is $|\det D\mathbf{G}|\cdot\text{(box volume)}$ by the
linear approximation.  Summing and passing to the limit yields the formula.
\end{proof}

The three classical coordinate systems are instances:

\begin{itemize}[noitemsep]
  \item \textbf{Polar} $(r,\theta)$: $x=r\cos\theta$, $y=r\sin\theta$, Jacobian $=r$.
  \item \textbf{Cylindrical} $(r,\theta,z)$: Jacobian $=r$.
  \item \textbf{Spherical} $(\rho,\phi,\theta)$: $x=\rho\sin\phi\cos\theta$,
        $y=\rho\sin\phi\sin\theta$, $z=\rho\cos\phi$, Jacobian $=\rho^2\sin\phi$.
\end{itemize}

\begin{example}
$\int\!\!\int_{x^2+y^2\leq 1} e^{-(x^2+y^2)}\,dA$.
In polar coordinates: $\int_0^{2\pi}\!\!\int_0^1 e^{-r^2}\,r\,dr\,d\theta
=2\pi\cdot\tfrac{1}{2}(1-e^{-1})=\pi(1-e^{-1})$.
\end{example}

\begin{exbox}
\begin{enumerate}
  \item Use Fubini's theorem to evaluate
        $\int_0^1\!\!\int_0^1 (x+y)^2\,dx\,dy$.
  \item Compute the volume of the unit ball $B^3=\{(x,y,z):x^2+y^2+z^2\leq 1\}$
        using spherical coordinates.
  \item Prove the Gaussian integral $\int_{-\infty}^{\infty}e^{-x^2}\,dx=\sqrt{\pi}$
        by computing $\left(\int_{-\infty}^\infty e^{-x^2}\,dx\right)^2$ in polar
        coordinates.
  \item Show that Fubini's theorem can fail if $f$ is not integrable (Tonelli's
        counterexample: $f(x,y)=(x^2-y^2)/(x^2+y^2)^2$ on $[0,1]^2$).
\end{enumerate}
\end{exbox}
\color{black}

%%-----------------------------------------------------------------------------
\chapter{Vector Fields and Line Integrals}
%%-----------------------------------------------------------------------------
\color{black}

\begin{defbox}{Vector Fields and Line Integrals}
A \textbf{vector field} on $U\subseteq\R^n$ is a function
$\mathbf{F}:U\to\R^n$.

For a smooth curve $\mathbf{r}:[a,b]\to\R^n$, the \textbf{line integral} of
$\mathbf{F}$ along $\mathbf{r}$ is:
\[\int_C \mathbf{F}\cdot d\mathbf{r}
  := \int_a^b \mathbf{F}(\mathbf{r}(t))\cdot\mathbf{r}'(t)\,dt.\]
This measures the total work done by $\mathbf{F}$ along the path.
\end{defbox}
\color{black}

\begin{defbox}{Conservative Fields and Potential Functions}
$\mathbf{F}$ is \textbf{conservative} if $\mathbf{F}=\nabla\phi$ for some scalar
function $\phi$ (called the \textbf{potential}).
\end{defbox}
\color{black}

\begin{thmbox}{Fundamental Theorem for Line Integrals}
If $\mathbf{F}=\nabla\phi$ on $U$, then for any curve $C$ from $A$ to $B$ in $U$:
\[\int_C \mathbf{F}\cdot d\mathbf{r} = \phi(B)-\phi(A).\]
In particular, line integrals of conservative fields are path-independent.
\end{thmbox}
\color{black}

\begin{proof}
$\int_a^b \nabla\phi(\mathbf{r}(t))\cdot\mathbf{r}'(t)\,dt
 = \int_a^b \frac{d}{dt}\phi(\mathbf{r}(t))\,dt
 = \phi(\mathbf{r}(b))-\phi(\mathbf{r}(a))$ by the chain rule and FTC.
\end{proof}

\begin{thmbox}{Characterisation of Conservative Fields}
On a simply connected domain, $\mathbf{F}$ is conservative iff
$\nabla\times\mathbf{F}=\mathbf{0}$ (the curl vanishes).
\end{thmbox}
\color{black}

\begin{exbox}
\begin{enumerate}
  \item Compute $\int_C \mathbf{F}\cdot d\mathbf{r}$ for $\mathbf{F}=(y,x)$ along
        the parabola $y=x^2$ from $(0,0)$ to $(1,1)$.
  \item Show $\mathbf{F}=(2xy+z^2,x^2,2xz)$ is conservative and find a potential.
  \item Prove: on a simply connected domain, $\oint_C \mathbf{F}\cdot d\mathbf{r}=0$
        for every closed curve iff $\mathbf{F}$ is conservative.
  \item Compute the work done by gravity $\mathbf{F}=(0,0,-mg)$ on a particle
        moving from $(0,0,h_1)$ to $(0,0,h_2)$ along any path.
\end{enumerate}
\end{exbox}
\color{black}

%%-----------------------------------------------------------------------------
\chapter{Surface Integrals}
%%-----------------------------------------------------------------------------
\color{black}

\begin{defbox}{Parametric Surfaces and Surface Area}
A \textbf{parametric surface} is a map $\mathbf{r}:D\subseteq\R^2\to\R^3$.
The \textbf{surface area element} is $d\mathbf{S}=(\mathbf{r}_u\times\mathbf{r}_v)\,du\,dv$.
The \textbf{surface area} is $\iint_D|\mathbf{r}_u\times\mathbf{r}_v|\,du\,dv$.

The \textbf{flux integral} (surface integral of a vector field) is:
\[\iint_S \mathbf{F}\cdot d\mathbf{S}
  = \iint_D \mathbf{F}(\mathbf{r}(u,v))\cdot(\mathbf{r}_u\times\mathbf{r}_v)\,du\,dv.\]
\end{defbox}
\color{black}

\begin{example}
The flux of $\mathbf{F}=(0,0,z)$ upward through the paraboloid $z=1-x^2-y^2$
over $D=\{x^2+y^2\leq 1\}$.  In polar: $\mathbf{r}=(r\cos\theta,r\sin\theta,1-r^2)$,
$\mathbf{r}_u\times\mathbf{r}_v=(2r^2\cos\theta,2r^2\sin\theta,r)$,
$\mathbf{F}\cdot d\mathbf{S}=(1-r^2)r\,dr\,d\theta$.
Integrating: $\int_0^{2\pi}\!\!\int_0^1 (r-r^3)\,dr\,d\theta
=2\pi(1/2-1/4)=\pi/2$.
\end{example}

\begin{exbox}
\begin{enumerate}
  \item Compute the surface area of the sphere $x^2+y^2+z^2=R^2$.
  \item Find the flux of $\mathbf{F}=(x,y,z)$ outward through the unit sphere.
  \item Show the surface area element for a graph $z=f(x,y)$ is
        $\sqrt{1+f_x^2+f_y^2}\,dx\,dy$.
\end{enumerate}
\end{exbox}
\color{black}

%%-----------------------------------------------------------------------------
\chapter{The Three Great Theorems}
%%-----------------------------------------------------------------------------
\color{black}

\begin{remarkbox}
\textbf{Historical Note: Green, Stokes, Gauss, and the Theorems Named After Them}

The three theorems of this chapter have a tangled history of credit and misattribution.

\textbf{Green's theorem} was published by George Green in 1828 in a privately printed pamphlet, \emph{An Essay on the Application of Mathematical Analysis to the Theories of Electricity and Magnetism}, sold by subscription in Nottingham. Green was a self-taught miller's son with almost no formal education. He distributed roughly fifty copies; the pamphlet was essentially unknown until William Thomson (later Lord Kelvin) rediscovered it in 1845 and recognised its importance. By then Green had died (1841), having finally been admitted to Cambridge in 1833 at age 40. Gauss had proved the same result independently for three dimensions (the Divergence Theorem) around 1813, but his proof was also not widely known.

\textbf{Stokes' theorem} appears, in modern form, in a letter by Kelvin to Stokes in 1850. Stokes set it as a problem on the Smith's Prize examination at Cambridge in 1854 --- an examination sat by Maxwell, among others. The theorem thus bears Stokes' name partly by examination accident. Earlier versions were known to Ampère and to Hankel, and the full proof in the modern sense was given by Brioschi in 1854.

The \textbf{Divergence Theorem} is variously attributed to Gauss, Green, Ostrogradsky (who published a proof in 1831 and is given credit in Russian textbooks), and Stokes. The pattern is typical of nineteenth-century mathematics: ideas were discovered repeatedly and independently, publications were slow and distribution limited, and national pride played a large role in determining who got credit in which country's textbooks.

The \textbf{unification} under the generalised Stokes' theorem on differential forms is due to \textbf{Cartan} (1899--1945), who developed the modern language of differential forms and exterior derivatives in the early twentieth century. Cartan's formulation shows that Green's, Stokes', the Divergence theorem, and the Fundamental Theorem of Calculus are all the same theorem $\int_M d\omega = \int_{\partial M} \omega$, differing only in dimension and the choice of differential form. This unification was one of the great conceptual achievements of modern mathematics.
\end{remarkbox}
\color{black}

The three great theorems of vector calculus --- Green's, Stokes', and the Divergence
theorem --- all say the same thing at different levels of generality: \emph{an
integral over a region equals an integral over its boundary}.

\section{Green's Theorem}

Green's theorem relates a line integral around a closed curve to a double integral over the enclosed region.  It is the two-dimensional Fundamental Theorem of Calculus: the integral of a ``derivative-like'' combination $(Q_x - P_y)$ over a region equals the integral of $(P, Q)$ around the boundary.  The theorem has an immediate payoff: it lets us compute areas using line integrals and circulation using area integrals, switching representations as convenient.

\begin{thmbox}{Green's Theorem}
Let $C$ be a positively oriented, piecewise smooth, simple closed curve bounding
a region $D$ in $\R^2$.  For continuously differentiable $P,Q$:
\[\oint_C (P\,dx + Q\,dy)
  = \iint_D \left(\frac{\partial Q}{\partial x} - \frac{\partial P}{\partial y}\right)dA.\]
\end{thmbox}
\color{black}

\begin{proof}
It suffices to prove $\oint_C P\,dx = -\iint_D P_y\,dA$ and
$\oint_C Q\,dy = \iint_D Q_x\,dA$ separately.
For a region $D=\{(x,y):a\leq x\leq b,\,g_1(x)\leq y\leq g_2(x)\}$:
\[\iint_D P_y\,dA = \int_a^b [P(x,g_2(x))-P(x,g_1(x))]\,dx = -\oint_C P\,dx,\]
where the last step parametrises the boundary in two pieces and applies FTC.
\end{proof}

\begin{corbox}{Area via Line Integral}
The area of region $D$ with boundary $C$ is
$A=\frac{1}{2}\oint_C (x\,dy - y\,dx)$.
\end{corbox}
\color{black}

\section{Stokes' Theorem}

Stokes' theorem extends Green's theorem from flat planar regions to curved surfaces in $\R^3$.  The curl $\nabla \times \mathbf{F}$ measures the local rotation of a vector field; Stokes says the total rotation through a surface equals the circulation around its boundary curve.  This is the mathematical core of Faraday's law in electromagnetism and of the Kelvin--Stokes theorem in fluid mechanics.

\begin{thmbox}{Stokes' Theorem}
Let $S$ be an oriented surface in $\R^3$ with positively oriented boundary curve
$\partial S$.  For a $C^1$ vector field $\mathbf{F}$:
\[\iint_S (\nabla\times\mathbf{F})\cdot d\mathbf{S}
  = \oint_{\partial S} \mathbf{F}\cdot d\mathbf{r}.\]
\end{thmbox}
\color{black}

\begin{proof}
Write $\mathbf{F}=(P,Q,R)$ and $S$ as a parametric surface.
Expand both sides in coordinates; after a careful calculation using Green's
theorem on each coordinate patch, the two sides match.
\end{proof}

Green's theorem is the special case where $S$ is a planar region.

\section{The Divergence Theorem}

The Divergence theorem relates the outward flux through a closed surface to the total source strength (divergence) inside the enclosed volume.  If $\mathbf{F}$ is a fluid velocity field, it says: the net flow out through the surface equals the total rate of fluid production inside.  This is the mathematical statement of conservation laws, and it underlies Gauss's law in electrostatics as well as the continuity equations of fluid dynamics.

\begin{thmbox}{Divergence Theorem (Gauss's Theorem)}
Let $E$ be a solid region in $\R^3$ with smooth closed boundary $\partial E$
(oriented outward).  For a $C^1$ vector field $\mathbf{F}$:
\[\iint_{\partial E} \mathbf{F}\cdot d\mathbf{S}
  = \iiint_E \nabla\cdot\mathbf{F}\,dV.\]
\end{thmbox}
\color{black}

\begin{proof}
Prove the component identity $\iint_{\partial E} R\,dA_z = \iiint_E R_z\,dV$
(where $dA_z$ is the $z$-component of $d\mathbf{S}$) by writing $E$ as a
region between two graphs and applying FTC.  The full theorem follows by adding
all three components.
\end{proof}

\section{Unification: Generalised Stokes' Theorem}

All three theorems are special cases of one statement on differential forms:

\begin{thmbox}{Generalised Stokes' Theorem}
Let $\omega$ be a smooth $(n-1)$-form on an oriented $n$-dimensional manifold $M$
with boundary $\partial M$.  Then:
\[\int_M d\omega = \int_{\partial M} \omega.\]
\end{thmbox}
\color{black}

Here $d$ is the exterior derivative.  The dictionary is:
\begin{itemize}[noitemsep]
  \item Green's theorem: $M$ is a 2D region, $\omega=P\,dx+Q\,dy$,
        $d\omega=(Q_x-P_y)\,dx\wedge dy$.
  \item Stokes' theorem: $M$ is a surface in $\R^3$, $\omega=\mathbf{F}\cdot d\mathbf{r}$,
        $d\omega=(\nabla\times\mathbf{F})\cdot d\mathbf{S}$.
  \item Divergence theorem: $M$ is a 3D region, $\omega=\mathbf{F}\cdot d\mathbf{S}$,
        $d\omega=\nabla\cdot\mathbf{F}\,dV$.
  \item FTC: $M=[a,b]$, $\omega=f$, $d\omega=f'\,dx$,
        $\partial M=\{b\}-\{a\}$.
\end{itemize}

\begin{exbox}
\begin{enumerate}
  \item Use Green's theorem to compute the area enclosed by the ellipse $x^2/a^2+y^2/b^2=1$.
  \item Use Stokes' theorem to evaluate
        $\oint_C (y^2\,dx+z\,dy+x\,dz)$ where $C$ is the unit circle in the plane $z=0$.
  \item Use the Divergence theorem to compute the flux of
        $\mathbf{F}=(x^3,y^3,z^3)$ outward through the unit sphere.
  \item Show that $\nabla\cdot(\nabla\times\mathbf{F})=0$ and
        $\nabla\times(\nabla f)=\mathbf{0}$ algebraically, then interpret them via
        the generalised Stokes' theorem.
  \item Derive the heat equation $u_t=\nabla^2 u$ from conservation of energy and
        the divergence theorem.
\end{enumerate}
\end{exbox}
\color{black}

%%=============================================================================
\part{The Capstone: \textit{Symboliculus}}
%%=============================================================================
\color{black}

The previous six parts built calculus from scratch and proved its core theorems.
Now we build something with those theorems: \textbf{Symboliculus}, a command-line
symbolic calculus solver in Lean~4 that computes derivatives and integrals with
full step-by-step traces --- just like Symbolab, but with every rule
traced back to the mathematics we proved.

%%-----------------------------------------------------------------------------
\chapter{Expression Trees}
%%-----------------------------------------------------------------------------
\color{black}

\section{The \li{Expr} Inductive Type}

A symbolic expression is a tree.  Each node is either a leaf (a constant or
variable) or an operator applied to one or two sub-expressions.

\begin{leanbox}
\begin{lstlisting}
-- symboliculus/Expr.lean
namespace Symboliculus

inductive Expr where
  | Const  : Float -> Expr               -- numeric constant
  | Var    : String -> Expr              -- variable (e.g. "x")
  | Add    : Expr -> Expr -> Expr
  | Sub    : Expr -> Expr -> Expr
  | Mul    : Expr -> Expr -> Expr
  | Div    : Expr -> Expr -> Expr
  | Pow    : Expr -> Expr -> Expr        -- e1 ^ e2
  | Neg    : Expr -> Expr
  | Sin    : Expr -> Expr
  | Cos    : Expr -> Expr
  | Tan    : Expr -> Expr
  | Exp    : Expr -> Expr                -- e^(...)
  | Ln     : Expr -> Expr
  | Abs    : Expr -> Expr
  deriving Repr, BEq

end Symboliculus
\end{lstlisting}
\end{leanbox}
\color{black}

\section{Evaluation}

Once we have the syntax tree, evaluation gives it meaning: given an assignment of numerical values to variables, \li{eval} traverses the tree and computes the corresponding \li{Float}.  This is the \textbf{denotational semantics} of our expression language --- exactly the semantics function from the lambda calculus book, now applied to a concrete arithmetic language.  Note that \li{eval} is total: every well-formed \li{Expr} has a value (though it may be \li{Float.nan} or \li{Float.inf} for divisions by zero).

\begin{leanbox}
\begin{lstlisting}
-- Evaluate an expression given a variable assignment
def Expr.eval (env : String -> Float) : Expr -> Float
  | Const c        => c
  | Var   x        => env x
  | Add   e1 e2    => e1.eval env + e2.eval env
  | Sub   e1 e2    => e1.eval env - e2.eval env
  | Mul   e1 e2    => e1.eval env * e2.eval env
  | Div   e1 e2    => e1.eval env / e2.eval env
  | Pow   e1 e2    => Float.pow (e1.eval env) (e2.eval env)
  | Neg   e        => -(e.eval env)
  | Sin   e        => Float.sin (e.eval env)
  | Cos   e        => Float.cos (e.eval env)
  | Tan   e        => Float.tan (e.eval env)
  | Exp   e        => Float.exp (e.eval env)
  | Ln    e        => Float.log (e.eval env)
  | Abs   e        => Float.abs (e.eval env)
\end{lstlisting}
\end{leanbox}
\color{black}

\section{Pretty-Printing}

Pretty-printing is the inverse of parsing: it turns an \li{Expr} tree back into a human-readable string.  The challenge is parenthesisation: we want $x + y \times z$ not $(x + y) \times z$, but we need parentheses where the tree structure would otherwise be ambiguous.  The precedence function assigns a priority to each operator, and the \li{paren} helper wraps a subexpression in parentheses only when its operator binds less tightly than the surrounding context requires.

\begin{leanbox}
\begin{lstlisting}
-- Pretty-print with minimal parentheses
def Expr.precedence : Expr -> Nat
  | Add _ _ | Sub _ _ => 1
  | Mul _ _ | Div _ _ => 2
  | Neg _              => 3
  | Pow _ _            => 4
  | _                  => 10

def Expr.pp : Expr -> String
  | Const c     => if c == Float.floor c then
                     toString (Int.ofFloat c)
                   else toString c
  | Var x       => x
  | Add e1 e2   => s!"{e1.pp} + {e2.pp}"
  | Sub e1 e2   => s!"{e1.pp} - {paren e2 2}"
  | Mul e1 e2   => s!"{paren e1 2} * {paren e2 3}"
  | Div e1 e2   => s!"{paren e1 2} / {paren e2 3}"
  | Pow e1 e2   => s!"{paren e1 5}^{paren e2 5}"
  | Neg e       => s!"-{paren e 4}"
  | Sin e       => s!"sin({e.pp})"
  | Cos e       => s!"cos({e.pp})"
  | Tan e       => s!"tan({e.pp})"
  | Exp e       => s!"exp({e.pp})"
  | Ln  e       => s!"ln({e.pp})"
  | Abs e       => s!"|{e.pp}|"
where
  paren (e : Expr) (minPrec : Nat) :=
    if e.precedence < minPrec then s!"({e.pp})" else e.pp
\end{lstlisting}
\end{leanbox}
\color{black}

\section{Parsing}

The parser turns a string like \li{"x\^{}2 * sin(x)"} into an \li{Expr}.
We implement a simple recursive-descent parser with operator precedence.

\begin{leanbox}
\begin{lstlisting}
-- Tokeniser
inductive Token where
  | TNum  : Float  -> Token
  | TVar  : String -> Token
  | TPlus | TMinus | TStar | TSlash | TCaret
  | TLParen | TRParen | TEOF
  deriving Repr, BEq

def tokenise (s : String) : List Token := sorry

-- Recursive-descent parser (Pratt / precedence climbing)
structure Parser where
  tokens : Array Token
  pos    : Nat

def Parser.peek (p : Parser) : Token :=
  p.tokens.getD p.pos Token.TEOF

def Parser.advance (p : Parser) : Parser :=
  { p with pos := p.pos + 1 }

-- parseExpr, parseTerm, parseFactor defined by mutual recursion
-- (full implementation in the Lake project)
def parseExpr  : Parser -> Except String (Expr × Parser) := sorry
def parseTerm  : Parser -> Except String (Expr × Parser) := sorry
def parseFactor: Parser -> Except String (Expr × Parser) := sorry

def parseStr (s : String) : Except String Expr :=
  match parseExpr { tokens := (tokenise s).toArray, pos := 0 } with
  | .ok (e, _) => .ok e
  | .error msg => .error msg
\end{lstlisting}
\end{leanbox}
\color{black}

\section{Algebraic Simplification}

Before and after differentiation or integration we run a simplification pass
to collapse trivial expressions.

\begin{leanbox}
\begin{lstlisting}
def Expr.simplify : Expr -> Expr
  -- Arithmetic on constants
  | Add (Const a) (Const b)  => Const (a + b)
  | Sub (Const a) (Const b)  => Const (a - b)
  | Mul (Const a) (Const b)  => Const (a * b)
  | Div (Const a) (Const b)  => if b != 0 then Const (a / b) else Div (Const a) (Const b)
  -- Additive identity
  | Add e (Const 0)          => e.simplify
  | Add (Const 0) e          => e.simplify
  -- Multiplicative identity / zero
  | Mul e (Const 1)          => e.simplify
  | Mul (Const 1) e          => e.simplify
  | Mul _ (Const 0)          => Const 0
  | Mul (Const 0) _          => Const 0
  -- Power rules
  | Pow e (Const 1)          => e.simplify
  | Pow _ (Const 0)          => Const 1
  | Pow (Const 1) _          => Const 1
  -- Double negation
  | Neg (Neg e)              => e.simplify
  -- Recurse
  | Add e1 e2                => Add e1.simplify e2.simplify
  | Sub e1 e2                => Sub e1.simplify e2.simplify
  | Mul e1 e2                => Mul e1.simplify e2.simplify
  | Div e1 e2                => Div e1.simplify e2.simplify
  | Pow e1 e2                => Pow e1.simplify e2.simplify
  | Neg e                    => Neg e.simplify
  | Sin e                    => Sin e.simplify
  | Cos e                    => Cos e.simplify
  | Tan e                    => Tan e.simplify
  | Exp e                    => Exp e.simplify
  | Ln  e                    => Ln  e.simplify
  | Abs e                    => Abs e.simplify
  | e                        => e

-- Repeat until fixed point
def Expr.fullSimplify (e : Expr) (fuel : Nat := 20) : Expr :=
  match fuel with
  | 0     => e
  | n + 1 =>
    let e' := e.simplify
    if e' == e then e else e'.fullSimplify n
\end{lstlisting}
\end{leanbox}
\color{black}

\begin{exbox}
\begin{enumerate}
  \item Implement structural equality for \li{Expr} (two expressions are equal iff
        they have the same tree shape and constants).
  \item Add \li{Arcsin}, \li{Arccos}, \li{Arctan} constructors and extend
        \li{eval} and \li{pp} accordingly.
  \item Extend \li{fullSimplify} to handle $e^{\ln x}=x$ and $\ln(e^x)=x$.
  \item Prove in Lean~4: for any \li{e : Expr} and environment \li{env},
        \li{e.simplify.eval env = e.eval env}.
\end{enumerate}
\end{exbox}
\color{black}

%%-----------------------------------------------------------------------------
\chapter{The Differentiation Engine}
%%-----------------------------------------------------------------------------
\color{black}

\section{The \li{DiffStep} Trace Type}

Each call to the differentiation engine returns not just the result but a
\emph{proof trace} recording which rule was applied and to which subexpression.

\begin{leanbox}
\begin{lstlisting}
-- symboliculus/Diff.lean
inductive DiffStep where
  | ConstRule   : DiffStep                          -- d/dx c = 0
  | VarRule     : DiffStep                          -- d/dx x = 1
  | VarOther    : DiffStep                          -- d/dx y = 0 (y != x)
  | SumRule     : DiffStep -> DiffStep -> DiffStep  -- (u+v)' = u'+v'
  | SubRule     : DiffStep -> DiffStep -> DiffStep
  | ProductRule : DiffStep -> DiffStep -> DiffStep  -- (uv)' = u'v+uv'
  | QuotRule    : DiffStep -> DiffStep -> DiffStep  -- (u/v)' = (u'v-uv')/v^2
  | ChainRule   : String   -> DiffStep -> DiffStep  -- outer fn, inner deriv
  | PowerRule   : DiffStep                          -- d/dx x^n = nx^{n-1}
  | NegRule     : DiffStep -> DiffStep
  deriving Repr
\end{lstlisting}
\end{leanbox}
\color{black}

\section{The \li{diff} Function}

The differentiation engine is a structural recursion over \li{Expr}.  Each constructor of \li{Expr} corresponds to exactly one differentiation rule from Part~III, and the Lean pattern match makes this correspondence explicit.  Crucially, each recursive call produces not just the derivative expression but also a \li{DiffStep} that records which rule was applied --- this is what will become the step-by-step trace the user sees.  After computing the derivative, we run \li{fullSimplify} to collapse trivial expressions before returning.

\begin{leanbox}
\begin{lstlisting}
-- diff e x returns (derivative, trace)
def diff (x : String) : Expr -> Expr × DiffStep
  -- Constants
  | Const _     => (Const 0, .ConstRule)

  -- Variables
  | Var y       => if y == x
                   then (Const 1, .VarRule)
                   else (Const 0, .VarOther)

  -- Negation: (-u)' = -u'
  | Neg u       =>
      let (u', s) := diff x u
      (Neg u', .NegRule s)

  -- Sum rule: (u + v)' = u' + v'
  | Add u v     =>
      let (u', su) := diff x u
      let (v', sv) := diff x v
      (Add u' v' |>.fullSimplify, .SumRule su sv)

  -- Sub rule: (u - v)' = u' - v'
  | Sub u v     =>
      let (u', su) := diff x u
      let (v', sv) := diff x v
      (Sub u' v' |>.fullSimplify, .SubRule su sv)

  -- Product rule: (uv)' = u'v + uv'
  | Mul u v     =>
      let (u', su) := diff x u
      let (v', sv) := diff x v
      ((Add (Mul u' v) (Mul u v')).fullSimplify, .ProductRule su sv)

  -- Quotient rule: (u/v)' = (u'v - uv') / v^2
  | Div u v     =>
      let (u', su) := diff x u
      let (v', sv) := diff x v
      ((Div (Sub (Mul u' v) (Mul u v')) (Pow v (Const 2))).fullSimplify,
       .QuotRule su sv)

  -- Power rule: (u^n)' = n * u^(n-1) * u'  (general, via chain rule)
  | Pow u (Const n) =>
      let (u', su) := diff x u
      ((Mul (Mul (Const n) (Pow u (Const (n-1)))) u').fullSimplify,
       .PowerRule)

  -- Trig: sin(u)' = cos(u) * u'
  | Sin u       =>
      let (u', su) := diff x u
      ((Mul (Cos u) u').fullSimplify, .ChainRule "sin" su)

  -- Trig: cos(u)' = -sin(u) * u'
  | Cos u       =>
      let (u', su) := diff x u
      ((Neg (Mul (Sin u) u')).fullSimplify, .ChainRule "cos" su)

  -- Tan: (tan u)' = sec^2(u) * u' = u' / cos^2(u)
  | Tan u       =>
      let (u', su) := diff x u
      ((Div u' (Pow (Cos u) (Const 2))).fullSimplify, .ChainRule "tan" su)

  -- Exp: (e^u)' = e^u * u'
  | Exp u       =>
      let (u', su) := diff x u
      ((Mul (Exp u) u').fullSimplify, .ChainRule "exp" su)

  -- Ln: (ln u)' = u' / u
  | Ln u        =>
      let (u', su) := diff x u
      ((Div u' u).fullSimplify, .ChainRule "ln" su)

  -- General Pow: (u^v)' = u^v * (v'*ln(u) + v*u'/u)
  | Pow u v     =>
      let (u', su) := diff x u
      let (v', sv) := diff x v
      let result := Mul (Pow u v)
                        (Add (Mul v' (Ln u)) (Div (Mul v u') u))
      (result.fullSimplify, .ChainRule "pow" (.ProductRule su sv))

  | Abs _       => (Const 0, .ConstRule)  -- simplified; full version uses sign
\end{lstlisting}
\end{leanbox}
\color{black}

\section{Rendering the Trace}

The \li{DiffStep} tree records the \emph{structure} of the derivation.  Rendering turns that structure into text the user can read.  The renderer is also a structural recursion: each \li{DiffStep} constructor maps to a human-readable description of the rule used, optionally recursing into sub-derivations.  This is the piece that makes Symboliculus feel like Symbolab --- without the trace renderer, the user would just see the final answer with no explanation of how it was reached.

\begin{leanbox}
\begin{lstlisting}
def DiffStep.render (orig result : Expr) : DiffStep -> String
  | .ConstRule    => s!"d/dx({orig.pp}) = 0  [constant rule]"
  | .VarRule      => s!"d/dx({orig.pp}) = 1  [variable rule]"
  | .VarOther     => s!"d/dx({orig.pp}) = 0  [independent variable]"
  | .NegRule s    =>
      s!"Negate: d/dx(-u) = -(d/dx u)\n  " ++ s.render orig result
  | .SumRule s t  =>
      s!"Sum rule: d/dx(u + v) = d/dx(u) + d/dx(v)"
  | .ProductRule s t =>
      s!"Product rule: d/dx(u*v) = u'*v + u*v'"
  | .QuotRule s t =>
      s!"Quotient rule: d/dx(u/v) = (u'v - uv') / v^2"
  | .ChainRule fn s =>
      s!"Chain rule on {fn}: d/dx({fn}(u)) = {fn}'(u) * d/dx(u)"
  | .PowerRule    =>
      s!"Power rule: d/dx(u^n) = n * u^(n-1) * u'"
  | _             => s!"[step applied to {orig.pp}]"
\end{lstlisting}
\end{leanbox}
\color{black}

\begin{exbox}
\begin{enumerate}
  \item Trace through \li{diff "x" (Mul (Var "x") (Sin (Var "x")))} by hand,
        showing the \li{DiffStep} tree produced and the rendered output.
  \item Extend \li{diff} to handle \li{Arctan} with the rule
        $(\arctan u)' = u'/(1+u^2)$.
  \item Prove in Lean~4: for \li{diff "x" (Add u v)}, the result equals
        \li{Add (diff "x" u).1 (diff "x" v).1} (sum rule is structurally correct).
  \item Add a \li{DiffStep.renderFull} that recursively renders the entire trace
        tree as a numbered list of steps.
\end{enumerate}
\end{exbox}
\color{black}

%%-----------------------------------------------------------------------------
\chapter{The Integration Engine}
%%-----------------------------------------------------------------------------
\color{black}

Integration is harder than differentiation: no single algorithm handles all cases,
and some elementary functions have no elementary antiderivative (Risch's theorem,
Chapter~36).  We model the engine as a collection of \emph{strategies}, each a
function that either succeeds with a trace or fails, allowing the engine to try
the next approach.

\section{Strategies and the \li{IntStep} Trace}

Integration is fundamentally different from differentiation: there is no single algorithm that handles all cases.  Instead, we model the engine as a collection of \emph{strategies}, each specialised to a class of integrands.  A strategy is a function that either succeeds --- returning an antiderivative and a trace --- or fails, signalling that the next strategy should try.  This is the strategy pattern from software engineering, but it also has a mathematical interpretation: the strategy chain is a propagator network over the lattice \li{Option (Expr $\times$ IntStep)}, ordered by $\bot <$ \li{Some}.

\begin{leanbox}
\begin{lstlisting}
-- symboliculus/Integrate.lean

inductive IntStep where
  | DirectLookup  : String -> IntStep     -- "integral of sin(x) is -cos(x)"
  | ConstMul      : IntStep -> IntStep    -- pull out constant
  | SumRule       : IntStep -> IntStep -> IntStep
  | USub          : Expr -> Expr -> IntStep -> IntStep
                    -- u-substitution: u=..., du=...
  | IBP           : Expr -> Expr -> IntStep -> IntStep -> IntStep
                    -- integration by parts: u=..., dv=...
  | PartialFrac   : List (Expr × Expr) -> List IntStep -> IntStep
                    -- partial fractions: [(factor, coeff)]
  | TrigSub       : String -> Expr -> IntStep -> IntStep
                    -- trig substitution: x = a*sin(t) etc.
  | TrigReduce    : IntStep -> IntStep    -- trig power reduction
  | Failed        : IntStep              -- no strategy worked
  deriving Repr

-- A strategy is a function that maybe integrates e w.r.t. x
def Strategy := String -> Expr -> Option (Expr × IntStep)
\end{lstlisting}
\end{leanbox}
\color{black}

\section{Direct Lookup Table}

The first strategy is the simplest: pattern-match the integrand against a table of known antiderivatives.  This handles atomic expressions directly --- constants, the identity, $\sin x$, $\cos x$, $e^x$, and power functions $x^n$ for $n \neq -1$.  When this strategy succeeds, the trace is a \li{DirectLookup} leaf with a human-readable description of the formula used.  When it fails, the dispatcher moves on to more sophisticated strategies.

\begin{leanbox}
\begin{lstlisting}
-- Antiderivatives of atomic expressions
def directLookup (x : String) : Expr -> Option (Expr × IntStep)
  | Const c           =>
      some (Mul (Const c) (Var x), .DirectLookup s!"∫{c} dx = {c}x")
  | Var y             =>
      if y == x
      then some (Div (Pow (Var x) (Const 2)) (Const 2),
                 .DirectLookup "∫x dx = x²/2")
      else none
  | Sin (Var y)       =>
      if y == x
      then some (Neg (Cos (Var x)), .DirectLookup "∫sin(x) dx = -cos(x)")
      else none
  | Cos (Var y)       =>
      if y == x
      then some (Sin (Var x), .DirectLookup "∫cos(x) dx = sin(x)")
      else none
  | Exp (Var y)       =>
      if y == x
      then some (Exp (Var x), .DirectLookup "∫eˣ dx = eˣ")
      else none
  | Pow (Var y) (Const n) =>
      if y == x && n != -1
      then some (Div (Pow (Var x) (Const (n+1))) (Const (n+1)),
                 .DirectLookup s!"∫xⁿ dx = xⁿ⁺¹/(n+1)")
      else none
  | _                 => none
\end{lstlisting}
\end{leanbox}
\color{black}

\section{U-Substitution Detector}

The $u$-substitution strategy looks for an integrand of the form $f(g(x))\cdot g'(x)$.
It tries common choices for $u$ and checks whether differentiating $u$ (and dividing)
simplifies the integrand.

\begin{leanbox}
\begin{lstlisting}
-- Try u = inner function, check if integrand / du simplifies
def tryUSub (x : String) (e : Expr) : Option (Expr × IntStep) :=
  -- Collect candidate u-substitutions from the expression tree
  let candidates := innerFunctions e
  candidates.findSome? fun u =>
    let du := (diff x u).1          -- du/dx
    -- Check if e / (du/dx) depends only on u (not x directly)
    let reduced := (Div e du).fullSimplify
    if isInTermsOf u reduced then
      -- Integrate the reduced expression in terms of u
      match directLookup "u" (substVar x "u" reduced) with
      | some (antideriv, step) =>
          -- Back-substitute x for u
          some ((substExpr "u" u antideriv).fullSimplify,
                .USub u du step)
      | none => none
    else none

-- Collect the innermost composite functions as u candidates
def innerFunctions : Expr -> List Expr
  | Sin u | Cos u | Tan u | Exp u | Ln u => [u]
  | Pow u _ | Pow _ u                    => [u]
  | Mul u v | Add u v                    => innerFunctions u ++ innerFunctions v
  | _                                    => []
\end{lstlisting}
\end{leanbox}
\color{black}

\section{Integration by Parts Heuristic}

The LIATE rule (Logarithm, Inverse trig, Algebraic, Trig, Exponential) guides
which factor to assign to $u$ and which to $dv$.

\begin{leanbox}
\begin{lstlisting}
-- LIATE priority for choosing u
def liatePriority : Expr -> Nat
  | Ln _               => 0   -- highest: logarithm
  | Arctan _ | Arcsin _ => 1  -- inverse trig
  | Pow (Var _) _      => 2   -- algebraic (polynomial)
  | Sin _ | Cos _      => 3   -- trigonometric
  | Exp _              => 4   -- exponential (lowest priority)
  | _                  => 5

def tryIBP (x : String) (e : Expr) : Option (Expr × IntStep) :=
  match e with
  | Mul u v =>
    -- Assign u and dv by LIATE priority
    let (uPart, dvPart) :=
      if liatePriority u <= liatePriority v
      then (u, v)
      else (v, u)
    let uPrime := (diff x uPart).1
    -- Try to integrate dvPart directly
    match directLookup x dvPart with
    | some (vPart, _) =>
        -- IBP: ∫u dv = uv - ∫v du
        let remaining := (Mul vPart uPrime).fullSimplify
        match integrate x remaining with
        | some (intRemainder, step) =>
            let result := Sub (Mul uPart vPart) intRemainder
            some (result.fullSimplify,
                  .IBP uPart dvPart (.DirectLookup "") step)
        | none => none
    | none => none
  | _ => none
\end{lstlisting}
\end{leanbox}
\color{black}

\section{Top-Level \li{integrate} Dispatcher}

The dispatcher tries each strategy in a fixed priority order and returns the first success.  The order reflects the complexity hierarchy of integration techniques: we try the cheapest checks first (direct lookup, constant multiple, sum) before invoking the more expensive searches (substitution, parts, partial fractions).  This greedy first-success approach is sound but not complete: some integrands require trying multiple strategies in combination, which a more sophisticated planner could handle.

\begin{leanbox}
\begin{lstlisting}
-- Try strategies in order; return first success
def integrate (x : String) (e : Expr) : Option (Expr × IntStep) :=
  let strategies : List Strategy :=
    [ directLookup,
      constMulStrategy,
      sumStrategy,
      tryUSub,
      tryIBP,
      partialFracStrategy,
      trigReduceStrategy ]
  strategies.findSome? (· x e)

-- Sum rule strategy
def sumStrategy (x : String) : Expr -> Option (Expr × IntStep)
  | Add u v =>
      match integrate x u, integrate x v with
      | some (U, su), some (V, sv) =>
          some ((Add U V).fullSimplify, .SumRule su sv)
      | _, _ => none
  | _ => none

-- Constant multiple strategy
def constMulStrategy (x : String) : Expr -> Option (Expr × IntStep)
  | Mul (Const c) e =>
      match integrate x e with
      | some (E, s) => some ((Mul (Const c) E).fullSimplify, .ConstMul s)
      | none        => none
  | Mul e (Const c) =>
      match integrate x e with
      | some (E, s) => some ((Mul (Const c) E).fullSimplify, .ConstMul s)
      | none        => none
  | _ => none
\end{lstlisting}
\end{leanbox}
\color{black}

\begin{exbox}
\begin{enumerate}
  \item Trace through \li{integrate "x" (Mul (Var "x") (Exp (Var "x")))} by hand,
        showing which strategy fires and the resulting \li{IntStep} tree.
  \item Implement \li{partialFracStrategy} for the case of two distinct linear
        factors: $\int \frac{1}{(x-a)(x-b)}\,dx$ for $a\neq b$.
  \item Implement \li{trigReduceStrategy} for $\int \sin^2(x)\,dx$ using the
        half-angle formula.
  \item Extend \li{directLookup} with antiderivatives for $1/x$, $\tan x$,
        $\arctan x$, and $1/\sqrt{1-x^2}$.
\end{enumerate}
\end{exbox}
\color{black}

%%-----------------------------------------------------------------------------
\chapter{Step-by-Step Trace Rendering}
%%-----------------------------------------------------------------------------
\color{black}

The most important user-facing feature is the step-by-step trace.  This is exactly
what Symbolab renders: each computation step labelled with the rule used.

\section{The \li{SolveTrace} Rose Tree}

A \li{SolveTrace} is a rose tree: each node carries the name of the rule applied, the resulting expression, and a list of child traces for the sub-problems that were solved recursively.  Leaf nodes represent atomic steps (a direct lookup, a base-case derivative); internal nodes represent compound rules (product rule, integration by parts) whose proofs depend on sub-derivations.  The tree structure mirrors the proof tree of the derivation --- rendering it is literally rendering a proof.

\begin{leanbox}
\begin{lstlisting}
-- A rose tree of computation steps
inductive SolveTrace where
  | Leaf   : String -> Expr -> SolveTrace          -- rule name, result
  | Node   : String -> Expr -> List SolveTrace -> SolveTrace

-- Convert DiffStep to SolveTrace
def DiffStep.toTrace (e result : Expr) : DiffStep -> SolveTrace
  | .ConstRule    =>
      .Leaf "Constant Rule" (Const 0)
  | .VarRule      =>
      .Leaf "Variable Rule" (Const 1)
  | .VarOther     =>
      .Leaf "Independent Variable" (Const 0)
  | .SumRule s t  =>
      .Node "Sum Rule" result [s.toTrace e result, t.toTrace e result]
  | .ProductRule s t =>
      .Node "Product Rule" result [s.toTrace e result, t.toTrace e result]
  | .QuotRule s t =>
      .Node "Quotient Rule" result [s.toTrace e result, t.toTrace e result]
  | .ChainRule fn s =>
      .Node s!"Chain Rule ({fn})" result [s.toTrace e result]
  | .PowerRule    =>
      .Leaf "Power Rule" result
  | .NegRule s    =>
      .Node "Negation Rule" result [s.toTrace e result]
\end{lstlisting}
\end{leanbox}
\color{black}

\section{Rendering to Text}

The text renderer performs an in-order traversal of the \li{SolveTrace} tree, printing each node with indentation proportional to its depth.  Each level of indentation corresponds to one level of sub-problem decomposition.  The result is a hierarchical, numbered derivation --- the same format Symbolab uses in its step-by-step display, where main steps appear at the top level and substeps are indented beneath them.

\begin{leanbox}
\begin{lstlisting}
-- Render a SolveTrace as an indented numbered list
def SolveTrace.renderLines (depth : Nat := 0) : SolveTrace -> List String
  | .Leaf rule result =>
      [String.replicate (depth * 2) ' ' ++ s!"• {rule}: {result.pp}"]
  | .Node rule result children =>
      let header := String.replicate (depth * 2) ' '
                      ++ s!"▶ {rule} → {result.pp}"
      let childLines := children.bind (·.renderLines (depth + 1))
      header :: childLines

def SolveTrace.render (t : SolveTrace) : String :=
  "\n".intercalate t.renderLines

-- Full solve-and-render for differentiation
def solveDerivative (x : String) (e : Expr) : String :=
  let (result, step) := diff x e
  let trace := step.toTrace e result
  s!"d/dx({e.pp}) = {result.pp}\n\nSteps:\n{trace.render}"

-- Full solve-and-render for integration
def solveIntegral (x : String) (e : Expr) : String :=
  match integrate x e with
  | some (result, step) =>
      s!"∫({e.pp}) dx = {result.pp} + C\n\n(step trace available)"
  | none =>
      s!"∫({e.pp}) dx — no elementary antiderivative found"
\end{lstlisting}
\end{leanbox}
\color{black}

\begin{example}[Sample Output]
Running \li{solveDerivative "x" (Mul (Var "x") (Sin (Var "x")))} produces:

\begin{verbatim}
d/dx(x * sin(x)) = sin(x) + x * cos(x)

Steps:
▶ Product Rule -> sin(x) + x * cos(x)
  • Variable Rule: 1
  ▶ Chain Rule (sin) -> cos(x)
    • Variable Rule: 1
\end{verbatim}
\end{example}

\begin{exbox}
\begin{enumerate}
  \item Implement \li{IntStep.toTrace} following the same pattern as
        \li{DiffStep.toTrace}.
  \item Add an HTML renderer that wraps each step in \li{<div class="step">}
        tags for browser display.
  \item Extend the trace to record the original sub-expression at each node,
        not just the rule and result.
  \item Run \li{solveIntegral "x" (Mul (Var "x") (Exp (Var "x")))} and verify
        the output matches $(x-1)e^x + C$.
\end{enumerate}
\end{exbox}
\color{black}

%%-----------------------------------------------------------------------------
\chapter{The CLI}
%%-----------------------------------------------------------------------------
\color{black}

\section{Lake Project Structure}

Lake is Lean~4's package manager and build system.  A Lake project specifies its modules (libraries) and executables in a \li{lakefile.lean} at the project root.  Our project has one library (\li{Symboliculus}, containing the expression, differentiation, integration, and trace modules) and one executable (\li{symboliculus}, the CLI entry point).  Running \li{lake build} compiles everything and places the binary at \li{./build/bin/symboliculus}.

\begin{leanbox}
\begin{lstlisting}
-- lakefile.lean
import Lake
open Lake DSL

package symboliculus where
  name    := "symboliculus"
  version := "0.1.0"

lean_lib Symboliculus where
  srcDir := "Symboliculus"

lean_exe symboliculus where
  root := `Main
\end{lstlisting}
\end{leanbox}
\color{black}

\begin{leanbox}
\begin{lstlisting}
-- Symboliculus/
--   Expr.lean       -- Expr type, eval, pp, simplify, parser
--   Diff.lean       -- diff, DiffStep, trace rendering
--   Integrate.lean  -- integrate, IntStep, strategies
--   Trace.lean      -- SolveTrace, render
-- Main.lean
\end{lstlisting}
\end{leanbox}
\color{black}

\section{\li{Main.lean} --- Argument Parsing and Entry Point}

The entry point parses command-line arguments, delegates to the appropriate engine (differentiation or integration), and prints the result.  The \li{--steps} flag controls whether the full trace tree is rendered or just the final answer.  Error handling covers both parse failures (the input string is not a valid expression) and unknown arguments.  The design keeps \li{Main.lean} thin: all the mathematical logic lives in the library modules, and \li{main} is purely plumbing.

\begin{leanbox}
\begin{lstlisting}
-- Main.lean
import Symboliculus.Expr
import Symboliculus.Diff
import Symboliculus.Integrate
import Symboliculus.Trace

open Symboliculus

def usage : String :=
  "symboliculus [--diff | --integral] --expr <expr> --var <var> [--steps]"

structure Args where
  mode  : Mode := .Diff
  expr  : String := ""
  var   : String := "x"
  steps : Bool   := false

inductive Mode where | Diff | Integral

def parseArgs : List String -> Except String Args
  | []         => .ok {}
  | "--diff"   :: rest => do let a <- parseArgs rest; .ok { a with mode := .Diff }
  | "--integral"::rest => do let a <- parseArgs rest; .ok { a with mode := .Integral }
  | "--expr" :: e :: rest =>
      do let a <- parseArgs rest; .ok { a with expr := e }
  | "--var"  :: v :: rest =>
      do let a <- parseArgs rest; .ok { a with var := v }
  | "--steps" :: rest =>
      do let a <- parseArgs rest; .ok { a with steps := true }
  | s :: _ => .error s!"Unknown argument: {s}\n{usage}"

def main (args : List String) : IO Unit := do
  match parseArgs args with
  | .error msg => IO.eprintln msg
  | .ok a =>
    if a.expr == "" then
      IO.eprintln s!"No expression given.\n{usage}"
    else
      match parseStr a.expr with
      | .error msg => IO.eprintln s!"Parse error: {msg}"
      | .ok e =>
        match a.mode with
        | .Diff =>
            let output := if a.steps
                          then solveDerivative a.var e
                          else
                            let (result, _) := diff a.var e
                            s!"d/d{a.var}({e.pp}) = {result.pp}"
            IO.println output
        | .Integral =>
            let output := if a.steps
                          then solveIntegral a.var e
                          else
                            match integrate a.var e with
                            | some (result, _) =>
                                s!"∫({e.pp}) d{a.var} = {result.pp} + C"
                            | none =>
                                s!"∫({e.pp}) d{a.var} = (no elementary form found)"
            IO.println output
\end{lstlisting}
\end{leanbox}
\color{black}

\section{Building and Running}

Below are the canonical invocations for the two main use cases: differentiation with full step trace, and integration with full step trace.  The examples match what you computed by hand in the exercises of Chapters~32 and~33 --- you can verify that the output matches your handwritten derivations.  Try feeding in more complex expressions; the engine should handle any combination of the functions and rules implemented in the library.

\begin{leanbox}
\begin{lstlisting}
-- Build the project
$ lake build

-- Differentiate x^2 * sin(x) with steps shown
$ ./build/bin/symboliculus \
    --diff \
    --expr "x^2 * sin(x)" \
    --var x \
    --steps

-- Output:
-- d/dx(x^2 * sin(x)) = 2*x*sin(x) + x^2*cos(x)
--
-- Steps:
-- ▶ Product Rule -> 2*x*sin(x) + x^2*cos(x)
--   ▶ Power Rule -> 2*x
--     • Variable Rule: 1
--   ▶ Chain Rule (sin) -> cos(x)
--     • Variable Rule: 1

-- Integrate x * exp(x) with steps
$ ./build/bin/symboliculus \
    --integral \
    --expr "x * exp(x)" \
    --var x \
    --steps

-- Output:
-- ∫(x * exp(x)) dx = (x - 1)*exp(x) + C
--
-- Steps:
-- ▶ Integration by Parts (u=x, dv=exp(x)dx)
--   • d/dx(x) = 1
--   • ∫exp(x) dx = exp(x)  [Direct Lookup]
\end{lstlisting}
\end{leanbox}
\color{black}

\begin{exbox}
\begin{enumerate}
  \item Build the project and verify the two examples above run correctly.
  \item Add a \li{--simplify} mode that just simplifies an expression without
        differentiating or integrating.
  \item Add \li{--definite a b} flags so the CLI evaluates
        $\int_a^b f(x)\,dx$ using FTC Part~2 (find antiderivative,
        evaluate at endpoints, subtract).
  \item Add a REPL mode (\li{--repl}) that reads expressions from stdin in a loop,
        displaying results interactively.
\end{enumerate}
\end{exbox}
\color{black}

%%-----------------------------------------------------------------------------
\chapter{Correctness, Incompleteness, and What Comes Next}
%%-----------------------------------------------------------------------------
\color{black}

\section{What We Can Prove: Correctness of Differentiation}

The differentiation engine is \emph{correct}: its output is the mathematical
derivative of its input.  We can state and prove this formally.

\begin{thmbox}{Correctness of \li{diff}}
For any expression \li{e} and variable \li{x}, if \li{f := fun t => e.eval (env[x := t])}
is the function represented by \li{e}, then:
\[\mathtt{(diff\ x\ e).1.eval\ env} = f'(\mathtt{env}\ x).\]
\end{thmbox}
\color{black}

\begin{proof}[Proof by structural induction on \li{e}]
The base cases (\li{Const}, \li{Var}) follow directly.  The inductive cases
correspond exactly to the mathematical differentiation rules proved in Part~III:
sum rule (Ch.\ 11), product rule (Ch.\ 11), chain rule (Ch.\ 11), power rule (Ch.\ 10).
Each case of \li{diff} was written to match its corresponding theorem.
\end{proof}

\begin{leanbox}
\begin{lstlisting}
-- The formal statement in Lean
theorem diff_correct (e : Expr) (x : String)
    (env : String -> Float) :
    let f := fun t => e.eval (fun y => if y == x then t else env y)
    (diff x e).1.eval env = deriv f (env x) := by
  induction e with
  | Const c  => simp [diff, Expr.eval, deriv_const]
  | Var y    =>
      by_cases hxy : y == x
      · simp [diff, hxy, Expr.eval, deriv_id']
      · simp [diff, hxy, Expr.eval, deriv_const]
  | Add u v ih_u ih_v =>
      simp [diff, Expr.eval, deriv_add, ih_u, ih_v]
  | Mul u v ih_u ih_v =>
      simp [diff, Expr.eval, deriv_mul, ih_u, ih_v]
  | Sin u ih_u =>
      simp [diff, Expr.eval, deriv_sin, ih_u,
            Function.comp, Real.deriv_sin]
  -- ... (remaining cases follow the same pattern)
  | _ => sorry
\end{lstlisting}
\end{leanbox}
\color{black}

\section{What We Cannot Decide: The Risch Algorithm}

\begin{remarkbox}
\textbf{Historical Note: The Undecidability of Integration}

The question of whether a given elementary function has an elementary antiderivative is one of the deepest algorithmic questions in mathematics.

\textbf{Liouville} proved in 1835 that $e^{-x^2}$, $\sin(x)/x$, and $\sqrt{1-k^2\sin^2 x}$ have no elementary antiderivatives --- no antiderivative expressible in terms of polynomials, exponentials, logarithms, and trigonometric functions via composition and arithmetic. His proof used an algebraic structure on \emph{differential fields} that constrains the form any antiderivative must take, then derives contradictions when no such form exists. This was among the first examples of a natural mathematical problem shown to be ``impossible'' within a given class of functions.

\textbf{Robert Risch} solved the general case in 1969--1970 with a complete algorithm: given any elementary function, the Risch algorithm either finds an elementary antiderivative or proves none exists. It is correct and always terminates, making symbolic integration of elementary functions \emph{decidable} in principle. In practice, a full implementation requires handling dozens of cases across algebraic, logarithmic, and exponential extensions. Complete implementations exist in Mathematica and Maple; they run to tens of thousands of lines.

\textbf{Symbolab and the pedagogical approach.} Symbolab --- which inspired this capstone --- does not implement Risch. Instead it pattern-matches against the techniques taught in calculus courses: substitution, integration by parts, partial fractions, trigonometric substitution. It displays the matching technique as a ``step'' in human-readable form. This is less complete than Risch but far more pedagogically useful: the steps shown are the ones a student would write. Symboliculus follows the same philosophy. Our correctness theorem guarantees that everything we \emph{do} return is correct; it says nothing about what we fail to find.
\end{remarkbox}
\color{black}

Unlike differentiation, integration is fundamentally harder.

\begin{thmbox}{Risch's Theorem (1969)}
There exists a decision procedure (the \textbf{Risch algorithm}) that, given
an elementary function $f$, either:
\begin{enumerate}[noitemsep]
  \item Computes an elementary antiderivative $F$ with $F'=f$, or
  \item Proves that no elementary antiderivative exists.
\end{enumerate}
\end{thmbox}
\color{black}

\begin{remarkbox}
The Risch algorithm is correct and decidable, but it is \emph{extremely} complex
to implement fully --- a complete implementation (as in Mathematica or Maple) runs
to thousands of lines of code and handles dozens of special cases.  Our engine
covers the cases taught in a calculus course (polynomial, trig, exponential,
integration by parts, $u$-sub, partial fractions) but not the general algorithm.
Examples of elementary functions with no elementary antiderivative:
$e^{-x^2}$, $\sin(x)/x$, $\sqrt{1-k^2\sin^2 x}$ (elliptic integral).
\end{remarkbox}
\color{black}

\section{Connections to the Rest of the Book Series}

Symboliculus is not an isolated project --- every major component connects back
to other books in this series:

\begin{itemize}
  \item \textbf{\li{Expr} and the lambda calculus.}  The \li{Expr} inductive type
        is a term algebra in exactly the sense of the lambda calculus book:
        constructors are introduction rules, pattern matching is elimination,
        and \li{eval} is the denotational semantics.  Adding variable capture and
        substitution makes it a full term rewriting system.

  \item \textbf{Strategy dispatch and propagator networks.}  The integration engine's
        strategy chain --- try \li{directLookup}, fall back to \li{tryUSub}, then
        \li{tryIBP}, etc.\ --- is a propagator network over the lattice
        \li{Option (Expr × IntStep)} ordered by $\bot < \text{Some}$.
        Each strategy is a monotone propagator; the dispatcher finds the least
        fixed point.

  \item \textbf{Correctness proofs and Galois theory.}  The field of rational
        functions over which the Risch algorithm operates is a differential field.
        The question of whether an elementary antiderivative exists is answered by
        Galois theory applied to differential equations --- the same Galois
        theory developed in the Galois theory book.

  \item \textbf{Numerical integration and ML.}  When the symbolic engine returns
        \li{none}, we could fall back to numerical integration (Gaussian quadrature,
        Simpson's rule) --- the numerical methods developed in the ML/AI book's
        optimisation chapter.

  \item \textbf{Optimisation and the Hessian.}  The multivariable optimisation
        chapters (Chs.\ 26, 27) connect directly to the optimisation book:
        gradient descent is the continuous analogue of the discrete propagator
        update, and the Hessian controls convergence speed.
\end{itemize}

\section{Open Problems and Extensions}

\begin{enumerate}
  \item \textbf{Implement the full Risch algorithm.}  Start with the
        ``transcendental case'' (integrands involving $e^x$ and $\ln x$ only);
        the algebraic case (involving $\sqrt{\cdot}$) is substantially harder.

  \item \textbf{Prove correctness of \li{integrate}.}  The correctness theorem
        would state: if \li{integrate x e = some (F, \_)}, then
        \li{diff x F = e} (up to simplification).  This requires a more careful
        treatment of \li{fullSimplify} preserving semantic equality.

  \item \textbf{Add multivariate support.}  Extend \li{diff} to handle
        partial derivatives ($\partial f/\partial x_i$) and then
        \li{integrate} to handle iterated integrals.

  \item \textbf{Connect to Lean's \li{Mathlib}.}  The definitions here are
        self-contained, but every theorem has a Mathlib counterpart.
        A natural project is to build a verified bridge: show that
        \li{MyReal} is isomorphic to Mathlib's \li{Real}, and that
        \li{diff x e} computes Mathlib's \li{HasDerivAt}.

  \item \textbf{Step-counting complexity.}  How many integration strategies
        does the engine try before giving up?  Analyse the worst-case and
        average-case complexity for polynomial, rational, and trigonometric
        families.
\end{enumerate}

\begin{exbox}
\begin{enumerate}
  \item Complete the proof of \li{diff\_correct} in Lean~4, filling in all
        remaining cases (\li{Mul}, \li{Div}, \li{Pow}, \li{Sin}, \li{Cos},
        \li{Exp}, \li{Ln}).
  \item Implement the simplest non-trivial case of the Risch algorithm: the
        ``rational function'' case, which decides when $\int p(x)/q(x)\,dx$ is
        elementary (always) and computes it via partial fractions.
  \item Add a ``verify'' mode to the CLI: given an alleged antiderivative $F$
        and integrand $f$, differentiate $F$ and check symbolically whether
        $F'=f$.
  \item Write a test suite (\li{Tests.lean}) with 20 derivative and 20 integral
        examples, running \li{diff} and \li{integrate} and checking the results
        match known values using \li{Expr.eval} at sample points.
\end{enumerate}
\end{exbox}
\color{black}

\backmatter

\end{document}
